% ============================================================================
% Fichier     : HE2_main.tex
% Titre       : Hacking Éthique II
% Auteur      : MINKA MI NGUIDJOÏ Thierry Emmanuel
% Institution : ENSP Yaoundé · Université de Yaoundé I · LIMSI
% Version     : 1.0 -- novembre 2026
% ============================================================================

\documentclass[12pt, twoside, openright, a4paper]{book}

% --- Feuille de style unifiée ---
\usepackage{Style_HE2}
\raggedbottom  % évite les blancs étirés entre un tableau et sa légende

% --- Bibliographie ---
\addbibresource{bibliographie/HE2_references.bib}

% ============================================================================
% MÉTADONNÉES PDF
% ============================================================================
\hypersetup{
    pdftitle={Hacking Éthique II},
    pdfauthor={MINKA MI NGUIDJOÏ Thierry Emmanuel},
    pdfsubject={SEC4062 · Génie Informatique, niveau 4 · ENSP Yaoundé I},
    pdfkeywords={hacking éthique, test d'intrusion, PTES, OWASP, MITRE ATT&CK, CVSS, équipe rouge}
}

\begin{document}

% --------------------------------------------------------------------------
% ÉLÉMENTS PRÉLIMINAIRES (pages romaines)
% --------------------------------------------------------------------------
\frontmatter

    % ==============================================================================
% Fichier : front/00_page_titre.tex
% Titre   : Page de titre - Hacking Éthique II
% Auteur  : MINKA MI NGUIDJOÏ Thierry Emmanuel
% Version : 1.0 -- novembre 2026 (charte des cours de l'ENSPY)
% ==============================================================================

\begin{titlepage}
\thispagestyle{empty}

% --- Fond décoratif : bandeaux bleu marine et filets or ---
\begin{tikzpicture}[remember picture, overlay]
    \fill[CouleurPrimaire]
        (current page.north west) rectangle
        ([yshift=-3.6cm] current page.north east);
    \fill[CouleurAccent]
        ([yshift=-3.6cm] current page.north west) rectangle
        ([yshift=-4.0cm] current page.north east);
    \fill[CouleurPrimaire]
        (current page.south west) rectangle
        ([yshift=2.5cm] current page.south east);
    \fill[CouleurAccent]
        ([yshift=2.5cm] current page.south west) rectangle
        ([yshift=2.9cm] current page.south east);
    % En-tête institutionnel, posé dans le bandeau supérieur
    \node[align=center, text=white] at ([yshift=-1.8cm] current page.north) {%
        {\fontsize{11}{13}\selectfont
        \textsc{République du Cameroun}}\\[0.3ex]
        {\color{CouleurAccent!60!white}\rule{0.4\paperwidth}{0.4pt}}\\[0.6ex]
        {\fontsize{10}{12}\selectfont\textbf{\tutelle}}\\[0.2ex]
        {\fontsize{9}{11}\selectfont \institutionprimaire}\\[0.2ex]
        {\fontsize{9}{11}\selectfont \unitepedago}};
    % Logos : Université de Yaoundé I à gauche, ENSPY à droite (cours de l'ENSPY)
    \node[fill=white, rounded corners=4pt, inner sep=2pt]
        at ([xshift=1.75cm, yshift=-1.8cm] current page.north west)
        {\includegraphics[height=2.6cm]{figures/logos/uy1_logo.png}};
    \node[fill=white, rounded corners=4pt, inner sep=2pt]
        at ([xshift=-1.75cm, yshift=-1.8cm] current page.north east)
        {\includegraphics[height=2.6cm]{figures/logos/enspy_logo.jpg}};
    % Année et version, posées dans le bandeau inférieur
    \node[align=center] at ([yshift=1.25cm] current page.south) {%
        {\color{white}\bfseries\large Année académique \anneescolaire}\\[0.3ex]
        {\color{CouleurAccent!70!white}\small
        Version 1.0 | Document pédagogique de l'UE \codeUE}};
\end{tikzpicture}

\vspace*{0.6cm}

% --- Titre principal (sur deux lignes) ---
\begin{center}
    {\color{CouleurPrimaire}\rule{0.85\textwidth}{1.2pt}}\\[1.2ex]
    {\fontsize{26}{32}\selectfont\bfseries\color{CouleurPrimaire}
        Hacking\\[0.4ex]
        Éthique II}\\[0.8ex]
    {\color{CouleurAccent}\rule{0.85\textwidth}{2pt}}\\[0.5ex]
    {\color{CouleurPrimaire}\rule{0.85\textwidth}{0.5pt}}
\end{center}

\vspace{0.8cm}

\begin{center}
    {\fontsize{13}{16}\selectfont\itshape\color{CouleurSecondaire}
        Autoriser, reconnaître, tester, restituer~:\\
        le test d'intrusion au service de la défense}
\end{center}

\vspace{0.5cm}

% --- Cartouche sigle et code UE ---
\begin{center}
    \begin{tikzpicture}
        \node[draw=CouleurPrimaire, fill=CouleurPrimaire, rounded corners=6pt,
              inner xsep=14pt, inner ysep=6pt]{%
            \color{white}\bfseries\large \sigleCours{} - UE~\codeUE
        };
    \end{tikzpicture}
\end{center}

\vspace{1.0cm}

% --- Schéma de la discipline : le cycle d'un test d'intrusion ---
\begin{center}
\begin{tikzpicture}[
    f/.style = {draw=CouleurBordDef, fill=CouleurDef, rounded corners=3pt,
        minimum width=2.2cm, minimum height=0.8cm, align=center,
        font=\small\bfseries\color{CouleurPrimaire}},
    arrow/.style = {-{Stealth[length=4pt]}, thick, color=CouleurSecondaire}]
\foreach \t [count=\i] in {Autoriser, Reconnaître, Exploiter, Restituer, Corriger}
    \node[f] (n\i) at (2.8*\i - 2.8, 0) {\t};
\foreach \i [evaluate=\i as \j using int(\i+1)] in {1,2,3,4}
    \draw[arrow] (n\i) -- (n\j);
\draw[arrow] (n5.south) to[bend left=18] node[below, font=\footnotesize\itshape,
    text=CouleurSecondaire] {contre-test} (n1.south);
\end{tikzpicture}
\end{center}

\vspace{0.6cm}

\begin{center}
    {\large\bfseries\color{CouleurPrimaire}
        \auteurcours}\\[0.5ex]
    {\normalsize\color{CouleurSecondaire}
       Enseignant Chercheur | \laboratoire{} | ENSPY}\\[0.2ex]
    {\small\color{CouleurBordInfo}
        \institutionprimaire{} · \tutelle}\\[0.2ex]
    {\small\color{CouleurBordInfo}
        \href{mailto:maletyon@proton.me}{maletyon@proton.me}}
\end{center}

\end{titlepage}

    \cleardoublepage

    % ============================================================================
% Fichier  : front/01_avant_propos.tex
% Titre    : Avant-propos
% Auteur   : MINKA MI NGUIDJOÏ Thierry Emmanuel
% Version  : 1.0 -- novembre 2026
% ============================================================================

\chapter*{Avant-propos}
\addcontentsline{toc}{chapter}{Avant-propos}
\markboth{Avant-propos}{Avant-propos}

\epigraph{« Ce qui distingue le testeur de l'attaquant,\\
ce n'est pas l'outil~: c'est l'autorisation. »}{Principe du cours}

Ce cours s'adresse aux élèves ingénieurs de quatrième année du Département
de Génie Informatique de l'École Nationale Supérieure Polytechnique de
Yaoundé. Il suppose acquis le cours de sécurité informatique (SEC3266), en
particulier les réseaux, les identités, les systèmes et les applications, et
s'articule avec l'audit des systèmes d'information (INF3226) et l'analyse
des risques (GES4062). Il est la seconde partie de l'enseignement du hacking
éthique~: il approfondit la démarche du test d'intrusion, des règles
d'engagement au rapport.

Il repose sur une conviction~: on ne défend bien que ce que l'on sait
attaquer. Mais il en tire une exigence~: chaque technique est enseignée avec
son cadre légal, sa détection et sa parade, et uniquement pour des systèmes
que l'on est autorisé à tester. Le cours ne fournit pas de code d'attaque
prêt à l'emploi~; il apprend à raisonner comme un attaquant pour conseiller
un défenseur.

\bigskip

\textbf{Ce que ce cours n'est pas.} Il n'est pas une préparation à une
certification, ni un catalogue d'outils, ni une invitation à tester sans
mandat. Un test non autorisé est une infraction au Cameroun, même animé par
la curiosité~: la loi vise expressément l'accès obtenu « par défi
intellectuel ».

\bigskip

\textbf{Ce que ce cours est.} Un traité de vingt-six chapitres, en six
parties précédées d'un prologue. La Partie~I pose le cadre~: définitions et
métiers, droit camerounais, règles d'engagement, méthodologies. La
Partie~II apprend à reconnaître~: renseignement en sources ouvertes,
balayage, analyse des vulnérabilités. La Partie~III traite des systèmes et
des réseaux, dont l'annuaire Active Directory. La Partie~IV porte sur les
applications web, les API, le mobile et le nuage. La Partie~V va au-delà des
systèmes~: sans fil, ingénierie sociale, équipe rouge, détection. La
Partie~VI apprend à restituer~: rapport, contre-test, divulgation,
laboratoire, mesure.

\bigskip

\textbf{Un fil rouge.} Après la crise de mars (Opération GAROUA, cours
SEC3266 et GES4062), LiZbiZ commande un test d'intrusion de son système
(Opération LOGONE). La BANQUE SANAGA conduit un exercice d'équipe rouge
(Opération MBAM), et l'ANPP fait tester son portail et son API (Opération
MEFOU). Les noms, les dates, les adresses et les montants sont fictifs~; les
mécanismes sont exacts, et les textes de loi sont cités d'après leur version
officielle.

\bigskip

\textbf{Comment travailler.} Lisez chaque chapitre avec sa fiche
d'auto-évaluation~; faites les exercices avant de lire leurs corrigés~;
pratiquez dans les labs de NDAP SUKULU, qui travaillent sur des pièces et
jamais sur des systèmes réels. Pour aller plus loin, construisez votre
propre laboratoire isolé (Chapitre~\ref{chap:labo}).

\bigskip

\begin{flushright}
    \IfFileExists{../../Logos/signature_avant_propos.png}{\includegraphics[height=1.6cm]{../../Logos/signature_avant_propos.png}}{\textit{MINKA MI NGUIDJOÏ Thierry Emmanuel}}\\
    \textit{Yaoundé, novembre 2026}
\end{flushright}

    \cleardoublepage

    % ============================================================================
% Fichier  : front/02_mode_emploi.tex
% Titre    : Mode d'emploi du cours
% Auteur   : MINKA MI NGUIDJOÏ Thierry Emmanuel
% Version  : 1.0 -- novembre 2026
% ============================================================================

\chapter*{Mode d'emploi du cours}
\addcontentsline{toc}{chapter}{Mode d'emploi du cours}
\markboth{Mode d'emploi}{Mode d'emploi}

\section*{Structure générale}

Le cours compte un prologue et six parties, soit 26~chapitres, pour
60~heures en présentiel réparties en trente séances de deux heures. Les
fiches d'auto-évaluation et les dix laboratoires de NDAP SUKULU sont des
travaux à faire hors de ces 60~heures.

\begin{tableaucours}{L{1.6cm}L{4.4cm}L{4.4cm}L{1.8cm}}{%
    \tableauHeader{Partie} &
    \tableauHeader{Thème} &
    \tableauHeader{Compétences visées} &
    \tableauHeader{Niveau}
}
Prologue & La lettre d'autorisation & Lire un mandat, séparer constatation
et opinion & \niveaubase \\
I & Le cadre du test & Métiers, droit, règles d'engagement, méthodes &
\niveaubase \\
II & Reconnaître & Sources ouvertes, balayage, vulnérabilités &
\niveaumoyen \\
III & Systèmes et réseaux & Exploitation, réseau, Active Directory,
privilèges, mots de passe & \niveaumoyen{} à \niveauexpert \\
IV & Applications et nuage & Web, injections, API, mobile, nuage &
\niveaumoyen{} à \niveauexpert \\
V & Au-delà des systèmes & Sans fil, ingénierie sociale, équipe rouge,
détection & \niveauexpert \\
VI & Restituer et progresser & Rapport, contre-test, divulgation,
laboratoire, mesure & \niveaumoyen \\
\end{tableaucours}

\section*{Les trente séances}

\begin{tableaucours}{L{1.4cm}L{8.4cm}L{2.4cm}}{%
    \tableauHeader{Séance} &
    \tableauHeader{Contenu} &
    \tableauHeader{Partie}
}
1 & Prologue~: la lettre d'autorisation & Prologue \\
2 & Le hacking éthique (Ch.~\ref{chap:ethique}) & I \\
3 & Le cadre juridique (Ch.~\ref{chap:droit}) & I \\
4 & Règles d'engagement (Ch.~\ref{chap:engagement}) & I \\
5 & Méthodologies et référentiels (Ch.~\ref{chap:methodes}) & I \\
6 & Reconnaissance passive (Ch.~\ref{chap:osint}) & II \\
7 & Balayage et énumération (Ch.~\ref{chap:balayage}) & II \\
8 & Analyser et prioriser (Ch.~\ref{chap:vulnerabilites}) & II \\
9 & L'exploitation (Ch.~\ref{chap:exploitation}) & III \\
10 & Attaques réseau (Ch.~\ref{chap:reseau}) & III \\
11--12 & Active Directory (Ch.~\ref{chap:ad}) & III \\
13 & Élévation de privilèges (Ch.~\ref{chap:privileges}) & III \\
14 & Post-exploitation (Ch.~\ref{chap:post}) & III \\
15 & Mots de passe (Ch.~\ref{chap:mdp}) & III \\
16--17 & Tester une application web (Ch.~\ref{chap:web}) & IV \\
18 & Injections et contrôle d'accès (Ch.~\ref{chap:injections}) & IV \\
19 & API et mobile (Ch.~\ref{chap:api}) & IV \\
20 & Nuage et conteneurs (Ch.~\ref{chap:nuage}) & IV \\
21 & Sans fil et objets connectés (Ch.~\ref{chap:sansfil}) & V \\
22 & Ingénierie sociale (Ch.~\ref{chap:social}) & V \\
23--24 & Équipe rouge et équipe violette (Ch.~\ref{chap:rouge}) & V \\
25 & Détection et contournement (Ch.~\ref{chap:detection}) & V \\
26 & Le rapport de test (Ch.~\ref{chap:rapport}) & VI \\
27 & Contre-test (Ch.~\ref{chap:retest}) & VI \\
28 & Divulgation coordonnée (Ch.~\ref{chap:divulgation}) & VI \\
29 & Le laboratoire (Ch.~\ref{chap:labo}) & VI \\
30 & Mesurer~; synthèse (Ch.~\ref{chap:mesurer}) & VI \\
\end{tableaucours}

\section*{Les encadrés}

\begin{boxdef}
\textbf{Définition}. Introduit un terme avec sa définition~; les termes
définis sont indexés et repris dans le glossaire (Annexe~\ref{ann:glossaire}).
\end{boxdef}

\begin{boxscenario}{Titre du scénario}
\textbf{Scénario opérationnel}. Un épisode d'une des opérations du
cours~: LOGONE (test d'intrusion de LiZbiZ), MBAM (équipe rouge de la BANQUE
SANAGA), MEFOU (portail et API de l'ANPP).
\end{boxscenario}

\begin{boxobjet}{T-xx · Intitulé}
\textbf{Fiche de constat}. Présente un constat de test sous une forme
constante~: description, preuve, gravité, impact, recommandation,
détection.
\end{boxobjet}

\begin{boxprocedure}
\textbf{Procédure de test}. Dit, pas à pas, ce qu'il faut faire, dans le
cadre d'une autorisation écrite.
\end{boxprocedure}

\begin{boxmetier}{Pentesteur}
\textbf{Application métier}. Traduit le chapitre en gestes concrets pour
l'un des métiers visés~: pentesteur, auditeur de sécurité, analyste SOC et
équipe rouge, RSSI.
\end{boxmetier}

\begin{boxjuris}
Cite un texte de loi d'après sa version officielle, puis en tire la
portée pour le testeur.
\end{boxjuris}

\begin{boxalert}
\textbf{Violation typique} et \textbf{règle}~: une erreur fréquente, et ce
qu'il faut faire à la place.
\end{boxalert}

\begin{boxexo}
\textbf{Exercice}. Dix par chapitre, gradués en trois niveaux~:
\niveaubase\ fondamental, \niveaumoyen\ intermédiaire, \niveauexpert\
expert. Les corrigés figurent dans le livret de correction.
\end{boxexo}

\section*{Conventions}

\begin{itemize}
    \item Toute technique est présentée pour des systèmes que l'on est
          autorisé à tester, avec sa détection et sa parade.
    \item Chaque analyse indique un niveau de confiance~: faible, moyenne ou
          élevée, et reste séparée de la constatation.
    \item Les adresses IP et les noms de domaine des exemples sont fictifs ou
          réservés à la documentation.
    \item Les \termecle{termes clés} apparaissent en gras~; les normes
          s'écrivent \norme{ISO/IEC~29147}, les lois camerounaises
          \loicam{n°~2010/012}.
\end{itemize}

\section*{Matériel d'accompagnement}

\begin{itemize}
    \item \textbf{Annexes}~: glossaire, outils, modèles et procédures,
          textes juridiques et normatifs, bibliographie commentée, contacts
          (Annexes~\ref{ann:glossaire} à~\ref{ann:contacts}).
    \item \textbf{Livret de correction}~: corrigés des exercices et des
          fiches, réservé au formateur.
    \item \textbf{NDAP SUKULU}~: fiches d'auto-évaluation et laboratoires,
          présentés dans la section qui suit.
\end{itemize}

    \cleardoublepage

    % ==============================================================================
% Fichier : front/03_ndap_sukulu.tex
% Titre   : NDAP SUKULU, l'environnement compagnon du cours
% Auteur  : MINKA MI NGUIDJOÏ Thierry Emmanuel
% Version : 1.0 -- novembre 2026
% ==============================================================================
% Texte commun aux cours de l'auteur ; seul le tableau des labs est propre au
% cours (fichier front/03b_labs.tex).
% ==============================================================================

\chapter*{NDAP SUKULU, la salle de classe compagnon}
\addcontentsline{toc}{chapter}{NDAP SUKULU, la salle de classe compagnon}
\markboth{NDAP SUKULU}{NDAP SUKULU}

\begin{center}
    \includegraphics[width=0.62\textwidth]{figures/logos/logo_ndap_sukulu.png}
\end{center}

\section*{Un nom venu du pays bassa}

\begin{quote}\itshape
Le soir tombe sur Éséka. Un enfant revient des champs avec son grand-père et
lui demande pourquoi il doit marcher si loin, chaque matin, pour aller à
l'école. Le vieil homme s'arrête au bord du chemin et répond, en bassa~:
« \textup{Ndap sukulu}, la salle de classe, ce ne sont pas les murs. Ce sont
les murs qui viennent après. La salle de classe, c'est l'endroit où quelqu'un
t'attend pour t'apprendre ce qu'il sait. » L'enfant insiste~: « Et qui
m'attend~? » Le grand-père sourit~: « \textup{Malet yon}. Ton enseignant. »
\end{quote}

En bassa, la langue paternelle de l'auteur, \textbf{\textit{ndap sukulu}}
désigne la salle de classe et \textbf{\textit{malet yon}} signifie « ton
enseignant ». L'environnement compagnon de ce cours porte ce nom parce qu'il
prolonge la salle de classe au-delà des heures de cours~: il attend
l'étudiant le soir, le week-end, à distance, avec des questions, des pièces à
examiner et une correction immédiate. Il est édité par \textbf{MaletYon},
« ton enseignant », et mis à la disposition des étudiants et des praticiens
qui souhaitent se perfectionner (\href{mailto:maletyon@proton.me}{maletyon@proton.me}).

\section*{Ce que l'on y trouve}

\begin{itemize}
    \item \textbf{Une fiche d'auto-évaluation par chapitre.} Chaque fiche
          puise vingt questions dans une banque de quarante (QCM, vrai ou
          faux, textes à trous). Elle dure trente minutes~; la copie validée
          est verrouillée et notée aussitôt. Le seuil de réussite est de
          60~\%.
    \item \textbf{Des laboratoires pratiques.} Chaque lab place l'étudiant
          devant les pièces d'une affaire du fil rouge (tableurs, journaux,
          contrats, procès-verbaux) dans un bureau d'analyste virtuel,
          ouvert dans le navigateur. Il dure soixante minutes, se termine
          par des points de contrôle notés automatiquement et renvoie aux
          sections du livre qu'il mobilise.
    \item \textbf{Un suivi pour l'enseignant}~: rapport de classe, notes par
          fiche et par lab, journal d'activité.
\end{itemize}

Les fiches et les labs sont des travaux à faire hors des heures de cours.
Ils ne remplacent ni le livre ni les exercices de fin de chapitre~: ils en
vérifient la maîtrise.

\section*{Comment y travailler}

\begin{enumerate}
    \item Se connecter au portail NDAP SUKULU avec ses identifiants et le
          code de cours \textbf{\codeUE}.
    \item Pour un lab, commencer par vérifier l'empreinte SHA-256 des pièces
          (\texttt{sha256sum -c MANIFESTE\_SHA256.txt}), puis travailler sur
          la copie du dossier \texttt{travail}~: les pièces originales sont
          en lecture seule et l'activité est journalisée, comme dans une
          vraie mission.
    \item Lire le rappel de cours affiché par le lab et relire les sections
          citées avant de répondre.
\end{enumerate}

\begin{boxinfo}
\textbf{Si le portail n'est pas encore ouvert.} En attendant le déploiement
de NDAP SUKULU, les fiches de questions sont fournies au format texte dans le
dossier \texttt{fiches} du cours, prêtes à être distribuées. Les corrections
des fiches figurent dans le livret de correction, à la suite des corrigés des
exercices de chaque chapitre.
\end{boxinfo}

\section*{Les laboratoires du cours}

% Tableau des labs du cours HE2 (SEC4062) dans NDAP SUKULU (propre au cours).
Le cours compte dix laboratoires, disponibles sous le code \codeUE. Ils
travaillent tous sur des pièces (relevés, journaux, exports, résultats de
balayage) et jamais sur un système réel~: ils entraînent le raisonnement du
testeur et du défenseur, pas le geste d'attaque.

\begin{tableaucours}{L{2.4cm}L{6.4cm}L{1.6cm}L{1.8cm}}{%
    \tableauHeader{Lab} &
    \tableauHeader{Sujet} &
    \tableauHeader{Chapitre} &
    \tableauHeader{Opération}
}
HE2-LAB-01 & Périmètre~: trier les cibles autorisées & \ref{chap:engagement} &
LOGONE \\
HE2-LAB-02 & OSINT~: sous-domaines et exposition & \ref{chap:osint} & LOGONE \\
HE2-LAB-03 & Lire un résultat de balayage & \ref{chap:balayage} & LOGONE \\
HE2-LAB-04 & Prioriser des vulnérabilités & \ref{chap:vulnerabilites} &
LOGONE \\
HE2-LAB-05 & Active Directory~: chemins vers les administrateurs &
\ref{chap:ad} & LOGONE \\
HE2-LAB-06 & Élévation de privilèges Linux~: relevé & \ref{chap:privileges}
& LOGONE \\
HE2-LAB-07 & Audit de mots de passe & \ref{chap:mdp} & LOGONE \\
HE2-LAB-08 & Historique HTTP~: contrôle d'accès et injection &
\ref{chap:injections} & MEFOU \\
HE2-LAB-09 & Équipe violette~: couverture de détection & \ref{chap:rouge} &
MBAM \\
HE2-LAB-10 & Rapport~: calculer les scores CVSS & \ref{chap:rapport} & LOGONE \\
\end{tableaucours}
\vspace{-0.8em}
\begin{center}{\small\itshape Les laboratoires du cours dans NDAP SUKULU}\end{center}


    \cleardoublepage

    \tableofcontents
    \cleardoublepage

    \listoffigures
    \listoftables
    \cleardoublepage

% CORPS (pages arabes)
% Numérotation des chapitres CONTINUE sur tout le cours : le prologue est le
% chapitre 0, puis C01 à C21, un numéro par fichier.
% --------------------------------------------------------------------------
\mainmatter

    \part*{Prologue Opérationnel}
    \addcontentsline{toc}{part}{Prologue Opérationnel}
        \setcounter{chapter}{-1}
        \IfFileExists{partie0/P00_la_lettre.tex}{% ============================================================================
% Fichier  : partie0/P00_la_lettre.tex
% Titre    : Prologue : la Lettre d'Autorisation
% Auteur   : MINKA MI NGUIDJOÏ Thierry Emmanuel
% Version  : 1.0 -- novembre 2026
% ============================================================================
% Opération LOGONE : test d'intrusion de LiZbiZ par NYONG SÉCURITÉ, du 19 au
% 30 octobre 2026. Cas fictif, cohérent avec SEC3266 et GES4062. Adresses de
% documentation (RFC 5737) ou fictives.
% ============================================================================

\chapter{Prologue : la Lettre d'Autorisation}
\label{chap:prologue}
\chaptermark{Prologue}

\epigraph{%
    « Sans lettre d'autorisation, le meilleur test d'intrusion\\
    est une infraction. »%
}{Principe du cours}

\niveauChapitre{Fondamental}{Cours SEC3266 (Sécurité informatique).}

\begin{boxobjectifs}
\begin{itemize}
    \item Lire un mandat de test d'intrusion et en repérer les éléments
          essentiels.
    \item Situer le test dans la vie d'une organisation~: après une crise,
          après un audit, avant une décision.
    \item Distinguer ce que le test prouve de ce qu'il ne prouve pas.
    \item Séparer constatations, suppositions et opinions dans un résultat
          de test.
    \item Formuler les questions auxquelles le cours répond.
\end{itemize}
\end{boxobjectifs}

\bigskip

% ============================================================================
\section*{Question directrice : que faut-il signer avant de toucher un seul système~?}
% ============================================================================

Le 12~octobre 2026, M.~ETEME signe la lettre qui autorise le cabinet NYONG
SÉCURITÉ à attaquer le système d'information de LiZbiZ. Sept mois plus tôt,
un attaquant avait fait la même chose sans demander la permission
(Opération GAROUA). Cette fois, tout est écrit~: qui teste, quoi, quand,
comment, et qui appeler si quelque chose tourne mal. Ce prologue suit la
mission, appelée Opération LOGONE, de la signature au rapport.

\medskip
\textbf{Plan du chapitre.}
\begin{enumerate}
    \item Pourquoi un test, et pourquoi maintenant (section~\ref{sec:p00-pourquoi}).
    \item La lettre d'autorisation (section~\ref{sec:p00-lettre}).
    \item Douze jours de test (section~\ref{sec:p00-test}).
    \item Ce que le test a montré (section~\ref{sec:p00-resultats}).
    \item Les questions du cours (section~\ref{sec:p00-questions}).
\end{enumerate}

\separateur

% ============================================================================
\section{Pourquoi un Test, et Pourquoi Maintenant}
\label{sec:p00-pourquoi}
% ============================================================================

LiZbiZ\footnote{LiZbiZ, NYONG SÉCURITÉ, la BANQUE SANAGA, l'ANPP et leurs
personnels sont fictifs, comme les dates, les adresses et les montants du
cours~; les techniques, les normes et les textes de loi cités sont réels.}
a beaucoup changé depuis mars~: second facteur sur la passerelle, bastion,
sauvegardes immuables, registre des risques, audit d'octobre. Le comité des
risques veut savoir si ces mesures résistent à un attaquant réel, et non
seulement à une liste de vérification. Il commande un \termecle{test
d'intrusion}~: une attaque autorisée, bornée et restituée.

\begin{tableaucours}{L{3.4cm}L{8.8cm}}{%
    \tableauHeader{Acteur} &
    \tableauHeader{Rôle dans l'Opération LOGONE}
}
M.~ETEME & Directeur général de LiZbiZ, signataire de la lettre
d'autorisation \\
M.~KOUNOU & Responsable informatique, contact technique d'urgence \\
Mme~MOUAFO & Gestionnaire des risques, commanditaire opérationnelle \\
M.~BIKORO & Chef de mission, NYONG SÉCURITÉ \\
Mme~NGO~MBOCK & Pentesteuse, tests internes et Active Directory \\
M.~TAMO & Analyste, tests web et ingénierie sociale \\
\end{tableaucours}
\vspace{-0.8em}
\begin{center}{\small\itshape Tableau~0.1. Les acteurs de l'Opération LOGONE}\end{center}

% ============================================================================
\section{La Lettre d'Autorisation}
\label{sec:p00-lettre}
% ============================================================================

\begin{tableaucours}{L{3.4cm}L{8.8cm}}{%
    \tableauHeader{Rubrique} &
    \tableauHeader{Contenu de la lettre du 12 octobre}
}
Parties & LiZbiZ SA (M.~ETEME) et NYONG SÉCURITÉ (M.~BIKORO) \\
Objectif & Vérifier qu'un attaquant externe ou interne ne peut pas obtenir
les droits d'administrateur du domaine ni lire la base des clients \\
Périmètre externe & Plage 198.51.100.0/27 et domaine lizbiz.cm \\
Périmètre interne & Réseaux de Douala (10.10.0.0/16) et de Garoua
(10.20.0.0/16), sauf 10.10.50.0/24 (paie, hébergée pour un tiers) \\
Exclusions & Déni de service~; système de paie~; intrusion physique \\
Ingénierie sociale & Courriels d'hameçonnage simulés à 60~adresses validées
par la direction des ressources humaines \\
Fenêtre & Du 19 au 30~octobre~; tests internes de 8~h à 18~h, sauf accord
écrit \\
Contacts d'urgence & M.~KOUNOU et M.~BIKORO, joignables jour et nuit~; mot
d'arrêt convenu \\
Données & Aucune copie de données personnelles~; preuves masquées~;
destruction à la remise du rapport \\
\end{tableaucours}
\vspace{-0.8em}
\begin{center}{\small\itshape Tableau~0.2. Extrait de la lettre d'autorisation de l'Opération LOGONE}\end{center}

\begin{boxalert}
\textbf{Violation typique.} Tester une adresse « voisine » du périmètre
parce qu'elle semble appartenir au client.

\smallskip
\textbf{Règle~:} seul ce qui est écrit est autorisé. Le réseau de la paie,
hébergé pour un tiers, n'appartient pas au signataire~: le tester, même par
erreur, serait un accès non autorisé (Chapitre~\ref{chap:droit}).
\end{boxalert}

% ============================================================================
\section{Douze Jours de Test}
\label{sec:p00-test}
% ============================================================================

\begin{tableaucours}{L{2.4cm}L{9.8cm}}{%
    \tableauHeader{Dates} &
    \tableauHeader{Activité}
}
19--20~oct. & Reconnaissance externe~: sources ouvertes, sous-domaines,
services exposés \\
21--22~oct. & Tests externes~: passerelle, portail de suivi, messagerie \\
23~oct. & Campagne d'hameçonnage simulée \\
26--28~oct. & Tests internes depuis un poste fourni~: balayage, Active
Directory, partages \\
29~oct. & Démonstration du chemin vers les administrateurs du domaine~;
arrêt des tests \\
30~oct. & Réunion de restitution à chaud \\
\end{tableaucours}
\vspace{-0.8em}
\begin{center}{\small\itshape Tableau~0.3. Le calendrier de l'Opération LOGONE}\end{center}

% ============================================================================
\section{Ce que le Test a Montré}
\label{sec:p00-resultats}
% ============================================================================

\begin{boxscenario}{Opération LOGONE. La restitution du 30 octobre}
\textbf{Constatations.}
\begin{itemize}
    \item Externe~: la passerelle exige bien un second facteur~; aucun accès
          obtenu depuis Internet.
    \item Hameçonnage~: 14~clics sur 60~destinataires, 5~saisies de mot de
          passe, 9~signalements au support.
    \item Interne~: en 5~heures, depuis un poste ordinaire, Mme~NGO~MBOCK
          obtient les droits d'administrateur du domaine en exploitant le
          mot de passe faible d'un compte de service.
    \item Détection~: le centre de surveillance de nuit signale le balayage
          interne en 40~minutes, mais ne voit pas l'obtention des droits.
    \item Bilan~: 14~constats, dont 2 critiques, 4 élevés, 5 moyens et
          3 faibles.
\end{itemize}

\medskip
\textbf{Analyse} (confiance élevée). La porte d'entrée de mars est fermée~;
la maison, elle, reste ouverte à qui y entre. Le test ne prouve pas
qu'aucune autre voie n'existe~: il prouve que celle-ci existait, à cette
date, dans ce périmètre.
\end{boxscenario}

\begin{boxmetier}{Pentesteur}
\textbf{Arrêter quand l'objectif est démontré.} Mme~NGO~MBOCK s'arrête
dès qu'elle obtient les droits~: elle capture la preuve (une capture
d'écran du groupe, sans données), prévient M.~BIKORO, qui prévient
M.~KOUNOU. Aller plus loin n'apporterait rien et créerait un risque.
\end{boxmetier}

\begin{boxmetier}{RSSI}
\textbf{Commander un test pour décider.} Le RSSI commande un test avec une
question précise (« peut-on devenir administrateur du domaine~? »), et non
pour obtenir une note. Le rapport doit l'aider à décider où investir.
\end{boxmetier}

% ============================================================================
\section{Les Questions du Cours}
\label{sec:p00-questions}
% ============================================================================

\begin{tableaucours}{L{3.4cm}L{6.6cm}L{2.2cm}}{%
    \tableauHeader{Partie} &
    \tableauHeader{Question posée par l'Opération LOGONE} &
    \tableauHeader{Chapitres}
}
I. Cadre & Qu'est-ce qui rend ce test légal et utile~? & \ref{chap:ethique}
à~\ref{chap:methodes} \\
II. Reconnaître & Qu'a-t-on appris sans toucher aux systèmes~? &
\ref{chap:osint} à~\ref{chap:vulnerabilites} \\
III. Systèmes & Comment un compte de service mène-t-il au domaine~? &
\ref{chap:exploitation} à~\ref{chap:mdp} \\
IV. Applications & Le portail et l'API sont-ils sûrs~? & \ref{chap:web}
à~\ref{chap:nuage} \\
V. Au-delà & Les personnes et le SOC tiennent-ils~? & \ref{chap:sansfil}
à~\ref{chap:detection} \\
VI. Restituer & Comment écrire, corriger, vérifier, mesurer~? &
\ref{chap:rapport} à~\ref{chap:mesurer} \\
\end{tableaucours}
\vspace{-0.8em}
\begin{center}{\small\itshape Tableau~0.4. Les questions du cours}\end{center}

\separateur

\begin{boxresume}
\textbf{Ce qu'il faut retenir de ce chapitre~:}
\begin{itemize}
    \item \textbf{Un test d'intrusion} est une attaque autorisée, bornée et
          restituée.
    \item \textbf{La lettre d'autorisation} fixe parties, objectif,
          périmètre, exclusions, fenêtre, contacts et traitement des données.
    \item \textbf{Seul ce qui est écrit est autorisé}~; un tiers non
          signataire est hors périmètre.
    \item \textbf{On s'arrête} quand l'objectif est démontré.
    \item \textbf{Un test prouve} l'existence d'une voie, pas l'absence
          d'autres voies.
\end{itemize}
\end{boxresume}

\begin{boxinfo}
\textbf{Points de contrôle de l'élève.} Je sais~: énumérer les rubriques
d'une lettre d'autorisation~; repérer une cible hors périmètre~; résumer
les résultats de l'Opération LOGONE~; dire ce qu'un test prouve et ne
prouve pas.
\end{boxinfo}

\bigskip

\begin{boxexo}[title={\ding{45}\quad Exercice~0.1. Le taux \niveaubase}]
Calculez le taux de clics, de saisies et de signalements de la campagne
d'hameçonnage.
\end{boxexo}

\begin{boxexo}[title={\ding{45}\quad Exercice~0.2. Hors périmètre \niveaubase}]
Les adresses 10.10.50.12, 10.20.3.50 et 198.51.100.40 sont-elles dans le
périmètre de la lettre~? Justifiez.
\end{boxexo}

\begin{boxexo}[title={\ding{45}\quad Exercice~0.3. Les registres \niveaubase}]
Classez en constatation, supposition ou opinion~: « le SOC est
incompétent »~; « le balayage a été signalé en 40~minutes »~; « un
attaquant réel aurait été plus discret ».
\end{boxexo}

\begin{boxexo}[title={\ding{45}\quad Exercice~0.4. Prouver \niveaumoyen}]
Pourquoi le test ne prouve-t-il pas que LiZbiZ est sûre depuis Internet~?
Que peut-on en dire honnêtement~?
\end{boxexo}

\begin{boxexo}[title={\ding{45}\quad Exercice~0.5. L'arrêt \niveaumoyen}]
Mme~NGO~MBOCK aurait pu lire la base des clients après avoir obtenu les
droits. Pourquoi s'arrête-t-elle~? Quelle preuve garde-t-elle~?
\end{boxexo}

\begin{boxexo}[title={\ding{45}\quad Exercice~0.6. Le tiers \niveaumoyen}]
Le système de paie est hébergé par LiZbiZ pour un tiers. Que faudrait-il
pour l'inclure dans un prochain test~?
\end{boxexo}

\begin{boxexo}[title={\ding{45}\quad Exercice~0.7. Les données \niveaumoyen}]
Pourquoi la lettre interdit-elle de copier des données personnelles~? Quels
risques courrait le cabinet sans cette clause~?
\end{boxexo}

\begin{boxexo}[title={\ding{45}\quad Exercice~0.8. La détection \niveaumoyen}]
Le SOC a vu le balayage mais pas l'obtention des droits. Qu'en concluez-vous
pour ses règles de détection~?
\end{boxexo}

\begin{boxexo}[title={\ding{45}\quad Exercice~0.9. Une lettre \niveauexpert}]
Rédigez la lettre d'autorisation complète de l'Opération LOGONE, à partir
du Tableau~0.2.
\end{boxexo}

\begin{boxexo}[title={\ding{45}\quad Exercice~0.10. La question \niveauexpert}]
Proposez trois questions précises que la BANQUE SANAGA pourrait poser à un
test d'intrusion, et pour chacune le périmètre nécessaire.
\end{boxexo}

\bigskip

\begin{boxinfo}
\textbf{Pour aller plus loin.}
\begin{itemize}
    \item NIST SP~800-115~\cite{nist800115}.
    \item Loi n°~2010/012, articles~65 et suivants~\cite{loi2010012}.
\end{itemize}
\end{boxinfo}

\begin{boxinfo}
\textbf{Transition.} Avant les techniques, les mots et les métiers. Le
Chapitre~\ref{chap:ethique} définit le hacking éthique et les types de
tests.
\end{boxinfo}
}{}

    \part{Le Cadre du Test}
    \label{part:cadre}

        \IfFileExists{partie1/P01_C01_hacking_ethique.tex}{% ============================================================================
% Fichier  : partie1/P01_C01_hacking_ethique.tex
% Titre    : Le Hacking Éthique : Définitions, Métiers, Types de Tests
% Auteur   : MINKA MI NGUIDJOÏ Thierry Emmanuel
% Version  : 1.0 -- novembre 2026
% ============================================================================
% Sources : NIST SP 800-115 (2008) ; PTES ; CIS Controls v8 (contrôle 18).
% ============================================================================

\chapter{Le Hacking Éthique : Définitions, Métiers, Types de Tests}
\label{chap:ethique}
\chaptermark{Le hacking éthique}

\epigraph{%
    « Le testeur et l'attaquant connaissent les mêmes techniques~;\\
    seul le premier rend compte de ce qu'il a fait. »%
}{Principe du cours}

\niveauChapitre{Fondamental}{Prologue~; cours SEC3266.}

\begin{boxobjectifs}
\begin{itemize}
    \item Définir le hacking éthique et le test d'intrusion.
    \item Distinguer test d'intrusion, analyse de vulnérabilités, audit,
          équipe rouge et programme de primes.
    \item Choisir une approche (boîte noire, grise, blanche) et un point de
          départ (externe, interne, présumé compromis).
    \item Décrire les métiers de la sécurité offensive et leurs
          compétences.
    \item Situer la valeur et les limites d'un test.
\end{itemize}
\end{boxobjectifs}

\bigskip

% ============================================================================
\section*{Question directrice : de quoi LiZbiZ avait-elle besoin~: un balayage, un audit ou un test d'intrusion~?}
% ============================================================================

En septembre, M.~KOUNOU propose d'acheter un scanner de vulnérabilités~;
Mme~NKENG rappelle que l'audit d'octobre est prévu~; Mme~MOUAFO veut savoir
si un attaquant peut encore devenir administrateur du domaine. Trois
besoins, trois réponses différentes. Ce chapitre apprend à choisir.

\medskip
\textbf{Plan du chapitre.}
\begin{enumerate}
    \item Définir (section~\ref{sec:c01-definir}).
    \item Les types d'évaluation offensive (section~\ref{sec:c01-types}).
    \item Approches et points de départ (section~\ref{sec:c01-approches}).
    \item Les métiers (section~\ref{sec:c01-metiers}).
    \item Valeur et limites d'un test (section~\ref{sec:c01-limites}).
\end{enumerate}

\separateur

% ============================================================================
\section{Définir}
\label{sec:c01-definir}
% ============================================================================

\begin{boxdef}
\termecle{Hacking éthique}~: emploi des techniques des attaquants, avec
l'autorisation écrite du propriétaire du système, dans un périmètre et une
durée fixés, pour en révéler les faiblesses et aider à les corriger.
\termecle{Test d'intrusion}~: évaluation dans laquelle le testeur imite un
attaquant réel pour contourner les mesures de sécurité et démontrer
l'impact d'une compromission~\cite{nist800115}.
\end{boxdef}

Trois conditions font le caractère éthique~: l'\emph{autorisation} (écrite,
par qui a le pouvoir de la donner), la \emph{limitation} (périmètre,
méthodes, fenêtre) et la \emph{restitution} (tout ce qui est trouvé est
rapporté au client, et à lui seul). Sans l'une d'elles, le même geste devient
une infraction ou une faute professionnelle.

% ============================================================================
\section{Les Types d'Évaluation Offensive}
\label{sec:c01-types}
% ============================================================================

\begin{tableaucours}{L{3.2cm}L{4.6cm}L{4.4cm}}{%
    \tableauHeader{Type} &
    \tableauHeader{Question} &
    \tableauHeader{Livrable}
}
Analyse de vulnérabilités & Quelles failles connues sont présentes~? &
Liste priorisée, souvent automatique \\
Test d'intrusion & Jusqu'où un attaquant peut-il aller~? & Chemins
d'attaque démontrés, constats, recommandations \\
Audit de sécurité & Les mesures sont-elles conformes à un référentiel~? &
Constats d'écarts (cours INF3226) \\
Exercice d'équipe rouge & Détecte-t-on et répond-on à un adversaire
réaliste~? & Scénario, détections, délais \\
Programme de primes & Des chercheurs externes trouvent-ils des failles~? &
Rapports de chercheurs, rémunérés selon la gravité \\
\end{tableaucours}
\vspace{-0.8em}
\begin{center}{\small\itshape Tableau~1.1. Les évaluations offensives}\end{center}

Les CIS Controls consacrent leur contrôle~18 aux tests d'intrusion~: ils
s'adressent aux organisations qui ont déjà mis en place les mesures de
base~; tester avant d'avoir inventorié et corrigé n'apprend rien que le
balayage n'aurait dit~\cite{cis2021controls}.

% ============================================================================
\section{Approches et Points de Départ}
\label{sec:c01-approches}
% ============================================================================

\begin{tableaucours}{L{3.0cm}L{4.8cm}L{4.4cm}}{%
    \tableauHeader{Approche} &
    \tableauHeader{Ce que sait le testeur} &
    \tableauHeader{Intérêt}
}
Boîte noire & Rien, sauf le périmètre & Imite l'attaquant externe~;
coûteuse en temps \\
Boîte grise & Un compte, des plans partiels & Bon compromis, le plus
fréquent \\
Boîte blanche & Code, configurations, schémas & Couverture maximale \\
\end{tableaucours}
\vspace{-0.8em}
\begin{center}{\small\itshape Tableau~1.2. Les approches}\end{center}

Le \termecle{point de départ} compte autant~: \emph{externe} (depuis
Internet), \emph{interne} (depuis un poste du réseau) ou \emph{présumé
compromis} (on suppose qu'un poste est déjà aux mains de l'attaquant, comme
après l'hameçonnage de février). Ce dernier scénario, dit d'\emph{intrusion
présumée}, mesure la défense en profondeur.

\begin{boxscenario}{Opération LOGONE. Le choix de l'approche}
\textbf{Décision} (septembre). Boîte grise~: un compte d'utilisateur
ordinaire et un poste fourni pour l'interne~; les schémas réseau sont
communiqués. Deux points de départ~: externe, et intrusion présumée.

\medskip
\textbf{Analyse} (confiance élevée). La boîte noire aurait consacré des jours
à retrouver ce que LiZbiZ sait déjà~; l'intrusion présumée répond à la vraie
question~: que se passe-t-il quand le prochain hameçonnage réussit~?
\end{boxscenario}

% ============================================================================
\section{Les Métiers}
\label{sec:c01-metiers}
% ============================================================================

\begin{tableaucours}{L{3.4cm}L{8.8cm}}{%
    \tableauHeader{Métier} &
    \tableauHeader{Activité}
}
Pentesteur & Conduit les tests, démontre les chemins d'attaque, rédige les
constats \\
Auditeur de sécurité & Vérifie la conformité aux référentiels, peut
intégrer des tests techniques \\
Analyste SOC et équipe rouge & L'un détecte et répond, l'autre émule
l'adversaire~; ensemble, ils forment l'équipe violette \\
RSSI & Commande, cadre, arbitre les corrections, rend compte à la direction \\
\end{tableaucours}
\vspace{-0.8em}
\begin{center}{\small\itshape Tableau~1.3. Les métiers visés par le cours}\end{center}

\begin{boxmetier}{Pentesteur}
\textbf{Trois compétences au même niveau.} La technique (systèmes, réseaux,
applications), la méthode (périmètre, traçabilité, arrêt) et l'écrit (un
constat qu'un directeur comprend et qu'un administrateur sait corriger).
\end{boxmetier}

% ============================================================================
\section{Valeur et Limites d'un Test}
\label{sec:c01-limites}
% ============================================================================

\begin{boxalert}
\textbf{Violation typique.} Présenter un test sans constat critique comme
une preuve de sécurité.

\smallskip
\textbf{Règle~:} un test est borné dans le temps, le périmètre et les
méthodes~; il prouve l'existence des failles trouvées, pas l'absence des
autres. Le rapport dit ce qui n'a pas été testé.
\end{boxalert}

\begin{boxobjet}{T-00 · Modèle de fiche de constat}
\textbf{Description~:} ce qui a été trouvé, en une phrase. \textbf{Preuve~:}
la trace (capture, requête, horodatage), sans données sensibles.
\textbf{Gravité~:} score CVSS et contexte. \textbf{Impact~:} ce qu'un
attaquant pourrait faire. \textbf{Recommandation~:} la correction et sa
vérification. \textbf{Détection~:} ce qui aurait dû alerter.
\end{boxobjet}

\separateur

\begin{boxresume}
\textbf{Ce qu'il faut retenir de ce chapitre~:}
\begin{itemize}
    \item \textbf{Hacking éthique}~: autorisation, limitation, restitution.
    \item \textbf{Cinq évaluations offensives} répondent à des questions
          différentes.
    \item \textbf{Boîte noire, grise, blanche}~; départ externe, interne ou
          présumé compromis.
    \item \textbf{Quatre métiers}~; l'équipe violette réunit attaque et
          défense.
    \item \textbf{Un test} prouve l'existence d'une faille, pas l'absence des
          autres.
\end{itemize}
\end{boxresume}

\begin{boxinfo}
\textbf{Points de contrôle de l'élève.} Je sais~: définir le hacking éthique
et le test d'intrusion~; choisir un type d'évaluation pour une question
donnée~; choisir une approche et un point de départ~; décrire les métiers~;
énoncer les limites d'un test.
\end{boxinfo}

\bigskip

\begin{boxexo}[title={\ding{45}\quad Exercice~1.1. Quel type~? \niveaubase}]
Quel type d'évaluation répond à~: « nos serveurs ont-ils des correctifs en
retard~? »~; « notre SOC verrait-il un attaquant~? »~; « sommes-nous
conformes à l'ISO/IEC~27001~? »
\end{boxexo}

\begin{boxexo}[title={\ding{45}\quad Exercice~1.2. Les trois conditions \niveaubase}]
Un étudiant teste le site de son école « pour l'aider ». Quelles conditions
du hacking éthique manquent~?
\end{boxexo}

\begin{boxexo}[title={\ding{45}\quad Exercice~1.3. Approche \niveaubase}]
Quelle approche choisir pour l'ANPP, qui fournit le code de son portail~?
\end{boxexo}

\begin{boxexo}[title={\ding{45}\quad Exercice~1.4. Intrusion présumée \niveaumoyen}]
Expliquez en quoi le point de départ « présumé compromis » répond mieux à la
question de LiZbiZ qu'un test externe seul.
\end{boxexo}

\begin{boxexo}[title={\ding{45}\quad Exercice~1.5. Contrôle 18 \niveaumoyen}]
Pourquoi les CIS Controls placent-ils les tests d'intrusion en dernier~?
Que doit faire une organisation avant~?
\end{boxexo}

\begin{boxexo}[title={\ding{45}\quad Exercice~1.6. Violette \niveaumoyen}]
Qu'apporte une équipe violette par rapport à un test d'intrusion classique~?
\end{boxexo}

\begin{boxexo}[title={\ding{45}\quad Exercice~1.7. Primes \niveaumoyen}]
Quels prérequis une organisation doit-elle réunir avant d'ouvrir un
programme de primes aux failles~?
\end{boxexo}

\begin{boxexo}[title={\ding{45}\quad Exercice~1.8. Limites \niveaumoyen}]
Rédigez le paragraphe « limites » du rapport de l'Opération LOGONE.
\end{boxexo}

\begin{boxexo}[title={\ding{45}\quad Exercice~1.9. Un programme \niveauexpert}]
Proposez à la BANQUE SANAGA un programme d'évaluations offensives sur deux
ans, en combinant les types du Tableau~1.1.
\end{boxexo}

\begin{boxexo}[title={\ding{45}\quad Exercice~1.10. Le profil \niveauexpert}]
Rédigez la fiche de poste d'un pentesteur junior pour NYONG SÉCURITÉ~:
missions, compétences, éthique, progression.
\end{boxexo}

\bigskip

\begin{boxinfo}
\textbf{Pour aller plus loin.}
\begin{itemize}
    \item NIST SP~800-115~\cite{nist800115}~; CIS Controls~\cite{cis2021controls}.
\end{itemize}
\end{boxinfo}

\begin{boxinfo}
\textbf{Transition.} L'autorisation n'est pas une formalité~: c'est elle
que la loi regarde. Le Chapitre~\ref{chap:droit} présente le droit
camerounais.
\end{boxinfo}
}{}
        \IfFileExists{partie1/P01_C02_droit.tex}{% ============================================================================
% Fichier  : partie1/P01_C02_droit.tex
% Titre    : Le Cadre Juridique Camerounais
% Auteur   : MINKA MI NGUIDJOÏ Thierry Emmanuel
% Version  : 1.0 -- novembre 2026
% ============================================================================
% Sources : loi n° 2010/012 du 21 décembre 2010 (texte officiel : art. 65,
% 66, 68, 69, 70, 71) ; loi n° 2024/017 (principes seulement).
% ============================================================================

\chapter{Le Cadre Juridique Camerounais}
\label{chap:droit}
\chaptermark{Le cadre juridique}

\epigraph{%
    « La curiosité n'est pas une autorisation. »%
}{Principe du cours}

\niveauChapitre{Fondamental}{Chapitre~\ref{chap:ethique}~; cours SEC3266,
chapitre sur le droit.}

\begin{boxobjectifs}
\begin{itemize}
    \item Citer les infractions de la loi n°~2010/012 qui concernent le
          testeur.
    \item Expliquer pourquoi l'autorisation écrite fait la différence entre
          un test et une infraction.
    \item Identifier qui peut autoriser un test, et qui ne le peut pas.
    \item Repérer les risques juridiques propres aux données personnelles,
          aux tiers et aux outils.
    \item Adopter les réflexes qui protègent le testeur et son client.
\end{itemize}
\end{boxobjectifs}

\bigskip

% ============================================================================
\section*{Question directrice : qu'est-ce qui sépare Mme NGO MBOCK de l'attaquant de mars~?}
% ============================================================================

En mars, un attaquant est entré dans le système de LiZbiZ, s'y est
maintenu et a obtenu les droits d'administrateur. En octobre, Mme~NGO~MBOCK
a fait presque la même chose. Les gestes se ressemblent~; les qualifications
juridiques sont opposées. La différence tient en une feuille signée. Ce
chapitre en explique la portée.

\medskip
\textbf{Plan du chapitre.}
\begin{enumerate}
    \item Les infractions qui concernent le testeur (section~\ref{sec:c02-infractions}).
    \item L'autorisation (section~\ref{sec:c02-autorisation}).
    \item Les données, les tiers, les outils (section~\ref{sec:c02-risques}).
    \item Les réflexes du testeur (section~\ref{sec:c02-reflexes}).
\end{enumerate}

\separateur

% ============================================================================
\section{Les Infractions qui Concernent le Testeur}
\label{sec:c02-infractions}
% ============================================================================

\begin{boxjuris}
\textbf{\loicam{n°~2010/012} du 21~décembre 2010 relative à la
cybersécurité et la cybercriminalité au Cameroun} (texte officiel)

\medskip
\textbf{Article~65, alinéa~2}~: \textit{« Est puni des peines prévues à
l'alinéa 1 ci-dessus, tout accès non autorisé, à l'ensemble ou à une partie
d'un réseau de communications électroniques ou d'un système d'information
ou d'un équipement terminal. »}

\medskip
\textbf{Article~65, alinéa~4}~: \textit{« Est puni des mêmes peines prévues
à l'alinéa 1 ci-dessus, celui qui, sans droit, permet l'accès dans un réseau
de communications électroniques ou dans un système d'information par défi
intellectuel. »}

\medskip
\textbf{Portée.} Les peines de l'alinéa~1 sont de cinq à dix ans
d'emprisonnement et de 5 à 10~millions de FCFA d'amende, ou l'une de ces
deux peines. L'alinéa~4 vise expressément le « défi intellectuel »~: la
bonne intention ou la curiosité ne justifient pas un accès sans droit.
\end{boxjuris}

\begin{tableaucours}{L{2.0cm}L{6.0cm}L{4.2cm}}{%
    \tableauHeader{Article} &
    \tableauHeader{Comportement visé} &
    \tableauHeader{Geste de test concerné}
}
65 (1) & Interception de données sans droit ni autorisation & Écoute du
réseau, interception d'échanges \\
65 (2) et (4) & Accès non autorisé~; accès permis par défi intellectuel &
Toute connexion à un système hors mandat \\
66 & Perturbation ou interruption du fonctionnement par l'altération des
données & Test qui fait tomber un service \\
68 & Accès ou maintien frauduleux avec perturbation grave & Exploitation
qui endommage \\
69 & Accès sans droit, en violation des mesures de sécurité, pour obtenir
des données & Lecture de données au-delà du nécessaire \\
70 & Attaque par saturation pour faire effondrer une ressource & Test de
charge ou de déni de service \\
71 & Introduction de données pour supprimer ou modifier les données &
Modification de données de production \\
\end{tableaucours}
\vspace{-0.8em}
\begin{center}{\small\itshape Tableau~2.1. Les articles de la loi n°~2010/012 qui concernent le testeur}\end{center}

% ============================================================================
\section{L'Autorisation}
\label{sec:c02-autorisation}
% ============================================================================

L'autorisation écrite rend l'accès « autorisé » au sens de l'article~65.
Encore faut-il qu'elle émane de la personne qui a le pouvoir de la
donner~: le propriétaire ou l'exploitant du système, par son représentant
habilité.

\begin{tableaucours}{L{4.4cm}L{7.8cm}}{%
    \tableauHeader{Qui peut autoriser~?} &
    \tableauHeader{Commentaire}
}
Le dirigeant de l'organisation & Oui, ou son délégataire désigné par écrit \\
Le responsable informatique & Seulement s'il a reçu délégation \\
Un salarié qui « veut savoir » & Non \\
L'hébergeur ou le fournisseur de nuage & Pour son infrastructure~; le
client peut avoir besoin de son accord selon le contrat \\
Un tiers dont les systèmes sont reliés & Non pour le système du client~;
il autorise seulement ses propres systèmes \\
\end{tableaucours}
\vspace{-0.8em}
\begin{center}{\small\itshape Tableau~2.2. Le pouvoir d'autoriser}\end{center}

\begin{boxalert}
\textbf{Violation typique.} Se contenter d'un accord oral ou d'un courriel
d'un technicien.

\smallskip
\textbf{Règle~:} l'autorisation est écrite, signée par une personne
habilitée, datée, et précise le périmètre et les méthodes. Le testeur en
garde une copie sur lui pendant les tests physiques et dans son dossier.
\end{boxalert}

% ============================================================================
\section{Les Données, les Tiers, les Outils}
\label{sec:c02-risques}
% ============================================================================

\begin{boxinfo}
\textbf{Les données personnelles.} Un test peut faire voir des données
personnelles (fichier des clients, paie). La loi n°~2024/017 relative à la
protection des données à caractère personnel s'applique à leur traitement~;
ses articles sont à lire dans le texte publié au Journal officiel. La règle
du cours~: ne consulter que ce qui prouve la faille, ne rien copier,
masquer les preuves, détruire ce qui reste à la fin de la mission.
\end{boxinfo}

\begin{tableaucours}{L{3.0cm}L{9.2cm}}{%
    \tableauHeader{Risque} &
    \tableauHeader{Parade}
}
Tiers touché & Vérifier la propriété de chaque adresse avant de tester~;
exclure ce qui appartient à d'autres \\
Interruption de service & Exclure le déni de service~; prévenir avant
toute action risquée~; mot d'arrêt \\
Outils & Utiliser des outils de provenance sûre, sur des postes dédiés~;
ne pas diffuser d'outil offensif hors cadre \\
Preuves & Tracer chaque action (horodatage, adresse source) pour pouvoir
prouver ce qui a été fait, et ce qui ne l'a pas été \\
\end{tableaucours}
\vspace{-0.8em}
\begin{center}{\small\itshape Tableau~2.3. Risques juridiques et parades}\end{center}

\begin{boxscenario}{Opération LOGONE. L'incident évité}
\textbf{Constatation.} Le 26~octobre, M.~TAMO s'apprête à balayer la plage
10.10.0.0/16 entière. Mme~NGO~MBOCK lui rappelle l'exclusion de
10.10.50.0/24 (paie, hébergée pour un tiers). La commande est corrigée avant
d'être lancée~; l'incident est consigné.

\medskip
\textbf{Analyse} (confiance élevée). Le balayage aurait touché un système
dont LiZbiZ n'est pas propriétaire~: un accès non autorisé au sens de
l'article~65, et une faute contractuelle du cabinet.
\end{boxscenario}

% ============================================================================
\section{Les Réflexes du Testeur}
\label{sec:c02-reflexes}
% ============================================================================

\begin{boxprocedure}
\textbf{Avant, pendant, après}
\begin{enumerate}
    \item \textbf{Avant}~: lettre signée par une personne habilitée~;
          périmètre vérifié adresse par adresse~; contrat avec clause de
          responsabilité et assurance du cabinet.
    \item \textbf{Pendant}~: rester dans le périmètre et la fenêtre~;
          journaliser chaque action~; s'arrêter et prévenir au moindre
          doute ou incident.
    \item \textbf{Après}~: remettre le rapport au seul client~; détruire les
          données et les accès créés~; conserver le journal des actions.
\end{enumerate}
\end{boxprocedure}

\begin{boxmetier}{Auditeur de sécurité}
\textbf{Vérifier le mandat avant le premier geste.} L'auditeur qui intègre
des tests techniques vérifie lui-même la lettre, le périmètre et le
pouvoir du signataire~; il ne se fie pas à ce que lui dit l'audité.
\end{boxmetier}

\separateur

\begin{boxresume}
\textbf{Ce qu'il faut retenir de ce chapitre~:}
\begin{itemize}
    \item \textbf{L'article~65} réprime l'accès non autorisé, y compris
          « par défi intellectuel ».
    \item \textbf{Les articles~66 à~71} visent la perturbation, le maintien
          frauduleux, l'obtention de données, la saturation, la modification.
    \item \textbf{L'autorisation} est écrite, signée par une personne
          habilitée, précise.
    \item \textbf{Données, tiers, outils, preuves}~: quatre risques à
          maîtriser.
    \item \textbf{Avant, pendant, après}~: les réflexes qui protègent le
          testeur et le client.
\end{itemize}
\end{boxresume}

\begin{boxinfo}
\textbf{Points de contrôle de l'élève.} Je sais~: citer l'article~65 et ses
alinéas~2 et~4~; relier un geste de test à un article~; dire qui peut
autoriser~; énumérer les réflexes avant, pendant et après.
\end{boxinfo}

\bigskip

\begin{boxexo}[title={\ding{45}\quad Exercice~2.1. Quel article~? \niveaubase}]
Quel article vise~: un test de charge qui fait tomber un site~; la lecture
de données clients au-delà du nécessaire~; un accès « pour voir »~?
\end{boxexo}

\begin{boxexo}[title={\ding{45}\quad Exercice~2.2. Défi intellectuel \niveaubase}]
Que signifie, pour un étudiant passionné, l'alinéa~4 de l'article~65~?
\end{boxexo}

\begin{boxexo}[title={\ding{45}\quad Exercice~2.3. Les peines \niveaubase}]
Quelles sont les peines prévues à l'article~65, alinéa~1~?
\end{boxexo}

\begin{boxexo}[title={\ding{45}\quad Exercice~2.4. Qui autorise~? \niveaumoyen}]
Un administrateur de la BANQUE SANAGA vous envoie un courriel « vous pouvez
tester notre serveur ». Que faites-vous~?
\end{boxexo}

\begin{boxexo}[title={\ding{45}\quad Exercice~2.5. Le nuage \niveaumoyen}]
LiZbiZ veut faire tester sa messagerie en ligne. Quelles autorisations
faut-il vérifier~?
\end{boxexo}

\begin{boxexo}[title={\ding{45}\quad Exercice~2.6. Les données \niveaumoyen}]
Pendant un test, vous accédez au fichier des clients. Que faites-vous,
étape par étape~?
\end{boxexo}

\begin{boxexo}[title={\ding{45}\quad Exercice~2.7. L'incident évité \niveaumoyen}]
Rédigez l'entrée du journal de mission sur l'incident évité du 26~octobre.
\end{boxexo}

\begin{boxexo}[title={\ding{45}\quad Exercice~2.8. L'outil \niveaumoyen}]
Un collègue veut publier sur Internet l'outil qu'il a écrit pendant une
mission. Quels risques juridiques et contractuels voyez-vous~?
\end{boxexo}

\begin{boxexo}[title={\ding{45}\quad Exercice~2.9. Le contrat \niveauexpert}]
Rédigez les clauses juridiques d'un contrat de test d'intrusion~:
autorisation, responsabilité, assurance, données, confidentialité.
\end{boxexo}

\begin{boxexo}[title={\ding{45}\quad Exercice~2.10. Comparer \niveauexpert}]
Comparez la situation de l'attaquant de mars et celle de Mme~NGO~MBOCK
au regard des articles~65, 68 et~69~: éléments matériels, autorisation,
conséquences.
\end{boxexo}

\bigskip

\begin{boxinfo}
\textbf{Pour aller plus loin.}
\begin{itemize}
    \item Loi n°~2010/012~\cite{loi2010012}~; loi n°~2024/017~\cite{loi2024017}.
\end{itemize}
\end{boxinfo}

\begin{boxinfo}
\textbf{Transition.} L'autorisation se traduit en règles concrètes~: les
règles d'engagement, objet du Chapitre~\ref{chap:engagement}.
\end{boxinfo}
}{}
        \IfFileExists{partie1/P01_C03_engagement.tex}{% ============================================================================
% Fichier  : partie1/P01_C03_engagement.tex
% Titre    : Règles d'Engagement et Éthique de la Mission
% Auteur   : MINKA MI NGUIDJOÏ Thierry Emmanuel
% Version  : 1.0 -- novembre 2026
% ============================================================================
% Sources : NIST SP 800-115 (2008), planification ; PTES (pré-engagement) ;
% codes de déontologie (ISC2, chapitre éthique de SEC3266).
% ============================================================================

\chapter{Règles d'Engagement et Éthique de la Mission}
\label{chap:engagement}
\chaptermark{Règles d'engagement}

\epigraph{%
    « Ce qui n'est pas écrit dans les règles d'engagement\\
    est interdit. »%
}{Principe du cours}

\niveauChapitre{Fondamental}{Chapitres~\ref{chap:ethique}
et~\ref{chap:droit}.}

\begin{boxobjectifs}
\begin{itemize}
    \item Rédiger les règles d'engagement d'un test.
    \item Définir un périmètre vérifiable, adresse par adresse.
    \item Organiser la communication et l'arrêt d'urgence.
    \item Traiter les données rencontrées et les découvertes fortuites.
    \item Appliquer les principes déontologiques du testeur.
\end{itemize}
\end{boxobjectifs}

\bigskip

% ============================================================================
\section*{Question directrice : que faire quand on trouve, au milieu d'un test, une trace d'attaque réelle~?}
% ============================================================================

Le 27~octobre, Mme~NGO~MBOCK découvre sur un serveur de Garoua un compte
« svc\_maj2 » créé le 4~mars~: la trace de l'attaquant de l'Opération
GAROUA, oubliée lors du nettoyage. Ce n'est pas l'objet du test. Que
faire~? Les règles d'engagement le disent, si elles ont été bien écrites.

\medskip
\textbf{Plan du chapitre.}
\begin{enumerate}
    \item Les règles d'engagement (section~\ref{sec:c03-regles}).
    \item Le périmètre (section~\ref{sec:c03-perimetre}).
    \item Communiquer et s'arrêter (section~\ref{sec:c03-communiquer}).
    \item Données et découvertes fortuites (section~\ref{sec:c03-donnees}).
    \item La déontologie du testeur (section~\ref{sec:c03-deontologie}).
\end{enumerate}

\separateur

% ============================================================================
\section{Les Règles d'Engagement}
\label{sec:c03-regles}
% ============================================================================

\begin{boxdef}
\termecle{Règles d'engagement}~: document qui précise, pour une mission de
test, ce qui est autorisé, comment, quand, par qui, et ce qu'il faut faire
en cas d'incident. Il complète la lettre d'autorisation et la rend
opérationnelle.
\end{boxdef}

\begin{tableaucours}{L{3.4cm}L{8.8cm}}{%
    \tableauHeader{Rubrique} &
    \tableauHeader{Contenu}
}
Objectifs & Questions auxquelles le test répond, drapeaux à atteindre \\
Périmètre & Adresses, domaines, applications, comptes~; exclusions \\
Méthodes & Autorisées (balayage, exploitation contrôlée, hameçonnage
simulé) et interdites (déni de service, destruction) \\
Fenêtre & Dates, heures, fuseau horaire~; tests de nuit sur accord \\
Équipe & Noms des testeurs, adresses IP sources, postes utilisés \\
Communication & Contacts, canal chiffré, points d'avancement \\
Arrêt & Mot d'arrêt, motifs d'arrêt, reprise \\
Données & Ce qu'on peut consulter, ce qu'on ne copie jamais, destruction \\
Découvertes & Conduite à tenir face à une compromission réelle ou à une
faille critique \\
\end{tableaucours}
\vspace{-0.8em}
\begin{center}{\small\itshape Tableau~3.1. Les rubriques des règles d'engagement}\end{center}

% ============================================================================
\section{Le Périmètre}
\label{sec:c03-perimetre}
% ============================================================================

Un périmètre vérifiable se décrit en adresses et en noms précis, avec leur
propriétaire. Une plage en notation CIDR (10.20.0.0/16) couvre 65\,536
adresses~; une exclusion (10.10.50.0/24) en retire 256. Avant chaque
action, le testeur vérifie que la cible est dedans.

\begin{boxprocedure}
\textbf{Vérifier qu'une cible est dans le périmètre}
\begin{enumerate}
    \item La cible appartient-elle à une plage ou à un domaine listés~?
    \item Appartient-elle à une exclusion~? Si oui, elle est hors périmètre.
    \item Le signataire en est-il propriétaire ou exploitant~? (vérifier les
          registres d'adresses et le contrat d'hébergement)
    \item La méthode envisagée est-elle autorisée pour cette cible~?
    \item L'heure est-elle dans la fenêtre~?
\end{enumerate}
\end{boxprocedure}

\begin{boxalert}
\textbf{Violation typique.} Suivre un lien ou une redirection vers un
domaine non listé (un prestataire de paiement, un fournisseur en nuage).

\smallskip
\textbf{Règle~:} un domaine découvert pendant le test n'entre pas dans le
périmètre de lui-même. On le note et on demande un avenant écrit.
\end{boxalert}

% ============================================================================
\section{Communiquer et S'Arrêter}
\label{sec:c03-communiquer}
% ============================================================================

\begin{tableaucours}{L{4.0cm}L{8.2cm}}{%
    \tableauHeader{Situation} &
    \tableauHeader{Conduite}
}
Service ralenti ou arrêté & Arrêt immédiat, appel du contact d'urgence \\
Faille critique exploitable depuis Internet & Prévenir le client sans
attendre le rapport \\
Trace d'une compromission réelle & Arrêt du test sur la cible, appel du
contact, préservation \\
Doute sur le périmètre & Ne pas agir~; demander par écrit \\
Mot d'arrêt prononcé par le client & Arrêt de toutes les actions \\
\end{tableaucours}
\vspace{-0.8em}
\begin{center}{\small\itshape Tableau~3.2. Quand prévenir, quand s'arrêter}\end{center}

% ============================================================================
\section{Données et Découvertes Fortuites}
\label{sec:c03-donnees}
% ============================================================================

\begin{boxscenario}{Opération LOGONE. Le compte du 4 mars}
\textbf{Constatation.} Le 27~octobre à 10~h~12, Mme~NGO~MBOCK trouve,
membre des administrateurs du domaine, le compte « svc\_maj2 », créé le
4~mars à 2~h~04 (cours SEC3266, lab de journaux).

\medskip
\textbf{Conduite.} Elle s'arrête, ne touche pas au compte, note l'heure et
la preuve, appelle M.~BIKORO, qui appelle M.~KOUNOU à 10~h~25 selon les
règles d'engagement. LiZbiZ désactive le compte, ouvre un incident et fait
vérifier l'absence d'autre persistance~; le test reprend le lendemain.

\medskip
\textbf{Analyse} (confiance élevée). La découverte relève de la réponse aux
incidents, pas du test. Exploiter ce compte aurait mêlé l'activité de
l'attaquant et celle du testeur, et détruit des preuves.
\end{boxscenario}

% ============================================================================
\section{La Déontologie du Testeur}
\label{sec:c03-deontologie}
% ============================================================================

\begin{tableaucours}{L{3.6cm}L{8.6cm}}{%
    \tableauHeader{Principe} &
    \tableauHeader{Application}
}
Légalité & Rien hors mandat~; rien d'interdit par la loi \\
Minimisation & Faire le minimum nécessaire pour prouver~; ne pas aller
au-delà de l'objectif \\
Confidentialité & Résultats au seul client~; aucune publication sans
accord \\
Loyauté & Rapporter tout, y compris ses propres erreurs pendant le test \\
Compétence & Ne pas employer une technique que l'on ne maîtrise pas sur un
système de production \\
\end{tableaucours}
\vspace{-0.8em}
\begin{center}{\small\itshape Tableau~3.3. Les principes déontologiques du testeur}\end{center}

\begin{boxobjet}{T-15 · Découverte fortuite (hors constats du test)}
\textbf{Description~:} compte « svc\_maj2 » membre des administrateurs du
domaine, créé le 4~mars par l'attaquant. \textbf{Preuve~:} liste des
membres du groupe et date de création (capture du 27~octobre, 10~h~12).
\textbf{Gravité~:} critique (accès persistant d'un attaquant).
\textbf{Recommandation~:} désactiver, ouvrir un incident, rechercher toute
autre persistance. \textbf{Détection~:} la revue des comptes à privilèges
aurait dû le voir.
\end{boxobjet}

\begin{boxmetier}{Analyste SOC et équipe rouge}
\textbf{Distinguer le testeur de l'attaquant.} Les règles d'engagement
donnent au SOC les adresses sources et les fenêtres du testeur. Toute
activité hors de ces adresses ou de ces heures se traite comme une attaque
réelle.
\end{boxmetier}

\separateur

\begin{boxresume}
\textbf{Ce qu'il faut retenir de ce chapitre~:}
\begin{itemize}
    \item \textbf{Les règles d'engagement} rendent l'autorisation
          opérationnelle~: objectifs, périmètre, méthodes, fenêtre, équipe,
          communication, arrêt, données, découvertes.
    \item \textbf{Le périmètre} se vérifie adresse par adresse avant chaque
          action.
    \item \textbf{On s'arrête et on prévient} en cas d'incident, de faille
          critique ou de doute.
    \item \textbf{Une compromission réelle} découverte relève de la réponse
          aux incidents.
    \item \textbf{Légalité, minimisation, confidentialité, loyauté,
          compétence.}
\end{itemize}
\end{boxresume}

\begin{boxinfo}
\textbf{Points de contrôle de l'élève.} Je sais~: rédiger des règles
d'engagement~; vérifier qu'une cible est dans le périmètre~; dire quand
s'arrêter et qui prévenir~; traiter une découverte fortuite~; appliquer
les principes déontologiques.
\end{boxinfo}

\bigskip

\begin{boxexo}[title={\ding{45}\quad Exercice~3.1. CIDR \niveaubase}]
Combien d'adresses comptent 198.51.100.0/27 et 10.10.50.0/24~? L'adresse
198.51.100.40 est-elle dans la première plage~?
\end{boxexo}

\begin{boxexo}[title={\ding{45}\quad Exercice~3.2. Les rubriques \niveaubase}]
Quelle rubrique des règles d'engagement répond à~: « puis-je tester à
22~h~? »~; « que faire si le site tombe~? »~; « puis-je copier un fichier
pour preuve~? »
\end{boxexo}

\begin{boxexo}[title={\ding{45}\quad Exercice~3.3. Redirection \niveaubase}]
Le portail de suivi redirige vers « paiement.exemple-prestataire.cm ».
Pouvez-vous le tester~? Que faites-vous~?
\end{boxexo}

\begin{boxexo}[title={\ding{45}\quad Exercice~3.4. Faille critique \niveaumoyen}]
Vous découvrez le premier jour une faille critique exploitable depuis
Internet. Attendez-vous le rapport final~? Rédigez le message au client.
\end{boxexo}

\begin{boxexo}[title={\ding{45}\quad Exercice~3.5. Le SOC \niveaumoyen}]
Faut-il prévenir le SOC d'un test~? Comparez les avantages d'un test annoncé
et d'un test non annoncé.
\end{boxexo}

\begin{boxexo}[title={\ding{45}\quad Exercice~3.6. Minimisation \niveaumoyen}]
Donnez trois exemples de gestes qui dépassent le minimum nécessaire pour
prouver une faille.
\end{boxexo}

\begin{boxexo}[title={\ding{45}\quad Exercice~3.7. Sa propre erreur \niveaumoyen}]
Un testeur a supprimé par erreur un fichier de configuration. Que dit le
principe de loyauté~? Que doit-il faire~?
\end{boxexo}

\begin{boxexo}[title={\ding{45}\quad Exercice~3.8. Mot d'arrêt \niveaumoyen}]
Proposez la procédure d'arrêt et de reprise d'un test, avec ses
responsables.
\end{boxexo}

\begin{boxexo}[title={\ding{45}\quad Exercice~3.9. Des règles complètes \niveauexpert}]
Rédigez les règles d'engagement de l'Opération MEFOU (portail et API de
l'ANPP), selon les rubriques du Tableau~3.1.
\end{boxexo}

\begin{boxexo}[title={\ding{45}\quad Exercice~3.10. Le dilemme \niveauexpert}]
Pendant un test, vous trouvez la preuve qu'un salarié détourne des fonds.
Analysez la situation selon la méthode de résolution d'un dilemme (cours
SEC3266) et dites ce que vous faites.
\end{boxexo}

\bigskip

\begin{boxinfo}
\textbf{Pour aller plus loin.}
\begin{itemize}
    \item NIST SP~800-115, planification~\cite{nist800115}.
    \item PTES, phase de pré-engagement~\cite{ptes}.
\end{itemize}
\end{boxinfo}

\begin{boxinfo}
\textbf{Transition.} Le cadre posé, il faut une méthode pour conduire le
test~: le Chapitre~\ref{chap:methodes} présente les méthodologies et les
référentiels.
\end{boxinfo}
}{}
        \IfFileExists{partie1/P01_C04_methodologies.tex}{% ============================================================================
% Fichier  : partie1/P01_C04_methodologies.tex
% Titre    : Méthodologies et Référentiels
% Auteur   : MINKA MI NGUIDJOÏ Thierry Emmanuel
% Version  : 1.0 -- novembre 2026
% ============================================================================
% Sources : NIST SP 800-115 (2008) ; PTES ; OSSTMM 3 (ISECOM) ; OWASP WSTG ;
% MITRE ATT&CK ; FIRST, CVSS.
% ============================================================================

\chapter{Méthodologies et Référentiels}
\label{chap:methodes}
\chaptermark{Méthodologies et référentiels}

\epigraph{%
    « Une méthode ne remplace pas le talent~;\\
    elle empêche d'oublier ce que le talent n'a pas vu. »%
}{Principe du cours}

\niveauChapitre{Intermédiaire}{Chapitres~\ref{chap:ethique}
à~\ref{chap:engagement}.}

\begin{boxobjectifs}
\begin{itemize}
    \item Décrire les phases d'un test selon le NIST SP~800-115 et le PTES.
    \item Situer l'OSSTMM, l'OWASP WSTG et les autres guides.
    \item Utiliser MITRE ATT\&CK pour décrire un comportement d'attaquant.
    \item Exprimer la gravité d'une faille par le CVSS.
    \item Tenir le journal de mission qui rend un test reproductible.
\end{itemize}
\end{boxobjectifs}

\bigskip

% ============================================================================
\section*{Question directrice : comment être sûr de ne rien avoir oublié~?}
% ============================================================================

Au premier jour, M.~BIKORO remet à son équipe une liste de vérification
fondée sur des référentiels publics, et un journal de mission vierge.
Chaque test y sera noté. Quand LiZbiZ demandera « avez-vous testé la
messagerie~? », la réponse sera écrite, avec l'heure.

\medskip
\textbf{Plan du chapitre.}
\begin{enumerate}
    \item Les phases d'un test (section~\ref{sec:c04-phases}).
    \item Les guides spécialisés (section~\ref{sec:c04-guides}).
    \item Décrire l'attaquant~: ATT\&CK (section~\ref{sec:c04-attack}).
    \item Coter la gravité~: CVSS (section~\ref{sec:c04-cvss}).
    \item Le journal de mission (section~\ref{sec:c04-journal}).
\end{enumerate}

\separateur

% ============================================================================
\section{Les Phases d'un Test}
\label{sec:c04-phases}
% ============================================================================

\begin{tableaucours}{L{3.0cm}L{4.6cm}L{4.6cm}}{%
    \tableauHeader{Phase} &
    \tableauHeader{NIST SP~800-115} &
    \tableauHeader{PTES}
}
Préparer & Planification & Pré-engagement \\
Reconnaître & Découverte & Collecte de renseignements, modélisation des
menaces \\
Analyser & Découverte (vulnérabilités) & Analyse des vulnérabilités \\
Exploiter & Attaque & Exploitation, post-exploitation \\
Restituer & Rapport & Rapport \\
\end{tableaucours}
\vspace{-0.8em}
\begin{center}{\small\itshape Tableau~4.1. Les phases d'un test (d'après le NIST SP~800-115 et le PTES)}\end{center}

Le NIST SP~800-115 décrit les techniques de test et d'évaluation de la
sécurité de l'information~\cite{nist800115}. Le PTES (\emph{Penetration
Testing Execution Standard}) détaille sept phases, du pré-engagement au
rapport~\cite{ptes}. Les phases ne sont pas strictement successives~: une
découverte en exploitation relance la reconnaissance.

% ============================================================================
\section{Les Guides Spécialisés}
\label{sec:c04-guides}
% ============================================================================

\begin{tableaucours}{L{3.4cm}L{8.8cm}}{%
    \tableauHeader{Référentiel} &
    \tableauHeader{Apport}
}
OSSTMM (ISECOM) & Méthode de mesure de la sécurité opérationnelle, avec
des métriques de surface d'attaque~\cite{osstmm3} \\
OWASP WSTG & Guide de test des applications web, scénario par
scénario~\cite{owaspwstg} \\
OWASP MASTG et MASVS & Test et exigences des applications mobiles \\
NIST SP~800-115 & Cadre général, techniques, rapport \\
\end{tableaucours}
\vspace{-0.8em}
\begin{center}{\small\itshape Tableau~4.2. Guides de test}\end{center}

% ============================================================================
\section{Décrire l'Attaquant : ATT\&CK}
\label{sec:c04-attack}
% ============================================================================

\begin{boxdef}
\termecle{MITRE ATT\&CK}~: base de connaissances publique, fondée sur des
observations réelles, qui décrit les tactiques (le but de l'attaquant à une
étape) et les techniques (la manière de l'atteindre) des
adversaires~\cite{mitreattack}.
\end{boxdef}

\begin{tableaucours}{L{3.6cm}L{8.6cm}}{%
    \tableauHeader{Tactique} &
    \tableauHeader{Exemple dans l'Opération GAROUA}
}
Accès initial & Comptes valides volés par hameçonnage \\
Exécution & Lancement d'un outil sur le serveur de fichiers \\
Persistance & Création du compte svc\_maj2 \\
Élévation de privilèges & Exploitation d'un serveur non corrigé \\
Découverte & Balayage du réseau interne \\
Mouvement latéral & Connexions avec le mot de passe local partagé \\
Exfiltration & Copie de 40~Go vers un stockage externe \\
Impact & Chiffrement des serveurs \\
\end{tableaucours}
\vspace{-0.8em}
\begin{center}{\small\itshape Tableau~4.3. Quelques tactiques d'ATT\&CK appliquées à l'attaque de mars}\end{center}

Décrire un test en techniques ATT\&CK permet au SOC de vérifier, technique
par technique, s'il aurait détecté chaque geste (Chapitre~\ref{chap:rouge}).

% ============================================================================
\section{Coter la Gravité : CVSS}
\label{sec:c04-cvss}
% ============================================================================

\begin{boxdef}
\termecle{CVSS} (\emph{Common Vulnerability Scoring System}), publié par le
FIRST~: système qui exprime la gravité technique d'une vulnérabilité par un
score de 0 à 10, calculé à partir de métriques (vecteur d'attaque,
complexité, privilèges requis, interaction de l'utilisateur, impacts sur la
confidentialité, l'intégrité et la disponibilité)~\cite{first2019cvss31}.
\end{boxdef}

\begin{tableaucours}{L{3.0cm}L{2.6cm}}{%
    \tableauHeader{Score} &
    \tableauHeader{Qualification}
}
0,0 & Aucune \\
0,1 à 3,9 & Faible \\
4,0 à 6,9 & Moyenne \\
7,0 à 8,9 & Élevée \\
9,0 à 10,0 & Critique \\
\end{tableaucours}
\vspace{-0.8em}
\begin{center}{\small\itshape Tableau~4.4. Les qualifications CVSS}\end{center}

\begin{boxalert}
\textbf{Violation typique.} Classer les constats par le seul score CVSS.

\smallskip
\textbf{Règle~:} le score dit la gravité technique~; la priorité dépend
aussi de l'exposition, de la valeur de l'actif et de l'exploitation connue
(Chapitre~\ref{chap:vulnerabilites}). Le score de base se calcule, le
contexte se juge.
\end{boxalert}

% ============================================================================
\section{Le Journal de Mission}
\label{sec:c04-journal}
% ============================================================================

\begin{boxprocedure}
\textbf{Tenir le journal de mission}
\begin{enumerate}
    \item Une ligne par action~: heure, testeur, adresse source, cible,
          outil ou commande, résultat.
    \item Les captures et les fichiers de preuve, nommés et horodatés.
    \item Les incidents et les arrêts, avec qui a été prévenu.
    \item Les tests prévus mais non faits, et pourquoi.
    \item Conservation selon le contrat, accès restreint.
\end{enumerate}
\end{boxprocedure}

\begin{boxscenario}{Opération LOGONE. La question de la messagerie}
\textbf{Constatation.} Le 5~novembre, M.~KOUNOU demande si la messagerie en
ligne a été testée. Le journal répond~: test de la configuration
(authentification, règles de transfert) le 22~octobre de 9~h~10 à 11~h~40,
aucune faille~; pas de test de l'interface d'administration, hors périmètre.

\medskip
\textbf{Analyse} (confiance élevée). Sans journal, la réponse aurait été un
souvenir. Avec lui, c'est une constatation vérifiable.
\end{boxscenario}

\begin{boxmetier}{Pentesteur}
\textbf{Écrire au fil de l'eau.} Le journal se tient pendant le test, pas
le soir. Une action non notée est une action que l'on ne pourra ni
expliquer au client, ni défendre en cas d'incident.
\end{boxmetier}

\separateur

\begin{boxresume}
\textbf{Ce qu'il faut retenir de ce chapitre~:}
\begin{itemize}
    \item \textbf{Préparer, reconnaître, analyser, exploiter, restituer}~:
          les phases du NIST SP~800-115 et du PTES.
    \item \textbf{OSSTMM, OWASP WSTG, MASTG} spécialisent la méthode.
    \item \textbf{ATT\&CK} décrit l'attaquant en tactiques et techniques.
    \item \textbf{Le CVSS} cote la gravité technique~; la priorité y ajoute
          le contexte.
    \item \textbf{Le journal de mission} rend le test vérifiable.
\end{itemize}
\end{boxresume}

\begin{boxinfo}
\textbf{Points de contrôle de l'élève.} Je sais~: nommer les phases d'un
test~; choisir un guide spécialisé~; décrire une attaque en tactiques
ATT\&CK~; qualifier un score CVSS~; tenir un journal de mission.
\end{boxinfo}

\bigskip

\begin{boxexo}[title={\ding{45}\quad Exercice~4.1. Les phases \niveaubase}]
Dans quelle phase se range~: la lettre d'autorisation~; le balayage des
ports~; la démonstration d'accès aux administrateurs du domaine~?
\end{boxexo}

\begin{boxexo}[title={\ding{45}\quad Exercice~4.2. Qualifier \niveaubase}]
Qualifiez les scores CVSS 9,8~; 7,5~; 5,3~; 3,1.
\end{boxexo}

\begin{boxexo}[title={\ding{45}\quad Exercice~4.3. Tactiques \niveaubase}]
Rangez dans les tactiques d'ATT\&CK~: la suppression des sauvegardes~; la
connexion avec le mot de passe du technicien~; la copie de 40~Go.
\end{boxexo}

\begin{boxexo}[title={\ding{45}\quad Exercice~4.4. Le guide \niveaumoyen}]
Quel guide utiliser pour tester l'API de l'ANPP~? Et son application mobile~?
\end{boxexo}

\begin{boxexo}[title={\ding{45}\quad Exercice~4.5. Le contexte \niveaumoyen}]
Deux failles ont un score de 8,1~: l'une sur un serveur exposé à Internet,
l'autre sur un poste isolé. Laquelle corriger d'abord~? Pourquoi~?
\end{boxexo}

\begin{boxexo}[title={\ding{45}\quad Exercice~4.6. Le journal \niveaumoyen}]
Rédigez cinq lignes du journal de mission du 26~octobre (tests internes).
\end{boxexo}

\begin{boxexo}[title={\ding{45}\quad Exercice~4.7. Non fait \niveaumoyen}]
Pourquoi noter les tests prévus mais non faits~? Donnez deux raisons
légitimes de ne pas faire un test prévu.
\end{boxexo}

\begin{boxexo}[title={\ding{45}\quad Exercice~4.8. ATT\&CK et SOC \niveaumoyen}]
Comment la description d'un test en techniques ATT\&CK aide-t-elle le SOC
à s'améliorer~?
\end{boxexo}

\begin{boxexo}[title={\ding{45}\quad Exercice~4.9. Liste de vérification \niveauexpert}]
Construisez la liste de vérification des tests internes de l'Opération
LOGONE, phase par phase.
\end{boxexo}

\begin{boxexo}[title={\ding{45}\quad Exercice~4.10. Comparer \niveauexpert}]
Comparez le NIST SP~800-115, le PTES et l'OSSTMM~: objet, niveau de détail,
usage recommandé.
\end{boxexo}

\bigskip

\begin{boxinfo}
\textbf{Pour aller plus loin.}
\begin{itemize}
    \item NIST SP~800-115~\cite{nist800115}~; PTES~\cite{ptes}~; OSSTMM~\cite{osstmm3}.
    \item MITRE ATT\&CK~\cite{mitreattack}~; CVSS~v3.1~\cite{first2019cvss31}.
\end{itemize}
\end{boxinfo}

\begin{boxinfo}
\textbf{Transition.} La Partie~II commence par ce que l'on apprend sans
toucher aux systèmes~: la reconnaissance passive (Chapitre~\ref{chap:osint}).
\end{boxinfo}
}{}

    \part{Reconnaître}
    \label{part:reconnaitre}

        \IfFileExists{partie2/P02_C05_osint.tex}{% ============================================================================
% Fichier  : partie2/P02_C05_osint.tex
% Titre    : Reconnaissance Passive et OSINT
% Auteur   : MINKA MI NGUIDJOÏ Thierry Emmanuel
% Version  : 1.0 -- novembre 2026
% ============================================================================
% Sources : PTES (collecte de renseignements) ; NIST SP 800-115 ; RFC 6962
% (Certificate Transparency) ; MITRE ATT&CK (tactique Reconnaissance).
% Résultats de l'Opération LOGONE : cas fictif.
% ============================================================================

\chapter{Reconnaissance Passive et OSINT}
\label{chap:osint}
\chaptermark{Reconnaissance passive et OSINT}

\epigraph{%
    « L'attaquant connaît souvent votre système d'information\\
    mieux que votre inventaire. »%
}{Principe du cours}

\niveauChapitre{Intermédiaire}{Chapitres~\ref{chap:engagement}
et~\ref{chap:methodes}.}

\begin{boxobjectifs}
\begin{itemize}
    \item Distinguer reconnaissance passive et active.
    \item Exploiter les sources ouvertes~: DNS, certificats publics,
          annonces, réseaux professionnels, dépôts de code.
    \item Construire la surface d'attaque externe d'une organisation.
    \item Rester dans le cadre légal et éthique de la collecte.
    \item Transformer les résultats en mesures de réduction de l'exposition.
\end{itemize}
\end{boxobjectifs}

\bigskip

% ============================================================================
\section*{Question directrice : que voyait-on de LiZbiZ depuis Internet, sans rien toucher~?}
% ============================================================================

Les 19 et 20~octobre, M.~TAMO ne se connecte à aucun système de LiZbiZ. Il
lit ce que le monde peut lire~: les enregistrements DNS, les journaux
publics de certificats, les offres d'emploi, les profils professionnels des
salariés. À la fin du deuxième jour, il connaît trois serveurs exposés que
l'inventaire de M.~KOUNOU ignore.

\medskip
\textbf{Plan du chapitre.}
\begin{enumerate}
    \item Passive ou active (section~\ref{sec:c05-passive}).
    \item Les sources ouvertes (section~\ref{sec:c05-sources}).
    \item La surface d'attaque externe (section~\ref{sec:c05-surface}).
    \item Les limites de la collecte (section~\ref{sec:c05-limites}).
    \item Réduire l'exposition (section~\ref{sec:c05-reduire}).
\end{enumerate}

\separateur

% ============================================================================
\section{Passive ou Active}
\label{sec:c05-passive}
% ============================================================================

\begin{boxdef}
\termecle{Reconnaissance passive}~: collecte d'informations sans interagir
avec les systèmes de la cible (sources publiques, services tiers).
\termecle{Reconnaissance active}~: collecte par interaction directe
(résolution de noms auprès des serveurs de la cible, balayage). Le
renseignement en sources ouvertes, ou \termecle{OSINT}, rassemble les
informations publiquement accessibles.
\end{boxdef}

La reconnaissance passive ne laisse presque aucune trace chez la cible~:
elle est aussi ce que fait un attaquant avant d'attaquer. ATT\&CK lui
consacre une tactique entière, la \emph{reconnaissance}~\cite{mitreattack}.

% ============================================================================
\section{Les Sources Ouvertes}
\label{sec:c05-sources}
% ============================================================================

\begin{tableaucours}{L{3.6cm}L{4.2cm}L{4.4cm}}{%
    \tableauHeader{Source} &
    \tableauHeader{Ce qu'elle révèle} &
    \tableauHeader{Exemple chez LiZbiZ}
}
Enregistrements DNS & Serveurs, messagerie, services & Serveurs de
messagerie hébergés en nuage \\
Journaux de transparence des certificats & Noms de domaine pour lesquels un
certificat a été émis & 14~sous-domaines de lizbiz.cm \\
Registres d'adresses & Plages attribuées, hébergeur & La plage
198.51.100.0/27 \\
Offres d'emploi & Technologies utilisées & « Administrateur de notre
pare-feu et de notre passerelle d'accès distant » \\
Réseaux professionnels & Organigramme, format des adresses &
prenom.nom@lizbiz.cm \\
Dépôts de code publics & Configurations, secrets oubliés & Aucun secret
trouvé \\
\end{tableaucours}
\vspace{-0.8em}
\begin{center}{\small\itshape Tableau~5.1. Sources ouvertes et apports}\end{center}

Les journaux de \termecle{transparence des certificats}, définis par la RFC~6962,
enregistrent publiquement les certificats émis par les autorités
participantes~\cite{rfc6962}. Ils servent à détecter les certificats émis
par erreur~; ils révèlent aussi les noms de serveurs qu'une organisation
croyait discrets.

% ============================================================================
\section{La Surface d'Attaque Externe}
\label{sec:c05-surface}
% ============================================================================

\begin{boxscenario}{Opération LOGONE. Ce que l'inventaire ignorait}
\textbf{Constatations} (20~octobre). Sur 14~sous-domaines trouvés dans les
journaux de certificats, 11 figurent à l'inventaire. Trois n'y sont pas~:
« ancien-portail.lizbiz.cm » (ancienne version du portail de suivi, encore
en ligne), « test-api.lizbiz.cm » (environnement de test) et
« vpn2.lizbiz.cm » (ancienne passerelle d'accès distant, non retirée).

\medskip
\textbf{Analyse} (confiance moyenne). Les trois serveurs oubliés sont
précisément ceux qui ne reçoivent ni correctifs ni second facteur. Leur
résolution et leur appartenance à la plage du périmètre restent à vérifier
avant tout test actif (Chapitre~\ref{chap:engagement}).
\end{boxscenario}

\begin{boxobjet}{T-03 · Ancienne passerelle d'accès distant non retirée}
\textbf{Description~:} « vpn2.lizbiz.cm » répond sur Internet et accepte une
authentification par mot de passe seul. \textbf{Preuve~:} entrée du journal
de certificats, page de connexion (capture du 21~octobre).
\textbf{Gravité~:} élevée. \textbf{Impact~:} contournement du second facteur
de la passerelle principale. \textbf{Recommandation~:} retirer le service,
révoquer son certificat, inscrire le retrait au processus de fin de vie.
\textbf{Détection~:} inventaire externe périodique.
\end{boxobjet}

% ============================================================================
\section{Les Limites de la Collecte}
\label{sec:c05-limites}
% ============================================================================

\begin{boxalert}
\textbf{Violation typique.} Se connecter à un compte trouvé dans une fuite
publique « pour vérifier qu'il fonctionne ».

\smallskip
\textbf{Règle~:} utiliser un identifiant volé, même publié, est un accès
non autorisé s'il n'est pas prévu par les règles d'engagement. La collecte
passive s'arrête à la lecture~; on signale la fuite au client.
\end{boxalert}

La collecte sur les personnes (salariés, dirigeants) touche à des données
personnelles~: on se limite à ce qui sert le test (format des adresses,
fonctions exposées), on ne constitue pas de dossier individuel, et on
détruit les données à la fin de la mission (Chapitre~\ref{chap:droit}).

% ============================================================================
\section{Réduire l'Exposition}
\label{sec:c05-reduire}
% ============================================================================

\begin{boxprocedure}
\textbf{Mesurer et réduire sa propre exposition}
\begin{enumerate}
    \item Surveiller les journaux de certificats pour ses domaines.
    \item Comparer chaque trimestre la surface externe observée à
          l'inventaire.
    \item Retirer les services anciens et les environnements de test exposés.
    \item Relire les offres d'emploi et les communications techniques
          publiques.
    \item Rechercher ses domaines dans les dépôts de code publics.
\end{enumerate}
\end{boxprocedure}

\begin{boxmetier}{RSSI}
\textbf{Faire faire l'OSINT de sa propre organisation.} Le RSSI commande
chaque trimestre une reconnaissance passive de son organisation~: ce qu'un
attaquant voit gratuitement, il doit le voir avant lui.
\end{boxmetier}

\separateur

\begin{boxresume}
\textbf{Ce qu'il faut retenir de ce chapitre~:}
\begin{itemize}
    \item \textbf{Passive}~: sans toucher la cible~; \textbf{active}~: en
          interagissant.
    \item \textbf{DNS, certificats, registres, annonces, réseaux, dépôts}~:
          les sources ouvertes.
    \item \textbf{La surface externe observée} se compare à l'inventaire~;
          les écarts sont souvent les failles.
    \item \textbf{La collecte s'arrête à la lecture}~; données personnelles
          minimisées.
    \item \textbf{Faire l'OSINT de soi-même}, avant l'attaquant.
\end{itemize}
\end{boxresume}

\begin{boxinfo}
\textbf{Points de contrôle de l'élève.} Je sais~: distinguer reconnaissance
passive et active~; citer six sources ouvertes~; expliquer la transparence
des certificats~; comparer surface observée et inventaire~; respecter les
limites de la collecte.
\end{boxinfo}

\bigskip

\begin{boxexo}[title={\ding{45}\quad Exercice~5.1. Passive ou active \niveaubase}]
Classez~: lire un journal de certificats~; résoudre un nom auprès du
serveur DNS de la cible~; consulter une offre d'emploi~; balayer un port.
\end{boxexo}

\begin{boxexo}[title={\ding{45}\quad Exercice~5.2. L'écart \niveaubase}]
Combien de sous-domaines trouvés ne figuraient pas à l'inventaire de
LiZbiZ~? Quelle part du total~?
\end{boxexo}

\begin{boxexo}[title={\ding{45}\quad Exercice~5.3. Les certificats \niveaubase}]
Pourquoi les journaux de transparence des certificats révèlent-ils des
serveurs « discrets »~?
\end{boxexo}

\begin{boxexo}[title={\ding{45}\quad Exercice~5.4. L'offre d'emploi \niveaumoyen}]
Réécrivez l'offre d'emploi de LiZbiZ pour qu'elle reste attractive sans
révéler ses technologies de sécurité.
\end{boxexo}

\begin{boxexo}[title={\ding{45}\quad Exercice~5.5. Le compte publié \niveaumoyen}]
Vous trouvez dans une fuite publique l'adresse et le mot de passe d'un
salarié de LiZbiZ. Que faites-vous, et que ne faites-vous pas~?
\end{boxexo}

\begin{boxexo}[title={\ding{45}\quad Exercice~5.6. Les personnes \niveaumoyen}]
Quelles informations sur les salariés le testeur peut-il collecter pour une
campagne d'hameçonnage simulée, et lesquelles doit-il s'interdire~?
\end{boxexo}

\begin{boxexo}[title={\ding{45}\quad Exercice~5.7. T-03 \niveaumoyen}]
Expliquez pourquoi l'ancienne passerelle annule l'effet du second facteur
de la passerelle principale.
\end{boxexo}

\begin{boxexo}[title={\ding{45}\quad Exercice~5.8. La procédure \niveaumoyen}]
Transformez la procédure « mesurer et réduire sa propre exposition » en
indicateur clé de risque avec des seuils (cours GES4062).
\end{boxexo}

\begin{boxexo}[title={\ding{45}\quad Exercice~5.9. L'ANPP \niveauexpert}]
Préparez le plan de reconnaissance passive de l'Opération MEFOU~: sources,
informations attendues, limites.
\end{boxexo}

\begin{boxexo}[title={\ding{45}\quad Exercice~5.10. ATT\&CK \niveauexpert}]
Pour la tactique « reconnaissance » d'ATT\&CK, proposez trois mesures
défensives qui la rendent plus difficile ou plus visible.
\end{boxexo}

\bigskip

\begin{boxinfo}
\textbf{Pour aller plus loin.}
\begin{itemize}
    \item RFC~6962~\cite{rfc6962}~; MITRE ATT\&CK~\cite{mitreattack}.
    \item PTES, collecte de renseignements~\cite{ptes}.
\end{itemize}
\end{boxinfo}

\begin{boxinfo}
\textbf{Transition.} La reconnaissance active commence par le balayage.
C'est l'objet du Chapitre~\ref{chap:balayage}.
\end{boxinfo}
}{}
        \IfFileExists{partie2/P02_C06_balayage.tex}{% ============================================================================
% Fichier  : partie2/P02_C06_balayage.tex
% Titre    : Balayage et Énumération
% Auteur   : MINKA MI NGUIDJOÏ Thierry Emmanuel
% Version  : 1.0 -- novembre 2026
% ============================================================================
% Sources : RFC 9293 (TCP) ; NIST SP 800-115 (découverte) ; documentation de
% l'outil Nmap. Résultats de l'Opération LOGONE : cas fictif.
% ============================================================================

\chapter{Balayage et Énumération}
\label{chap:balayage}
\chaptermark{Balayage et énumération}

\epigraph{%
    « Un port ouvert est une question~:\\
    qui répond, et pourquoi est-il là~? »%
}{Principe du cours}

\niveauChapitre{Intermédiaire}{Chapitre~\ref{chap:osint}~; cours SEC3266,
chapitres sur les réseaux.}

\begin{boxobjectifs}
\begin{itemize}
    \item Expliquer la découverte d'hôtes et le balayage de ports TCP et UDP.
    \item Interpréter les états d'un port~: ouvert, fermé, filtré.
    \item Identifier les services et leurs versions, et énumérer leurs
          informations.
    \item Conduire un balayage maîtrisé, sans perturber la production.
    \item Relier le balayage à sa détection par le défenseur.
\end{itemize}
\end{boxobjectifs}

\bigskip

% ============================================================================
\section*{Question directrice : combien de portes ouvertes comptait vraiment LiZbiZ~?}
% ============================================================================

Le 21~octobre, M.~TAMO balaye la plage externe de LiZbiZ~; le 26, depuis le
poste fourni, Mme~NGO~MBOCK balaye les réseaux de Garoua. Les résultats
mettent en regard ce qui devrait être ouvert, selon la matrice des flux du
cours SEC3266, et ce qui l'est.

\medskip
\textbf{Plan du chapitre.}
\begin{enumerate}
    \item Découvrir les hôtes (section~\ref{sec:c06-decouverte}).
    \item Balayer les ports (section~\ref{sec:c06-ports}).
    \item Identifier les services (section~\ref{sec:c06-services}).
    \item Balayer sans nuire (section~\ref{sec:c06-maitrise}).
    \item Ce que voit le défenseur (section~\ref{sec:c06-defenseur}).
\end{enumerate}

\separateur

% ============================================================================
\section{Découvrir les Hôtes}
\label{sec:c06-decouverte}
% ============================================================================

La découverte cherche les adresses qui répondent~: requêtes ICMP, requêtes
ARP sur le réseau local, tentatives de connexion TCP vers des ports
courants. Un hôte qui ne répond pas au « ping » n'est pas forcément éteint~:
un pare-feu peut filtrer l'ICMP.

\begin{boxdef}
\termecle{Énumération}~: étape qui, après la découverte, extrait d'un
service les informations qu'il livre (noms de partages, comptes, versions,
bannières, configuration).
\end{boxdef}

% ============================================================================
\section{Balayer les Ports}
\label{sec:c06-ports}
% ============================================================================

TCP établit une connexion en trois temps (SYN, SYN-ACK, ACK) défini par la
RFC~9293~\cite{rfc9293}. Le balayage exploite les réponses à un SYN.

\begin{tableaucours}{L{2.4cm}L{4.6cm}L{5.2cm}}{%
    \tableauHeader{État} &
    \tableauHeader{Réponse observée} &
    \tableauHeader{Signification}
}
Ouvert & SYN-ACK & Un service écoute \\
Fermé & RST & L'hôte répond, aucun service \\
Filtré & Aucune réponse ou message ICMP d'interdiction & Un filtrage
intercepte \\
\end{tableaucours}
\vspace{-0.8em}
\begin{center}{\small\itshape Tableau~6.1. Les états d'un port TCP}\end{center}

Le balayage UDP est plus lent et plus incertain~: un service UDP ouvert ne
répond souvent rien, ce qui le confond avec un port filtré. On le réserve
aux ports utiles (DNS, SNMP, NTP).

% ============================================================================
\section{Identifier les Services}
\label{sec:c06-services}
% ============================================================================

\begin{boxscenario}{Opération LOGONE. Prévu et constaté}
\textbf{Constatations.} Externe~: 7~hôtes actifs sur 32~adresses, 18~ports
ouverts, dont 4 non prévus par la matrice des flux (dont l'ancienne
passerelle, constat T-03). Interne, Garoua~: 214~hôtes actifs~; le
protocole de partage de fichiers SMBv1 est actif sur deux serveurs~; une
imprimante expose son interface d'administration avec le mot de passe par
défaut (constat déjà signalé en mars, cours SEC3266).

\medskip
\textbf{Analyse} (confiance élevée). Les écarts entre la matrice et la
réalité sont autant de constats~: un service non prévu n'est ni surveillé,
ni corrigé.
\end{boxscenario}

\begin{tableaucours}{L{2.0cm}L{3.0cm}L{7.2cm}}{%
    \tableauHeader{Port} &
    \tableauHeader{Service} &
    \tableauHeader{Information à énumérer}
}
22 & SSH & Version, méthodes d'authentification acceptées \\
80, 443 & HTTP, HTTPS & Serveur, application, certificat, pages \\
445 & SMB & Version du protocole, partages accessibles, signature \\
389 & LDAP & Annuaire~: domaine, possibilité d'interrogation anonyme \\
161 (UDP) & SNMP & Communauté par défaut, informations système \\
\end{tableaucours}
\vspace{-0.8em}
\begin{center}{\small\itshape Tableau~6.2. Ports courants et énumération}\end{center}

% ============================================================================
\section{Balayer sans Nuire}
\label{sec:c06-maitrise}
% ============================================================================

\begin{boxprocedure}
\textbf{Un balayage maîtrisé}
\begin{enumerate}
    \item Vérifier le périmètre et les exclusions (Chapitre~\ref{chap:engagement}).
    \item Commencer par la découverte, puis les ports courants, puis le reste.
    \item Limiter le débit sur les réseaux fragiles (automates, imprimantes,
          liaisons lentes).
    \item Éviter les scripts intrusifs sur la production sans accord.
    \item Enregistrer les résultats bruts, horodatés, dans le journal.
\end{enumerate}
\end{boxprocedure}

\begin{boxalert}
\textbf{Violation typique.} Lancer un balayage complet à plein débit sur
une liaison inter-sites lente.

\smallskip
\textbf{Règle~:} un balayage peut saturer un équipement ou une liaison et
interrompre un service~: c'est la conséquence que réprime l'article~66 si
elle n'est pas maîtrisée. On adapte le débit au réseau.
\end{boxalert}

% ============================================================================
\section{Ce que Voit le Défenseur}
\label{sec:c06-defenseur}
% ============================================================================

Un balayage est bruyant~: une source qui contacte de nombreux ports ou de
nombreux hôtes en peu de temps est un motif classique (cours SEC3266). Le
SOC de LiZbiZ a détecté le balayage interne en 40~minutes.

\begin{boxmetier}{Analyste SOC et équipe rouge}
\textbf{Corréler les balayages internes.} Une règle simple~: alerte si une
même source interne contacte plus de 100~ports ou 50~hôtes en dix minutes.
Elle aurait vu PC-GAR-07 en mars (SEC3266) et le testeur en octobre.
\end{boxmetier}

\separateur

\begin{boxresume}
\textbf{Ce qu'il faut retenir de ce chapitre~:}
\begin{itemize}
    \item \textbf{Découverte, balayage, énumération}~: trois étapes.
    \item \textbf{Ouvert, fermé, filtré}~: trois états d'un port TCP~; UDP est
          plus incertain.
    \item \textbf{Comparer la réalité à la matrice des flux}~: chaque écart
          est un constat.
    \item \textbf{Balayer sans nuire}~: périmètre, progressivité, débit,
          scripts prudents.
    \item \textbf{Le balayage est bruyant}~: le SOC doit le voir.
\end{itemize}
\end{boxresume}

\begin{boxinfo}
\textbf{Points de contrôle de l'élève.} Je sais~: expliquer la découverte
d'hôtes~; interpréter l'état d'un port~; choisir quoi énumérer selon le
service~; conduire un balayage maîtrisé~; écrire une règle de détection
simple.
\end{boxinfo}

\bigskip

\begin{boxexo}[title={\ding{45}\quad Exercice~6.1. Les états \niveaubase}]
Quel état attribuer à un port qui répond par RST~? par SYN-ACK~? qui ne
répond pas~?
\end{boxexo}

\begin{boxexo}[title={\ding{45}\quad Exercice~6.2. Le ping \niveaubase}]
Pourquoi un hôte qui ne répond pas au « ping » peut-il être actif~?
\end{boxexo}

\begin{boxexo}[title={\ding{45}\quad Exercice~6.3. Les écarts \niveaubase}]
Sur les 18~ports ouverts à l'externe, combien étaient prévus par la
matrice des flux~?
\end{boxexo}

\begin{boxexo}[title={\ding{45}\quad Exercice~6.4. UDP \niveaumoyen}]
Pourquoi un balayage UDP est-il plus lent et plus incertain qu'un balayage
TCP~?
\end{boxexo}

\begin{boxexo}[title={\ding{45}\quad Exercice~6.5. SMBv1 \niveaumoyen}]
Rédigez la fiche de constat « SMBv1 actif sur deux serveurs ».
\end{boxexo}

\begin{boxexo}[title={\ding{45}\quad Exercice~6.6. L'imprimante \niveaumoyen}]
L'imprimante avait déjà été signalée en mars. Que révèle sa présence en
octobre sur la gestion des corrections~? (cours GES4062)
\end{boxexo}

\begin{boxexo}[title={\ding{45}\quad Exercice~6.7. Le débit \niveaumoyen}]
Proposez un plan de balayage des 214~hôtes de Garoua qui ne sature pas la
liaison avec Douala.
\end{boxexo}

\begin{boxexo}[title={\ding{45}\quad Exercice~6.8. Détecter \niveaumoyen}]
Le SOC a mis 40~minutes à voir le balayage. Proposez deux améliorations
pour réduire ce délai.
\end{boxexo}

\begin{boxexo}[title={\ding{45}\quad Exercice~6.9. Énumération \niveauexpert}]
Pour chacun des ports du Tableau~6.2, dites quelle mesure défensive
réduit l'information livrée à un attaquant.
\end{boxexo}

\begin{boxexo}[title={\ding{45}\quad Exercice~6.10. Inventaire continu \niveauexpert}]
Concevez un processus de balayage interne mensuel, autorisé, qui alimente
l'inventaire et la matrice des flux de LiZbiZ.
\end{boxexo}

\bigskip

\begin{boxinfo}
\textbf{Pour aller plus loin.}
\begin{itemize}
    \item RFC~9293 (TCP)~\cite{rfc9293}~; NIST SP~800-115~\cite{nist800115}.
\end{itemize}
\end{boxinfo}

\begin{boxinfo}
\textbf{Transition.} Un service identifié doit être confronté aux failles
connues~: le Chapitre~\ref{chap:vulnerabilites} apprend à les analyser et
à les prioriser.
\end{boxinfo}
}{}
        \IfFileExists{partie2/P02_C07_vulnerabilites.tex}{% ============================================================================
% Fichier  : partie2/P02_C07_vulnerabilites.tex
% Titre    : Analyser et Prioriser les Vulnérabilités
% Auteur   : MINKA MI NGUIDJOÏ Thierry Emmanuel
% Version  : 1.0 -- novembre 2026
% ============================================================================
% Sources : FIRST, CVSS v3.1 et EPSS ; CISA, catalogue Known Exploited
% Vulnerabilities ; programme CVE (MITRE) ; NIST SP 800-40 Rev. 4.
% Résultats de l'Opération LOGONE : cas fictif.
% ============================================================================

\chapter{Analyser et Prioriser les Vulnérabilités}
\label{chap:vulnerabilites}
\chaptermark{Analyser et prioriser}

\epigraph{%
    « Un rapport de trois cents failles\\
    est un rapport qu'on ne lira pas. »%
}{Principe du cours}

\niveauChapitre{Intermédiaire}{Chapitres~\ref{chap:methodes}
et~\ref{chap:balayage}.}

\begin{boxobjectifs}
\begin{itemize}
    \item Utiliser un scanner de vulnérabilités et en connaître les limites.
    \item Écarter les faux positifs par la vérification.
    \item Comprendre les identifiants CVE et les catalogues de failles
          exploitées.
    \item Prioriser par la gravité, l'exploitation connue, l'exposition et
          la valeur de l'actif.
    \item Présenter une liste courte, ordonnée et justifiée.
\end{itemize}
\end{boxobjectifs}

\bigskip

% ============================================================================
\section*{Question directrice : par quoi commencer, parmi 312 résultats~?}
% ============================================================================

Le scanner lancé le 26~octobre sur les réseaux de Garoua produit 312
résultats. M.~KOUNOU n'a ni le temps ni l'équipe de tout corriger. Mme~NGO
MBOCK doit lui remettre une liste courte~: ce qui est vrai, ce qui est
dangereux, et dans quel ordre.

\medskip
\textbf{Plan du chapitre.}
\begin{enumerate}
    \item Les scanners de vulnérabilités (section~\ref{sec:c07-scanners}).
    \item Vérifier~: faux positifs et faux négatifs (section~\ref{sec:c07-verifier}).
    \item CVE et catalogues (section~\ref{sec:c07-cve}).
    \item Prioriser (section~\ref{sec:c07-prioriser}).
    \item Présenter (section~\ref{sec:c07-presenter}).
\end{enumerate}

\separateur

% ============================================================================
\section{Les Scanners de Vulnérabilités}
\label{sec:c07-scanners}
% ============================================================================

Un \termecle{scanner de vulnérabilités} compare les services et les
versions observés à une base de failles connues, et tente parfois des
vérifications. Il peut travailler \emph{sans authentification} (vue
d'attaquant externe) ou \emph{avec authentification} (il lit les versions
installées~: bien plus précis).

\begin{tableaucours}{L{3.6cm}L{8.6cm}}{%
    \tableauHeader{Limite} &
    \tableauHeader{Conséquence}
}
Déduction par la version & Une version corrigée par un rétroportage peut
être signalée vulnérable \\
Bases incomplètes & Les failles de configuration et de logique sont mal vues \\
Absence de contexte & Le scanner ne sait pas ce qui est critique pour le
métier \\
Bruit & Des centaines de résultats de faible intérêt \\
\end{tableaucours}
\vspace{-0.8em}
\begin{center}{\small\itshape Tableau~7.1. Limites des scanners}\end{center}

% ============================================================================
\section{Vérifier : Faux Positifs et Faux Négatifs}
\label{sec:c07-verifier}
% ============================================================================

\begin{boxdef}
\termecle{Faux positif}~: faille signalée qui n'existe pas.
\termecle{Faux négatif}~: faille présente que l'outil n'a pas signalée. Le
premier fait perdre du temps~; le second donne une fausse assurance.
\end{boxdef}

\begin{boxscenario}{Opération LOGONE. De 312 à 41}
\textbf{Constatations.} Sur 312~résultats bruts~: 158 sont des
informations sans gravité~; 96 sont des doublons (même faille sur plusieurs
hôtes, regroupée)~; 17 sont des faux positifs (versions corrigées par
rétroportage). Restent 41~vulnérabilités distinctes et vérifiées, dont 3
figurent au catalogue des failles exploitées.

\medskip
\textbf{Analyse} (confiance élevée). Le travail du testeur n'est pas de
recopier le scanner, mais de transformer ses résultats en constats vrais et
ordonnés.
\end{boxscenario}

% ============================================================================
\section{CVE et Catalogues}
\label{sec:c07-cve}
% ============================================================================

\begin{tableaucours}{L{3.4cm}L{8.8cm}}{%
    \tableauHeader{Référence} &
    \tableauHeader{Rôle}
}
CVE & Identifiant public unique d'une vulnérabilité (programme CVE, MITRE)~\cite{mitrecve} \\
CVSS & Gravité technique, de 0 à 10 (FIRST)~\cite{first2019cvss31} \\
EPSS & Probabilité estimée d'exploitation dans les 30~jours (FIRST)~\cite{firstepss} \\
Catalogue KEV & Liste des vulnérabilités dont l'exploitation est avérée (CISA)~\cite{cisakev} \\
\end{tableaucours}
\vspace{-0.8em}
\begin{center}{\small\itshape Tableau~7.2. Les références utiles à la priorisation}\end{center}

% ============================================================================
\section{Prioriser}
\label{sec:c07-prioriser}
% ============================================================================

\begin{boxprocedure}
\textbf{Une règle de priorisation (convention du cours)}
\begin{enumerate}
    \item \textbf{Priorité~1}~: exploitation avérée (catalogue KEV) et actif
          exposé à Internet ou critique~; ou CVSS d'au moins 9 sur un actif
          exposé.
    \item \textbf{Priorité~2}~: CVSS d'au moins 7 sur un actif interne
          critique, ou probabilité d'exploitation élevée.
    \item \textbf{Priorité~3}~: le reste, traité dans le cycle normal des
          correctifs (cours SEC3266, Chapitre sur les systèmes).
\end{enumerate}
\end{boxprocedure}

\begin{boxalert}
\textbf{Violation typique.} Trier uniquement par le score CVSS décroissant.

\smallskip
\textbf{Règle~:} une faille de score 7,5 exploitée activement sur la
passerelle passe avant une faille de score 9,1 sur un poste isolé, sans
exploitation connue. On croise gravité, exploitation, exposition et valeur.
\end{boxalert}

% ============================================================================
\section{Présenter}
\label{sec:c07-presenter}
% ============================================================================

\begin{boxobjet}{T-06 · SMBv1 actif sur deux serveurs de Garoua}
\textbf{Description~:} le protocole SMBv1, obsolète, est accepté par
FS-02 et APP-01. \textbf{Preuve~:} négociation du protocole (relevé du
26~octobre). \textbf{Gravité~:} élevée (vulnérabilités connues de SMBv1,
exploitées par le passé à grande échelle). \textbf{Impact~:} exécution de
code à distance si un correctif manque. \textbf{Recommandation~:}
désactiver SMBv1, vérifier les clients qui en dépendent.
\textbf{Détection~:} journaliser les négociations SMBv1.
\end{boxobjet}

\begin{boxmetier}{RSSI}
\textbf{Demander peu, mais juste.} Le RSSI attend du rapport une liste de
priorité~1 courte et argumentée, qu'il peut faire corriger dans le mois, et
une annexe complète pour le cycle normal.
\end{boxmetier}

\separateur

\begin{boxresume}
\textbf{Ce qu'il faut retenir de ce chapitre~:}
\begin{itemize}
    \item \textbf{Un scanner} compare des versions à une base~; avec
          authentification, il est plus précis.
    \item \textbf{On vérifie}~: faux positifs écartés, faux négatifs
          recherchés à la main.
    \item \textbf{CVE, CVSS, EPSS, KEV}~: identifier, coter, estimer,
          savoir ce qui est exploité.
    \item \textbf{On priorise} en croisant gravité, exploitation, exposition
          et valeur.
    \item \textbf{On présente} une liste courte de priorité~1 et une annexe.
\end{itemize}
\end{boxresume}

\begin{boxinfo}
\textbf{Points de contrôle de l'élève.} Je sais~: dire les limites d'un
scanner~; distinguer faux positif et faux négatif~; expliquer CVE, CVSS,
EPSS et KEV~; appliquer une règle de priorisation~; rédiger un constat.
\end{boxinfo}

\bigskip

\begin{boxexo}[title={\ding{45}\quad Exercice~7.1. Les chiffres \niveaubase}]
Vérifiez que 312 résultats moins les informations, les doublons et les faux
positifs laissent 41~vulnérabilités.
\end{boxexo}

\begin{boxexo}[title={\ding{45}\quad Exercice~7.2. Faux positif \niveaubase}]
Pourquoi une version corrigée par rétroportage est-elle souvent signalée à
tort~?
\end{boxexo}

\begin{boxexo}[title={\ding{45}\quad Exercice~7.3. Les références \niveaubase}]
Associez à chaque question la bonne référence (CVE, CVSS, EPSS, KEV)~:
« est-elle exploitée~? »~; « quelle gravité~? »~; « quel nom~? »~; « quelle
probabilité d'être exploitée bientôt~? »
\end{boxexo}

\begin{boxexo}[title={\ding{45}\quad Exercice~7.4. Prioriser \niveaumoyen}]
Appliquez la règle du cours à~: CVSS 7,5 au catalogue KEV sur la
passerelle~; CVSS 9,1 sans exploitation sur un poste isolé~; CVSS 5,3 sur le
portail.
\end{boxexo}

\begin{boxexo}[title={\ding{45}\quad Exercice~7.5. Authentifié \niveaumoyen}]
Pourquoi un balayage authentifié réduit-il les faux positifs~? Quelles
précautions exige-t-il~?
\end{boxexo}

\begin{boxexo}[title={\ding{45}\quad Exercice~7.6. Faux négatifs \niveaumoyen}]
Citez trois types de failles qu'un scanner voit mal, et comment le testeur
les recherche.
\end{boxexo}

\begin{boxexo}[title={\ding{45}\quad Exercice~7.7. T-06 \niveaumoyen}]
Quels clients pourraient encore dépendre de SMBv1 chez LiZbiZ~? Comment le
vérifier avant de le désactiver~?
\end{boxexo}

\begin{boxexo}[title={\ding{45}\quad Exercice~7.8. Regrouper \niveaumoyen}]
Pourquoi regrouper une même faille présente sur 40~hôtes en un seul
constat~? Que faut-il indiquer pour ne rien perdre~?
\end{boxexo}

\begin{boxexo}[title={\ding{45}\quad Exercice~7.9. Une politique \niveauexpert}]
Rédigez la politique de priorisation des vulnérabilités de LiZbiZ, avec les
délais associés à chaque priorité.
\end{boxexo}

\begin{boxexo}[title={\ding{45}\quad Exercice~7.10. Mesurer \niveauexpert}]
Proposez trois indicateurs du processus de gestion des vulnérabilités et
leurs seuils.
\end{boxexo}

\bigskip

\begin{boxinfo}
\textbf{Pour aller plus loin.}
\begin{itemize}
    \item FIRST, CVSS~\cite{first2019cvss31} et EPSS~\cite{firstepss}~;
          CISA, catalogue KEV~\cite{cisakev}.
    \item NIST SP~800-40 Rev.~4~\cite{nist80040}.
\end{itemize}
\end{boxinfo}

\begin{boxinfo}
\textbf{Transition.} Une faille vérifiée se démontre par l'exploitation,
dans un cadre maîtrisé. La Partie~III commence par ses principes
(Chapitre~\ref{chap:exploitation}).
\end{boxinfo}
}{}

    \part{Systèmes et Réseaux}
    \label{part:systemes}

        \IfFileExists{partie3/P03_C08_exploitation.tex}{% ============================================================================
% Fichier  : partie3/P03_C08_exploitation.tex
% Titre    : L'Exploitation : Principes et Cadres
% Auteur   : MINKA MI NGUIDJOÏ Thierry Emmanuel
% Version  : 1.0 -- novembre 2026
% ============================================================================
% Sources : NIST SP 800-115 (attaque) ; PTES (exploitation) ; MITRE ATT&CK.
% Démonstration de l'Opération LOGONE : cas fictif, aucun code d'attaque.
% ============================================================================

\chapter{L'Exploitation : Principes et Cadres}
\label{chap:exploitation}
\chaptermark{L'exploitation}

\epigraph{%
    « Démontrer, pas détruire. »%
}{Principe du cours}

\niveauChapitre{Intermédiaire}{Chapitres~\ref{chap:engagement}
et~\ref{chap:vulnerabilites}.}

\begin{boxobjectifs}
\begin{itemize}
    \item Distinguer vulnérabilité, exploit, charge utile et preuve de
          concept.
    \item Choisir de démontrer ou de ne pas démontrer une faille.
    \item Évaluer le risque d'une exploitation pour la production.
    \item Utiliser un cadre d'exploitation de manière maîtrisée.
    \item Produire une preuve minimale, vérifiable et sans dommage.
\end{itemize}
\end{boxobjectifs}

\bigskip

% ============================================================================
\section*{Question directrice : faut-il exploiter l'ancien portail pour prouver qu'il est vulnérable~?}
% ============================================================================

L'ancien portail de suivi, oublié depuis 2024, utilise une version
d'application connue pour une faille d'exécution de code à distance. Le
scanner le dit. M.~KOUNOU répond que « ce serveur ne sert plus à rien ».
M.~BIKORO décide de démontrer la faille, avec la plus petite preuve
possible. Pourquoi, et comment~?

\medskip
\textbf{Plan du chapitre.}
\begin{enumerate}
    \item Le vocabulaire de l'exploitation (section~\ref{sec:c08-vocabulaire}).
    \item Démontrer ou non (section~\ref{sec:c08-decider}).
    \item Le risque pour la production (section~\ref{sec:c08-risque}).
    \item Les cadres d'exploitation (section~\ref{sec:c08-cadres}).
    \item La preuve minimale (section~\ref{sec:c08-preuve}).
\end{enumerate}

\separateur

% ============================================================================
\section{Le Vocabulaire de l'Exploitation}
\label{sec:c08-vocabulaire}
% ============================================================================

\begin{boxdef}
\termecle{Exploit}~: moyen (programme, requête, séquence d'actions) qui
tire parti d'une vulnérabilité pour obtenir un effet non prévu.
\termecle{Charge utile}~: ce que l'exploit exécute une fois la faille
déclenchée. \termecle{Preuve de concept}~: démonstration minimale qu'une
faille est exploitable, sans en tirer tout le parti possible.
\end{boxdef}

\begin{tableaucours}{L{3.4cm}L{8.8cm}}{%
    \tableauHeader{Effet recherché} &
    \tableauHeader{Exemple de preuve minimale}
}
Exécution de code & Afficher l'identité du compte d'exécution et le nom de
la machine \\
Lecture de fichier & Lire un fichier sans valeur (version, fichier de test
déposé par accord) \\
Contournement d'authentification & Afficher la page d'accueil d'un compte
de test fourni \\
Élévation de privilèges & Afficher l'appartenance au groupe obtenu \\
\end{tableaucours}
\vspace{-0.8em}
\begin{center}{\small\itshape Tableau~8.1. Effets et preuves minimales}\end{center}

% ============================================================================
\section{Démontrer ou Non}
\label{sec:c08-decider}
% ============================================================================

Démontrer une faille la rend indiscutable et en mesure l'impact réel~; ne
pas la démontrer évite tout risque pour la production. La décision dépend
de l'objectif, des règles d'engagement et du risque.

\begin{boxprocedure}
\textbf{Décider d'exploiter}
\begin{enumerate}
    \item L'exploitation est-elle autorisée pour cette cible~?
    \item L'exploit est-il connu, stable, compris~? Sait-on ce qu'il fait~?
    \item Quel est le pire effet possible (arrêt, corruption)~? Est-il
          acceptable~?
    \item Le client est-il prévenu, la fenêtre adaptée, une sauvegarde
          faite~?
    \item La preuve minimale suffit-elle à l'objectif~?
\end{enumerate}
\end{boxprocedure}

\begin{boxalert}
\textbf{Violation typique.} Lancer un exploit téléchargé sans l'avoir lu.

\smallskip
\textbf{Règle~:} un exploit d'origine inconnue peut contenir une porte
dérobée, détruire des données ou faire tomber le service. On n'emploie sur
la production que du code dont on connaît la provenance et l'effet, après
l'avoir essayé dans un laboratoire (Chapitre~\ref{chap:labo}).
\end{boxalert}

% ============================================================================
\section{Le Risque pour la Production}
\label{sec:c08-risque}
% ============================================================================

\begin{tableaucours}{L{3.4cm}L{4.6cm}L{4.2cm}}{%
    \tableauHeader{Type de faille} &
    \tableauHeader{Risque de l'exploitation} &
    \tableauHeader{Précaution}
}
Dépassement de mémoire & Arrêt du service, instabilité & Laboratoire
d'abord~; fenêtre creuse \\
Injection dans une base & Modification de données & Requêtes en lecture
seule \\
Exécution de commande & Commande destructrice par erreur & Commandes
d'identification seulement \\
Déni de service & Interruption certaine & Exclu par défaut \\
\end{tableaucours}
\vspace{-0.8em}
\begin{center}{\small\itshape Tableau~8.2. Risques et précautions}\end{center}

% ============================================================================
\section{Les Cadres d'Exploitation}
\label{sec:c08-cadres}
% ============================================================================

Les \termecle{cadres d'exploitation} réunissent des modules d'exploits, de
charges utiles et d'outils de post-exploitation, documentés et maintenus.
Ils réduisent le risque des exploits isolés, mais ne dispensent ni de
comprendre le module, ni de vérifier qu'il est adapté à la version ciblée.
Le cours ne détaille pas leur emploi~: il en retient l'usage
professionnel~: modules connus, options vérifiées, journal de chaque
lancement.

% ============================================================================
\section{La Preuve Minimale}
\label{sec:c08-preuve}
% ============================================================================

\begin{boxscenario}{Opération LOGONE. L'ancien portail}
\textbf{Constatation.} Le 22~octobre à 14~h~05, après accord de M.~KOUNOU et
sauvegarde du serveur, M.~TAMO déclenche la faille de l'ancien portail avec
un module connu. La preuve~: l'affichage de l'identité du compte de service
web et du nom de la machine. Aucune autre commande n'est exécutée~; la
session est fermée à 14~h~08.

\medskip
\textbf{Analyse} (confiance élevée). « Ce serveur ne sert plus à rien »~:
mais il est en ligne, dans la DMZ, et le compte web peut lire la
configuration de connexion à la base du portail actuel. La démonstration a
transformé une alerte du scanner en constat incontestable.
\end{boxscenario}

\begin{boxobjet}{T-09 · Exécution de code sur l'ancien portail de suivi}
\textbf{Description~:} « ancien-portail.lizbiz.cm » exécute une version
d'application vulnérable à une exécution de code à distance.
\textbf{Preuve~:} capture du 22~octobre, 14~h~06 (identité du compte, nom
de la machine). \textbf{Gravité~:} élevée dans le contexte (serveur exposé).
\textbf{Recommandation~:} retirer le serveur~; à défaut, le mettre à jour et
l'isoler. \textbf{Détection~:} processus inhabituel lancé par le compte web.
\end{boxobjet}

\begin{boxmetier}{Pentesteur}
\textbf{Une commande, une preuve, on ferme.} Le pentesteur prépare à
l'avance la commande de preuve, l'exécute une fois, capture, ferme la
session et note l'heure. Explorer « tant qu'on y est » n'est pas une preuve,
c'est un risque.
\end{boxmetier}

\separateur

\begin{boxresume}
\textbf{Ce qu'il faut retenir de ce chapitre~:}
\begin{itemize}
    \item \textbf{Vulnérabilité, exploit, charge utile, preuve de concept}~:
          quatre notions distinctes.
    \item \textbf{On décide d'exploiter} selon l'autorisation, la stabilité,
          le pire effet et l'objectif.
    \item \textbf{Jamais d'exploit inconnu} sur la production~; laboratoire
          d'abord.
    \item \textbf{Les cadres d'exploitation} réduisent le risque sans
          dispenser de comprendre.
    \item \textbf{Une preuve minimale} suffit~: une commande, une capture,
          on ferme.
\end{itemize}
\end{boxresume}

\begin{boxinfo}
\textbf{Points de contrôle de l'élève.} Je sais~: définir exploit, charge
utile et preuve de concept~; dérouler la décision d'exploiter~; associer
une précaution à un type de faille~; choisir une preuve minimale.
\end{boxinfo}

\bigskip

\begin{boxexo}[title={\ding{45}\quad Exercice~8.1. Vocabulaire \niveaubase}]
Dans l'épisode de l'ancien portail, identifiez la vulnérabilité, l'exploit,
la charge utile et la preuve.
\end{boxexo}

\begin{boxexo}[title={\ding{45}\quad Exercice~8.2. Preuve minimale \niveaubase}]
Proposez une preuve minimale pour une lecture de fichier arbitraire sur un
serveur web.
\end{boxexo}

\begin{boxexo}[title={\ding{45}\quad Exercice~8.3. Déni de service \niveaubase}]
Pourquoi le déni de service est-il exclu par défaut~? Quel article de la loi
n°~2010/012 le concerne~?
\end{boxexo}

\begin{boxexo}[title={\ding{45}\quad Exercice~8.4. Décider \niveaumoyen}]
Appliquez la procédure de décision à une faille de dépassement de mémoire
sur le serveur de base de données de la facturation.
\end{boxexo}

\begin{boxexo}[title={\ding{45}\quad Exercice~8.5. L'exploit inconnu \niveaumoyen}]
Citez trois dangers d'un exploit téléchargé sur un site non vérifié.
\end{boxexo}

\begin{boxexo}[title={\ding{45}\quad Exercice~8.6. « Ne sert plus à rien » \niveaumoyen}]
Comment répondre à M.~KOUNOU, qui juge l'ancien portail sans importance~?
\end{boxexo}

\begin{boxexo}[title={\ding{45}\quad Exercice~8.7. La détection \niveaumoyen}]
Quels journaux auraient permis de détecter l'exploitation de l'ancien
portail~?
\end{boxexo}

\begin{boxexo}[title={\ding{45}\quad Exercice~8.8. Ne pas démontrer \niveaumoyen}]
Donnez deux situations où le testeur doit renoncer à démontrer une faille,
et comment il la rapporte alors.
\end{boxexo}

\begin{boxexo}[title={\ding{45}\quad Exercice~8.9. Une procédure \niveauexpert}]
Rédigez la procédure d'exploitation contrôlée de NYONG SÉCURITÉ~: accord,
préparation, exécution, preuve, fermeture, journal.
\end{boxexo}

\begin{boxexo}[title={\ding{45}\quad Exercice~8.10. Laboratoire \niveauexpert}]
Décrivez le laboratoire qui permet de valider un module d'exploitation avant
son emploi chez un client.
\end{boxexo}

\bigskip

\begin{boxinfo}
\textbf{Pour aller plus loin.}
\begin{itemize}
    \item NIST SP~800-115~\cite{nist800115}~; PTES~\cite{ptes}.
\end{itemize}
\end{boxinfo}

\begin{boxinfo}
\textbf{Transition.} Sur le réseau interne, l'attaquant cherche des secrets
qui circulent~: le Chapitre~\ref{chap:reseau} traite des attaques réseau et
de l'authentification.
\end{boxinfo}
}{}
        \IfFileExists{partie3/P03_C09_reseau.tex}{% ============================================================================
% Fichier  : partie3/P03_C09_reseau.tex
% Titre    : Attaques Réseau et Authentification
% Auteur   : MINKA MI NGUIDJOÏ Thierry Emmanuel
% Version  : 1.0 -- novembre 2026
% ============================================================================
% Sources : RFC 4795 (LLMNR) ; RFC 4120 (Kerberos) ; documentation publique
% de Microsoft sur NTLM et la signature SMB ; MITRE ATT&CK (T1557). Résultats
% de l'Opération LOGONE : cas fictif.
% ============================================================================

\chapter{Attaques Réseau et Authentification}
\label{chap:reseau}
\chaptermark{Attaques réseau et authentification}

\epigraph{%
    « Un réseau qui crie ses questions à tout le monde\\
    finit par recevoir une mauvaise réponse. »%
}{Principe du cours}

\niveauChapitre{Intermédiaire à expert}{Chapitre~\ref{chap:exploitation}~;
cours SEC3266, chapitres sur les réseaux et les identités.}

\begin{boxobjectifs}
\begin{itemize}
    \item Expliquer les protocoles d'authentification d'un réseau Windows
          (NTLM, Kerberos).
    \item Comprendre l'empoisonnement des résolutions de noms locales.
    \item Comprendre le relais d'authentification et ses conditions.
    \item Relier chaque attaque à sa détection et à sa parade.
    \item Rédiger les constats correspondants.
\end{itemize}
\end{boxobjectifs}

\bigskip

% ============================================================================
\section*{Question directrice : comment un poste ordinaire peut-il recueillir des secrets sans rien attaquer~?}
% ============================================================================

Le 26~octobre, depuis le poste fourni, Mme~NGO~MBOCK observe le trafic
local de Garoua. Des postes demandent à tout le réseau « qui est
srv-imprim~? », un serveur qui n'existe plus. Une machine qui répondrait
« c'est moi » recevrait leur tentative d'authentification. Ce chapitre
explique pourquoi, et comment fermer cette porte.

\medskip
\textbf{Plan du chapitre.}
\begin{enumerate}
    \item NTLM et Kerberos (section~\ref{sec:c09-protocoles}).
    \item L'empoisonnement des résolutions locales (section~\ref{sec:c09-llmnr}).
    \item Le relais d'authentification (section~\ref{sec:c09-relais}).
    \item Détecter et corriger (section~\ref{sec:c09-parades}).
\end{enumerate}

\separateur

% ============================================================================
\section{NTLM et Kerberos}
\label{sec:c09-protocoles}
% ============================================================================

\begin{tableaucours}{L{2.6cm}L{4.8cm}L{4.8cm}}{%
    \tableauHeader{Protocole} &
    \tableauHeader{Principe} &
    \tableauHeader{Faiblesse exploitable}
}
NTLM & Défi et réponse calculée à partir du haché du mot de passe & La
réponse peut être relayée ou cassée hors ligne si le mot de passe est
faible \\
Kerberos & Tickets délivrés par un centre de distribution de clés
(RFC~4120) & Tickets de service chiffrés avec le secret du compte de
service (Chapitre~\ref{chap:ad}) \\
\end{tableaucours}
\vspace{-0.8em}
\begin{center}{\small\itshape Tableau~9.1. Deux protocoles d'authentification}\end{center}

Kerberos, défini par la RFC~4120, est le protocole par défaut d'un domaine
Active Directory~\cite{rfc4120}~; NTLM subsiste pour la compatibilité, et
c'est souvent lui que vise l'attaquant.

% ============================================================================
\section{L'Empoisonnement des Résolutions Locales}
\label{sec:c09-llmnr}
% ============================================================================

Quand un nom ne se résout pas par le DNS, Windows peut interroger tout le
réseau local par LLMNR (RFC~4795) ou NBT-NS~\cite{rfc4795}. N'importe quelle
machine peut répondre. Un attaquant présent sur le réseau répond à ces
questions et reçoit en retour une tentative d'authentification NTLM.

\begin{boxdef}
\termecle{Empoisonnement de résolution}~: réponse frauduleuse à une
requête de résolution de nom diffusée sur le réseau local, qui détourne la
connexion de la victime vers l'attaquant (ATT\&CK, T1557).
\end{boxdef}

\begin{boxscenario}{Opération LOGONE. Les questions sans réponse}
\textbf{Constatation.} En une heure d'écoute passive, 37~requêtes LLMNR
pour des noms inexistants (dont « srv-imprim », retiré en 2025), émises par
22~postes. Avec l'accord de M.~KOUNOU, une réponse contrôlée de dix minutes
recueille 6~réponses NTLM de 4~comptes~; aucune n'est cassée, les mots de
passe concernés étant longs.

\medskip
\textbf{Analyse} (confiance élevée). La faille est structurelle (LLMNR et
NBT-NS actifs)~; son impact dépend de la force des mots de passe et de la
possibilité de relais.
\end{boxscenario}

% ============================================================================
\section{Le Relais d'Authentification}
\label{sec:c09-relais}
% ============================================================================

Au lieu de casser la réponse NTLM, l'attaquant peut la \emph{relayer}
aussitôt vers un autre serveur, et s'y authentifier comme la victime. Le
relais est possible lorsque le serveur cible n'exige pas la
\termecle{signature} des échanges (SMB) ou la liaison au canal (HTTP,
LDAP).

\begin{tableaucours}{L{4.0cm}L{8.2cm}}{%
    \tableauHeader{Condition} &
    \tableauHeader{Situation chez LiZbiZ}
}
Requêtes diffusées détournables & LLMNR et NBT-NS actifs \\
Serveurs sans signature SMB exigée & 31 serveurs sur 34 à Garoua \\
Comptes à privilèges qui s'authentifient & Comptes d'administration
utilisés depuis des postes ordinaires \\
\end{tableaucours}
\vspace{-0.8em}
\begin{center}{\small\itshape Tableau~9.2. Les conditions du relais (Opération LOGONE)}\end{center}

\begin{boxalert}
\textbf{Violation typique.} Relayer une authentification vers un serveur de
production sans accord.

\smallskip
\textbf{Règle~:} le relais crée une session réelle au nom d'un salarié~:
il n'est démontré que vers une cible et un compte convenus, et la session
est fermée aussitôt la preuve faite.
\end{boxalert}

% ============================================================================
\section{Détecter et Corriger}
\label{sec:c09-parades}
% ============================================================================

\begin{tableaucours}{L{4.4cm}L{7.8cm}}{%
    \tableauHeader{Parade} &
    \tableauHeader{Effet}
}
Désactiver LLMNR et NBT-NS & Supprime les requêtes diffusées détournables \\
Exiger la signature SMB & Empêche le relais vers SMB \\
Protection étendue de l'authentification (LDAP, HTTP) & Lie
l'authentification au canal \\
Réduire NTLM au profit de Kerberos & Réduit la surface \\
Comptes d'administration séparés, bastion & Les secrets d'administration ne
circulent pas depuis les postes \\
Détection & Répondant inconnu aux requêtes LLMNR~; authentifications NTLM
d'un compte depuis une machine inhabituelle \\
\end{tableaucours}
\vspace{-0.8em}
\begin{center}{\small\itshape Tableau~9.3. Parades et détection}\end{center}

\begin{boxobjet}{T-07 · Signature SMB non exigée et LLMNR actif}
\textbf{Description~:} 31~serveurs sur 34 n'exigent pas la signature SMB~;
LLMNR et NBT-NS sont actifs sur les postes. \textbf{Preuve~:} relevés du
26~octobre~; réponses NTLM recueillies (masquées). \textbf{Gravité~:}
moyenne (élevée si un compte à privilèges est capturé).
\textbf{Recommandation~:} désactiver LLMNR et NBT-NS par stratégie de
groupe~; exiger la signature SMB. \textbf{Détection~:} surveiller les
réponses LLMNR d'hôtes non autorisés.
\end{boxobjet}

\begin{boxmetier}{Analyste SOC et équipe rouge}
\textbf{Poser un piège.} Un poste « leurre » émet régulièrement une requête
LLMNR pour un nom inexistant~: toute réponse signale un empoisonneur sur le
réseau.
\end{boxmetier}

\separateur

\begin{boxresume}
\textbf{Ce qu'il faut retenir de ce chapitre~:}
\begin{itemize}
    \item \textbf{NTLM} se relaie ou se casse si le mot de passe est faible~;
          \textbf{Kerberos} est le protocole par défaut du domaine.
    \item \textbf{LLMNR et NBT-NS} diffusent des questions auxquelles
          n'importe qui peut répondre.
    \item \textbf{Le relais} exige l'absence de signature ou de liaison au
          canal.
    \item \textbf{Désactiver, signer, lier, séparer, détecter}~: les parades.
    \item \textbf{Démontrer} seulement vers une cible convenue, et fermer
          aussitôt.
\end{itemize}
\end{boxresume}

\begin{boxinfo}
\textbf{Points de contrôle de l'élève.} Je sais~: comparer NTLM et
Kerberos~; expliquer l'empoisonnement LLMNR~; énoncer les conditions du
relais~; proposer les parades et une détection.
\end{boxinfo}

\bigskip

\begin{boxexo}[title={\ding{45}\quad Exercice~9.1. LLMNR \niveaubase}]
Pourquoi un poste envoie-t-il une requête LLMNR pour « srv-imprim »~?
Quelle mesure d'hygiène l'aurait évité~?
\end{boxexo}

\begin{boxexo}[title={\ding{45}\quad Exercice~9.2. Proportion \niveaubase}]
Quelle part des serveurs de Garoua n'exige pas la signature SMB~?
\end{boxexo}

\begin{boxexo}[title={\ding{45}\quad Exercice~9.3. Kerberos \niveaubase}]
Quel organisme a défini Kerberos version~5 dans la RFC~4120~? Quel est le
rôle du centre de distribution de clés~?
\end{boxexo}

\begin{boxexo}[title={\ding{45}\quad Exercice~9.4. Casser ou relayer \niveaumoyen}]
Pourquoi les réponses recueillies n'ont-elles pas été cassées~? Que
faudrait-il pour qu'elles le soient~?
\end{boxexo}

\begin{boxexo}[title={\ding{45}\quad Exercice~9.5. Les conditions \niveaumoyen}]
Laquelle des trois conditions du Tableau~9.2 est la plus simple à
supprimer~? Laquelle protège le mieux~?
\end{boxexo}

\begin{boxexo}[title={\ding{45}\quad Exercice~9.6. Le piège \niveaumoyen}]
Décrivez le poste « leurre » et la règle d'alerte du SOC.
\end{boxexo}

\begin{boxexo}[title={\ding{45}\quad Exercice~9.7. Effets de bord \niveaumoyen}]
Quels risques de dysfonctionnement faut-il vérifier avant d'exiger la
signature SMB partout~?
\end{boxexo}

\begin{boxexo}[title={\ding{45}\quad Exercice~9.8. ATT\&CK \niveaumoyen}]
Dans quelle tactique d'ATT\&CK se range l'empoisonnement de résolution~?
Pourquoi~?
\end{boxexo}

\begin{boxexo}[title={\ding{45}\quad Exercice~9.9. Plan de correction \niveauexpert}]
Rédigez le plan de correction de T-07 en trois vagues, avec les tests de
non-régression.
\end{boxexo}

\begin{boxexo}[title={\ding{45}\quad Exercice~9.10. Réduire NTLM \niveauexpert}]
Proposez une démarche pour mesurer l'usage de NTLM chez LiZbiZ et le
réduire sans interrompre les applications.
\end{boxexo}

\bigskip

\begin{boxinfo}
\textbf{Pour aller plus loin.}
\begin{itemize}
    \item RFC~4120 (Kerberos)~\cite{rfc4120}~; RFC~4795 (LLMNR)~\cite{rfc4795}.
    \item MITRE ATT\&CK~\cite{mitreattack}.
\end{itemize}
\end{boxinfo}

\begin{boxinfo}
\textbf{Transition.} Kerberos et l'annuaire forment le cœur du réseau
Windows~: le Chapitre~\ref{chap:ad} est consacré à Active Directory.
\end{boxinfo}
}{}
        \IfFileExists{partie3/P03_C10_active_directory.tex}{% ============================================================================
% Fichier  : partie3/P03_C10_active_directory.tex
% Titre    : Active Directory
% Auteur   : MINKA MI NGUIDJOÏ Thierry Emmanuel
% Version  : 1.0 -- novembre 2026
% ============================================================================
% Sources : RFC 4120 (Kerberos) ; documentation publique de Microsoft (comptes
% de service gérés, modèle d'administration par niveaux, LAPS, groupe
% Protected Users) ; MITRE ATT&CK (T1558.003 Kerberoasting, T1558.004 AS-REP
% Roasting). Opération LOGONE : cas fictif, aucun code d'attaque.
% ============================================================================

\chapter{Active Directory}
\label{chap:ad}
\chaptermark{Active Directory}

\epigraph{%
    « Dans un domaine, l'attaquant ne cherche pas la faille~;\\
    il cherche le chemin. »%
}{Principe du cours}

\niveauChapitre{Expert}{Chapitre~\ref{chap:reseau}~; cours SEC3266,
chapitre sur les identités.}

\begin{boxobjectifs}
\begin{itemize}
    \item Décrire les objets d'un annuaire Active Directory et les groupes à
          privilèges.
    \item Expliquer l'obtention de tickets Kerberos et ses détournements
          (kerberoasting, AS-REP roasting).
    \item Raisonner en chemins d'attaque dans un graphe de droits.
    \item Appliquer le modèle d'administration par niveaux.
    \item Détecter et corriger les faiblesses les plus exploitées.
\end{itemize}
\end{boxobjectifs}

\bigskip

% ============================================================================
\section*{Question directrice : comment un compte d'utilisateur ordinaire devient-il administrateur du domaine en cinq heures~?}
% ============================================================================

Le 28~octobre, Mme~NGO~MBOCK part du compte d'utilisateur ordinaire fourni
par LiZbiZ. Elle n'exploite aucune faille logicielle. Elle interroge
l'annuaire, comme tout utilisateur peut le faire, demande un ticket pour un
service, casse hors ligne le mot de passe d'un compte de service, et
découvre que ce compte est administrateur du domaine. Ce chapitre, sur deux
séances, déroule ce chemin et ses parades.

\medskip
\textbf{Plan du chapitre.}
\begin{enumerate}
    \item L'annuaire et ses privilèges (section~\ref{sec:c10-annuaire}).
    \item Kerberos en bref (section~\ref{sec:c10-kerberos}).
    \item Le kerberoasting (section~\ref{sec:c10-kerberoasting}).
    \item L'AS-REP roasting (section~\ref{sec:c10-asrep}).
    \item Les chemins d'attaque (section~\ref{sec:c10-chemins}).
    \item Le modèle par niveaux et les défenses (section~\ref{sec:c10-defense}).
\end{enumerate}

\separateur

% ============================================================================
\section{L'Annuaire et ses Privilèges}
\label{sec:c10-annuaire}
% ============================================================================

\begin{boxdef}
\termecle{Active Directory}~: service d'annuaire de Microsoft qui gère, dans
un \emph{domaine}, les comptes, les groupes, les ordinateurs et les
stratégies, et qui authentifie les utilisateurs par Kerberos. Les
\termecle{contrôleurs de domaine} en hébergent la base.
\end{boxdef}

\begin{tableaucours}{L{4.4cm}L{7.8cm}}{%
    \tableauHeader{Objet ou groupe} &
    \tableauHeader{Pourquoi il intéresse l'attaquant}
}
Admins du domaine & Contrôle total du domaine \\
Administrateurs de l'entreprise & Contrôle de toute la forêt \\
Opérateurs de sauvegarde, de serveur & Accès étendu, souvent oublié \\
Comptes de service & Mots de passe rarement changés, droits souvent
excessifs \\
Stratégies de groupe & Exécution de code sur toutes les machines liées \\
Listes de contrôle d'accès & Un droit d'écriture sur un objet peut suffire à
en prendre le contrôle \\
\end{tableaucours}
\vspace{-0.8em}
\begin{center}{\small\itshape Tableau~10.1. Ce que vise l'attaquant dans l'annuaire}\end{center}

Tout utilisateur authentifié peut, par défaut, lire l'essentiel de
l'annuaire~: comptes, groupes, attributs. C'est utile au fonctionnement, et
c'est ce qui rend la reconnaissance interne si facile.

% ============================================================================
\section{Kerberos en Bref}
\label{sec:c10-kerberos}
% ============================================================================

\begin{tableaucours}{cL{5.8cm}L{5.2cm}}{%
    \tableauHeader{Étape} &
    \tableauHeader{Échange} &
    \tableauHeader{Secret utilisé}
}
1 & L'utilisateur demande un ticket d'accès général (TGT) au centre de
distribution & Clé dérivée du mot de passe de l'utilisateur
(pré-authentification) \\
2 & Il présente le TGT et demande un ticket de service pour un nom de
service (SPN) & Clé du centre de distribution \\
3 & Le centre renvoie un ticket de service chiffré & Clé dérivée du mot de
passe du \emph{compte de service} \\
4 & L'utilisateur présente le ticket au service & Le service le déchiffre \\
\end{tableaucours}
\vspace{-0.8em}
\begin{center}{\small\itshape Tableau~10.2. Les échanges Kerberos (d'après la RFC~4120)}\end{center}

% ============================================================================
\section{Le Kerberoasting}
\label{sec:c10-kerberoasting}
% ============================================================================

\begin{boxdef}
\termecle{Kerberoasting}~: technique (ATT\&CK T1558.003) par laquelle un
utilisateur authentifié demande des tickets de service pour des comptes
associés à un SPN, puis tente de retrouver hors ligne leur mot de passe à
partir du ticket chiffré.
\end{boxdef}

Demander un ticket de service est un geste normal~: n'importe quel
utilisateur le fait. La faiblesse est ailleurs~: un compte de service dont
le mot de passe est court ou prévisible se casse hors ligne, sans laisser de
trace d'échec.

\begin{boxscenario}{Opération LOGONE. Le compte de la facturation}
\textbf{Constatations} (28~octobre).
\begin{itemize}
    \item 9~h~10~: interrogation de l'annuaire~; 6~comptes ont un SPN.
    \item 9~h~25~: tickets de service demandés pour ces 6~comptes.
    \item 11~h~40~: le mot de passe de « svc\_sql\_facturation » est retrouvé
          hors ligne~: 11~caractères, formé du nom de l'application, d'un
          symbole et d'une année.
    \item 14~h~05~: ce compte est membre des Admins du domaine (déjà relevé
          à la revue des comptes de mars, cours SEC3266).
    \item 14~h~10~: preuve capturée, test arrêté, contacts prévenus.
\end{itemize}

\medskip
\textbf{Analyse} (confiance élevée). Aucune faille logicielle~: trois
faiblesses de configuration se combinent (mot de passe faible, mot de passe
jamais changé, droits excessifs). Le SOC n'a rien vu~: aucune règle ne
surveillait les demandes de tickets.
\end{boxscenario}

% ============================================================================
\section{L'AS-REP Roasting}
\label{sec:c10-asrep}
% ============================================================================

Lorsqu'un compte a la \emph{pré-authentification Kerberos désactivée},
n'importe qui peut demander pour lui une réponse chiffrée avec une clé
dérivée de son mot de passe, et tenter de la casser hors ligne (ATT\&CK
T1558.004). Cette option, parfois activée pour d'anciennes applications,
devrait être exceptionnelle et tracée. Chez LiZbiZ, deux comptes l'avaient
encore, dont un compte de test oublié.

% ============================================================================
\section{Les Chemins d'Attaque}
\label{sec:c10-chemins}
% ============================================================================

L'annuaire forme un \termecle{graphe}~: des utilisateurs membres de groupes,
des groupes ayant des droits sur des objets, des sessions ouvertes sur des
machines. Un chemin d'attaque est une suite d'arêtes qui mène d'un compte
ordinaire à un groupe à privilèges. Des outils publics de cartographie
calculent ces chemins~; les défenseurs les utilisent pour les couper.

\begin{tableaucours}{L{1.0cm}L{5.6cm}L{5.6cm}}{%
    \tableauHeader{Arête} &
    \tableauHeader{Relation} &
    \tableauHeader{Exemple chez LiZbiZ}
}
1 & Membre de & svc\_sql\_facturation \(\rightarrow\) Admins du domaine \\
2 & Session ouverte sur & Un administrateur connecté sur PC-GAR-12 \\
3 & Administrateur local de & Le groupe « Support Garoua » sur tous les
postes de Garoua \\
4 & Droit d'écriture sur & Le groupe « Support Garoua » peut modifier le
groupe « Admins serveurs » \\
\end{tableaucours}
\vspace{-0.8em}
\begin{center}{\small\itshape Tableau~10.3. Des relations qui forment des chemins}\end{center}

\begin{boxalert}
\textbf{Violation typique.} Corriger le seul compte exploité et déclarer le
problème résolu.

\smallskip
\textbf{Règle~:} le test a montré \emph{un} chemin. La correction consiste à
cartographier tous les chemins vers les groupes à privilèges et à couper les
plus courts, puis à refaire la cartographie.
\end{boxalert}

% ============================================================================
\section{Le Modèle par Niveaux et les Défenses}
\label{sec:c10-defense}
% ============================================================================

Le \termecle{modèle d'administration par niveaux} sépare les comptes et les
machines en niveaux (contrôleurs de domaine et identités~; serveurs~;
postes)~: un compte d'un niveau ne s'ouvre jamais sur une machine d'un
niveau inférieur, pour que ses secrets n'y soient jamais exposés.

\begin{tableaucours}{L{4.6cm}L{7.6cm}}{%
    \tableauHeader{Défense} &
    \tableauHeader{Effet}
}
Comptes de service gérés & Mots de passe longs, aléatoires, changés
automatiquement~: le kerberoasting devient inutile \\
Mots de passe de 25~caractères au moins pour les autres comptes de service
& Rend le cassage hors ligne irréaliste \\
Retrait des droits excessifs & Un compte de service n'est pas
administrateur du domaine \\
Pré-authentification obligatoire & Supprime l'AS-REP roasting \\
Mot de passe administrateur local unique (LAPS) & Bloque le mouvement
latéral (cours SEC3266) \\
Groupe Protected Users, modèle par niveaux & Limite l'exposition des
secrets à privilèges \\
Détection & Événements de demande de ticket inhabituels (nombreux SPN en
peu de temps, chiffrement faible demandé) \\
\end{tableaucours}
\vspace{-0.8em}
\begin{center}{\small\itshape Tableau~10.4. Les défenses de l'annuaire}\end{center}

\begin{boxobjet}{T-01 · Compte de service à privilèges et mot de passe faible}
\textbf{Description~:} « svc\_sql\_facturation » a un SPN, un mot de passe de
11~caractères prévisible, jamais changé, et appartient aux Admins du
domaine. \textbf{Preuve~:} appartenance au groupe (capture du 28~octobre,
14~h~08), mot de passe non reproduit. \textbf{Gravité~:} critique.
\textbf{Impact~:} contrôle total du domaine par tout utilisateur
authentifié. \textbf{Recommandation~:} retirer le compte des Admins du
domaine, le convertir en compte de service géré, cartographier les autres
chemins. \textbf{Détection~:} alerte sur les demandes de tickets en rafale.
\end{boxobjet}

\begin{boxmetier}{Analyste SOC et équipe rouge}
\textbf{Placer un compte « appât ».} Un compte de service fictif, avec un
SPN et aucun usage légitime~: toute demande de ticket le concernant signale
une tentative de kerberoasting.
\end{boxmetier}

\separateur

\begin{boxresume}
\textbf{Ce qu'il faut retenir de ce chapitre~:}
\begin{itemize}
    \item \textbf{L'annuaire} est lisible par tout utilisateur~; les
          groupes à privilèges, les comptes de service et les droits sont
          les cibles.
    \item \textbf{Kerberos} chiffre le ticket de service avec le secret du
          compte de service~: d'où le kerberoasting.
    \item \textbf{AS-REP roasting}~: la pré-authentification désactivée
          expose le compte.
    \item \textbf{On raisonne en chemins}~; on cartographie et on coupe les
          plus courts.
    \item \textbf{Comptes de service gérés, moindre privilège, niveaux,
          détection}~: les défenses.
\end{itemize}
\end{boxresume}

\begin{boxinfo}
\textbf{Points de contrôle de l'élève.} Je sais~: nommer les groupes à
privilèges~; décrire les quatre échanges Kerberos~; expliquer le
kerberoasting et l'AS-REP roasting~; lire un chemin d'attaque~; proposer
les défenses et une détection.
\end{boxinfo}

\bigskip

\begin{boxexo}[title={\ding{45}\quad Exercice~10.1. Les groupes \niveaubase}]
Citez trois groupes ou types de comptes d'un domaine qui intéressent en
priorité l'attaquant, et pourquoi.
\end{boxexo}

\begin{boxexo}[title={\ding{45}\quad Exercice~10.2. Le secret \niveaubase}]
Avec quel secret le ticket de service est-il chiffré~? Pourquoi est-ce la
source du kerberoasting~?
\end{boxexo}

\begin{boxexo}[title={\ding{45}\quad Exercice~10.3. Chronologie \niveaubase}]
Combien de temps s'écoule entre l'interrogation de l'annuaire et la preuve
d'accès aux Admins du domaine~?
\end{boxexo}

\begin{boxexo}[title={\ding{45}\quad Exercice~10.4. Les trois faiblesses \niveaumoyen}]
Identifiez les trois faiblesses de configuration combinées dans T-01, et la
mesure qui supprime chacune.
\end{boxexo}

\begin{boxexo}[title={\ding{45}\quad Exercice~10.5. Pourquoi le SOC n'a rien vu \niveaumoyen}]
Pourquoi le kerberoasting est-il discret~? Quelle règle aurait pu le voir~?
\end{boxexo}

\begin{boxexo}[title={\ding{45}\quad Exercice~10.6. Le chemin \niveaumoyen}]
À partir du Tableau~10.3, décrivez un second chemin, sans le compte de la
facturation, qui mène un membre de « Support Garoua » vers des droits
élevés.
\end{boxexo}

\begin{boxexo}[title={\ding{45}\quad Exercice~10.7. Niveaux \niveaumoyen}]
Un administrateur du domaine lit ses courriels sur PC-GAR-12. Quelle règle
du modèle par niveaux viole-t-il~? Quel risque crée-t-il~?
\end{boxexo}

\begin{boxexo}[title={\ding{45}\quad Exercice~10.8. AS-REP \niveaumoyen}]
Comment recenser les comptes dont la pré-authentification est désactivée~?
Que faire de chacun~?
\end{boxexo}

\begin{boxexo}[title={\ding{45}\quad Exercice~10.9. Le plan \niveauexpert}]
Rédigez le plan de durcissement de l'annuaire de LiZbiZ sur trois mois, à
partir du Tableau~10.4.
\end{boxexo}

\begin{boxexo}[title={\ding{45}\quad Exercice~10.10. Mesurer les chemins \niveauexpert}]
Proposez un indicateur fondé sur les chemins d'attaque (nombre, longueur)
et un objectif pour le comité des risques.
\end{boxexo}

\bigskip

\begin{boxinfo}
\textbf{Pour aller plus loin.}
\begin{itemize}
    \item RFC~4120~\cite{rfc4120}~; MITRE ATT\&CK~\cite{mitreattack}.
\end{itemize}
\end{boxinfo}

\begin{boxinfo}
\textbf{Transition.} Sur une machine, l'attaquant cherche aussi à passer
d'un compte limité à un compte d'administration~: le
Chapitre~\ref{chap:privileges} traite de l'élévation de privilèges.
\end{boxinfo}
}{}
        \IfFileExists{partie3/P03_C11_privileges.tex}{% ============================================================================
% Fichier  : partie3/P03_C11_privileges.tex
% Titre    : Élévation de Privilèges
% Auteur   : MINKA MI NGUIDJOÏ Thierry Emmanuel
% Version  : 1.0 -- novembre 2026
% ============================================================================
% Sources : MITRE ATT&CK (tactique Privilege Escalation) ; CIS Benchmarks ;
% documentation de sudo et des capacités Linux. Opération LOGONE : cas fictif.
% ============================================================================

\chapter{Élévation de Privilèges}
\label{chap:privileges}
\chaptermark{Élévation de privilèges}

\epigraph{%
    « L'élévation de privilèges exploite rarement le noyau~;\\
    elle exploite le plus souvent une négligence de configuration. »%
}{Principe du cours}

\niveauChapitre{Intermédiaire à expert}{Chapitre~\ref{chap:exploitation}~;
cours SEC3266, chapitre sur les systèmes.}

\begin{boxobjectifs}
\begin{itemize}
    \item Distinguer élévation verticale et mouvement horizontal.
    \item Recenser les vecteurs d'élévation sous Linux et sous Windows.
    \item Lire un relevé de configuration pour y repérer un vecteur.
    \item Classer les vecteurs par probabilité et par risque de dommage.
    \item Proposer la correction et la détection de chaque vecteur.
\end{itemize}
\end{boxobjectifs}

\bigskip

% ============================================================================
\section*{Question directrice : un compte web limité peut-il devenir administrateur du serveur~?}
% ============================================================================

Après la démonstration sur l'ancien portail (Chapitre~\ref{chap:exploitation}),
M.~TAMO dispose d'un compte web sans privilège. Plutôt que d'exploiter le
serveur, il en lit la configuration, avec l'accord de M.~KOUNOU~: droits
sudo, programmes SUID, tâches planifiées, fichiers modifiables. Trois
vecteurs apparaissent~; un seul suffit.

\medskip
\textbf{Plan du chapitre.}
\begin{enumerate}
    \item Vertical et horizontal (section~\ref{sec:c11-types}).
    \item Les vecteurs sous Linux (section~\ref{sec:c11-linux}).
    \item Les vecteurs sous Windows (section~\ref{sec:c11-windows}).
    \item Lire un relevé (section~\ref{sec:c11-releve}).
    \item Corriger et détecter (section~\ref{sec:c11-corriger}).
\end{enumerate}

\separateur

% ============================================================================
\section{Vertical et Horizontal}
\label{sec:c11-types}
% ============================================================================

\begin{boxdef}
\termecle{Élévation de privilèges}~: obtention de droits supérieurs à ceux
du compte de départ (verticale~: utilisateur vers administrateur).
\termecle{Mouvement horizontal}~: passage d'un compte à un autre de même
niveau, souvent pour atteindre des données ou une autre machine.
\end{boxdef}

% ============================================================================
\section{Les Vecteurs sous Linux}
\label{sec:c11-linux}
% ============================================================================

\begin{tableaucours}{L{3.6cm}L{4.6cm}L{4.0cm}}{%
    \tableauHeader{Vecteur} &
    \tableauHeader{Faiblesse} &
    \tableauHeader{Correction}
}
sudo trop large & Commande autorisée sans mot de passe qui permet d'ouvrir
un interpréteur & Restreindre aux commandes nécessaires \\
Programme SUID inhabituel & S'exécute avec les droits de son propriétaire
& Retirer le bit SUID, supprimer le fichier \\
Tâche planifiée modifiable & Script exécuté par root, modifiable par
d'autres & Droits stricts sur les scripts \\
Capacités excessives & Un binaire dispose d'une capacité qui contourne les
contrôles & Retirer la capacité \\
Fichier sensible lisible & Clé, mot de passe dans un fichier accessible &
Droits stricts, coffre de secrets \\
Noyau non corrigé & Faille du noyau & Correctifs~; vecteur risqué à
exploiter \\
\end{tableaucours}
\vspace{-0.8em}
\begin{center}{\small\itshape Tableau~11.1. Vecteurs d'élévation sous Linux}\end{center}

% ============================================================================
\section{Les Vecteurs sous Windows}
\label{sec:c11-windows}
% ============================================================================

\begin{tableaucours}{L{3.6cm}L{4.6cm}L{4.0cm}}{%
    \tableauHeader{Vecteur} &
    \tableauHeader{Faiblesse} &
    \tableauHeader{Correction}
}
Service mal protégé & Binaire ou configuration du service modifiable par
un utilisateur & Droits stricts~; comptes de service dédiés \\
Chemin de service sans guillemets & Chemin avec espaces interprété par
morceaux & Mettre le chemin entre guillemets \\
Installation toujours élevée & Stratégie qui installe tout paquet avec les
droits système & Désactiver la stratégie \\
Mot de passe en clair & Script ou fichier de configuration lisible &
Supprimer, changer le mot de passe \\
Administrateur local partagé & Même mot de passe sur toutes les machines
& Mot de passe unique par machine \\
\end{tableaucours}
\vspace{-0.8em}
\begin{center}{\small\itshape Tableau~11.2. Vecteurs d'élévation sous Windows}\end{center}

\begin{boxscenario}{Opération LOGONE. Le script du partage}
\textbf{Constatation.} Sur le partage de fichiers \texttt{\textbackslash\textbackslash FS-01\textbackslash scripts},
lisible par tous les utilisateurs, un script de déploiement contient en
clair le mot de passe du compte administrateur local des postes de Douala.

\medskip
\textbf{Analyse} (confiance élevée). Depuis mars, chaque poste a son propre
mot de passe administrateur local (LAPS)~; mais ce script, antérieur, a
servi à leur installation, et 23 postes non réinstallés gardent l'ancien
mot de passe commun. C'est le constat critique T-02.
\end{boxscenario}

% ============================================================================
\section{Lire un Relevé}
\label{sec:c11-releve}
% ============================================================================

\begin{boxprocedure}
\textbf{Recenser les vecteurs à partir d'un relevé}
\begin{enumerate}
    \item Comptes et groupes~: qui a l'UID~0 (Linux), qui est administrateur
          local (Windows)~?
    \item Droits délégués~: \texttt{sudo -l}~; services et leurs droits.
    \item Programmes SUID et capacités hors des chemins du système.
    \item Tâches planifiées et scripts qu'elles exécutent~: qui peut les
          modifier~?
    \item Fichiers lisibles contenant des secrets.
    \item Version du noyau et correctifs, en dernier recours.
\end{enumerate}
\end{boxprocedure}

\begin{boxalert}
\textbf{Violation typique.} Commencer par une faille du noyau.

\smallskip
\textbf{Règle~:} une exploitation du noyau peut faire tomber le serveur. On
cherche d'abord les vecteurs de configuration, sans risque et plus
probables~; on rapporte la faille du noyau sans l'exploiter, sauf accord
explicite.
\end{boxalert}

% ============================================================================
\section{Corriger et Détecter}
\label{sec:c11-corriger}
% ============================================================================

\begin{boxobjet}{T-02 · Mot de passe administrateur local en clair}
\textbf{Description~:} un script lisible par tous sur
\texttt{\textbackslash\textbackslash FS-01\textbackslash scripts} contient le
mot de passe administrateur local commun de 23~postes de Douala.
\textbf{Preuve~:} chemin et date du fichier, mot de passe masqué.
\textbf{Gravité~:} critique. \textbf{Impact~:} administration de 23~postes
et mouvement latéral. \textbf{Recommandation~:} supprimer le script,
changer le mot de passe sur les 23~postes (LAPS), rechercher d'autres
secrets dans les partages. \textbf{Détection~:} connexions avec le compte
administrateur local depuis le réseau.
\end{boxobjet}

\begin{boxmetier}{Auditeur de sécurité}
\textbf{Chercher les secrets là où on travaille.} L'auditeur recherche les
mots de passe dans les partages, les scripts, les stratégies de groupe et
les dépôts internes~: c'est souvent là, et non dans les serveurs, que se
trouvent les clés du domaine.
\end{boxmetier}

\separateur

\begin{boxresume}
\textbf{Ce qu'il faut retenir de ce chapitre~:}
\begin{itemize}
    \item \textbf{Vertical} (plus de droits) et \textbf{horizontal} (autre
          compte).
    \item \textbf{Linux}~: sudo, SUID, tâches, capacités, fichiers, noyau.
    \item \textbf{Windows}~: services, chemins, stratégies, mots de passe en
          clair, administrateur local partagé.
    \item \textbf{On lit la configuration d'abord}~; le noyau en dernier.
    \item \textbf{Chaque vecteur} a sa correction et sa détection.
\end{itemize}
\end{boxresume}

\begin{boxinfo}
\textbf{Points de contrôle de l'élève.} Je sais~: distinguer élévation et
mouvement horizontal~; citer les vecteurs Linux et Windows~; lire un relevé
dans l'ordre~; rédiger un constat avec correction et détection.
\end{boxinfo}

\bigskip

\begin{boxexo}[title={\ding{45}\quad Exercice~11.1. Vertical ou horizontal \niveaubase}]
Classez~: passer du compte web au compte root~; utiliser le compte d'un
collègue de même service~; devenir administrateur local d'un poste.
\end{boxexo}

\begin{boxexo}[title={\ding{45}\quad Exercice~11.2. SUID \niveaubase}]
Pourquoi un fichier SUID appartenant à root dans \texttt{/tmp} est-il
suspect~? (cours SEC3266, lab de durcissement)
\end{boxexo}

\begin{boxexo}[title={\ding{45}\quad Exercice~11.3. Chemin \niveaubase}]
Expliquez le risque d'un chemin de service non entouré de guillemets.
\end{boxexo}

\begin{boxexo}[title={\ding{45}\quad Exercice~11.4. sudo \niveaumoyen}]
Un utilisateur peut lancer un éditeur de texte en sudo sans mot de passe.
Pourquoi est-ce dangereux~? Que proposez-vous~?
\end{boxexo}

\begin{boxexo}[title={\ding{45}\quad Exercice~11.5. Le script \niveaumoyen}]
Pourquoi 23~postes gardent-ils l'ancien mot de passe malgré LAPS~?
Comment les retrouver~?
\end{boxexo}

\begin{boxexo}[title={\ding{45}\quad Exercice~11.6. Ordre \niveaumoyen}]
Justifiez l'ordre de la procédure de recensement.
\end{boxexo}

\begin{boxexo}[title={\ding{45}\quad Exercice~11.7. Tâche planifiée \niveaumoyen}]
Une tâche exécutée par root lance un script modifiable par le groupe
« web ». Rédigez le constat.
\end{boxexo}

\begin{boxexo}[title={\ding{45}\quad Exercice~11.8. Détecter \niveaumoyen}]
Quels événements Windows permettent de détecter l'usage du compte
administrateur local depuis le réseau~?
\end{boxexo}

\begin{boxexo}[title={\ding{45}\quad Exercice~11.9. Recherche de secrets \niveauexpert}]
Concevez une campagne interne, autorisée, de recherche de secrets dans les
partages de LiZbiZ~: périmètre, méthode, traitement des résultats.
\end{boxexo}

\begin{boxexo}[title={\ding{45}\quad Exercice~11.10. Standard \niveauexpert}]
Rédigez cinq règles d'un standard de configuration qui ferment les vecteurs
des Tableaux~11.1 et~11.2.
\end{boxexo}

\bigskip

\begin{boxinfo}
\textbf{Pour aller plus loin.}
\begin{itemize}
    \item MITRE ATT\&CK~\cite{mitreattack}~; CIS Benchmarks~\cite{cis2021controls}.
\end{itemize}
\end{boxinfo}

\begin{boxinfo}
\textbf{Transition.} Une fois des droits obtenus, l'attaquant s'étend et
s'installe~: c'est la post-exploitation, objet du Chapitre~\ref{chap:post}.
\end{boxinfo}
}{}
        \IfFileExists{partie3/P03_C12_post_exploitation.tex}{% ============================================================================
% Fichier  : partie3/P03_C12_post_exploitation.tex
% Titre    : Post-Exploitation, Mouvement Latéral, Persistance
% Auteur   : MINKA MI NGUIDJOÏ Thierry Emmanuel
% Version  : 1.0 -- novembre 2026
% ============================================================================
% Sources : MITRE ATT&CK (tactiques Lateral Movement, Persistence, Collection,
% Exfiltration) ; journaux de sécurité Windows (événements 4624, 4672, 4698,
% 4720, 7045). Opérations GAROUA et LOGONE : cas fictifs.
% ============================================================================

\chapter{Post-Exploitation, Mouvement Latéral, Persistance}
\label{chap:post}
\chaptermark{Post-exploitation}

\epigraph{%
    « Le testeur laisse le système tel qu'il l'a trouvé~;\\
    l'attaquant y laisse une porte. »%
}{Principe du cours}

\niveauChapitre{Expert}{Chapitres~\ref{chap:ad} et~\ref{chap:privileges}.}

\begin{boxobjectifs}
\begin{itemize}
    \item Décrire les objectifs de la post-exploitation~: comprendre,
          s'étendre, durer, atteindre les données.
    \item Expliquer les techniques de mouvement latéral et leurs traces.
    \item Reconnaître les mécanismes de persistance.
    \item Fixer les limites du testeur~: démontrer sans installer.
    \item Relier chaque technique aux événements qui la détectent.
\end{itemize}
\end{boxobjectifs}

\bigskip

% ============================================================================
\section*{Question directrice : que fait l'attaquant après être devenu administrateur, et que doit faire le testeur~?}
% ============================================================================

En mars, l'attaquant devenu administrateur a créé un compte, installé un
service, copié 40~Go et détruit les sauvegardes. En octobre, Mme~NGO~MBOCK
devenue administratrice s'est arrêtée dans les cinq minutes. Ce chapitre
décrit ce qu'un attaquant fait après l'élévation, pour que le défenseur le
voie, et ce que le testeur s'interdit.

\medskip
\textbf{Plan du chapitre.}
\begin{enumerate}
    \item Les objectifs de la post-exploitation (section~\ref{sec:c12-objectifs}).
    \item Le mouvement latéral (section~\ref{sec:c12-lateral}).
    \item La persistance (section~\ref{sec:c12-persistance}).
    \item Les limites du testeur (section~\ref{sec:c12-limites}).
    \item Ce que voient les journaux (section~\ref{sec:c12-journaux}).
\end{enumerate}

\separateur

% ============================================================================
\section{Les Objectifs de la Post-Exploitation}
\label{sec:c12-objectifs}
% ============================================================================

\begin{tableaucours}{L{3.4cm}L{4.4cm}L{4.4cm}}{%
    \tableauHeader{Objectif} &
    \tableauHeader{Tactique ATT\&CK} &
    \tableauHeader{Dans l'Opération GAROUA}
}
Comprendre l'environnement & Découverte & Balayage du réseau, interrogation
de l'annuaire \\
S'étendre & Mouvement latéral & Connexions avec le mot de passe local
partagé \\
Durer & Persistance & Compte svc\_maj2, service WinUpd32 \\
Atteindre les données & Collecte, exfiltration & Copie de 40~Go \\
Agir & Impact & Suppression des sauvegardes, chiffrement \\
\end{tableaucours}
\vspace{-0.8em}
\begin{center}{\small\itshape Tableau~12.1. Les objectifs de la post-exploitation~\cite{mitreattack}}\end{center}

% ============================================================================
\section{Le Mouvement Latéral}
\label{sec:c12-lateral}
% ============================================================================

\begin{boxdef}
\termecle{Mouvement latéral}~: ensemble des techniques par lesquelles un
attaquant, présent sur une machine, prend le contrôle d'autres machines du
réseau, le plus souvent avec des identifiants valides.
\end{boxdef}

\begin{tableaucours}{L{4.0cm}L{4.4cm}L{3.8cm}}{%
    \tableauHeader{Technique} &
    \tableauHeader{Principe} &
    \tableauHeader{Trace principale}
}
Services d'administration à distance & Exécution d'une commande par
création de service distant & Événement 7045 (service installé) \\
Bureau à distance & Session graphique & Ouverture de session de type~10 \\
Instrumentation à distance & Exécution par l'interface d'administration
Windows & Processus enfant inhabituel \\
Réutilisation d'empreinte & Authentification NTLM avec le haché, sans le
mot de passe & Ouverture de session réseau avec NTLM depuis une machine
inhabituelle \\
\end{tableaucours}
\vspace{-0.8em}
\begin{center}{\small\itshape Tableau~12.2. Techniques de mouvement latéral et traces}\end{center}

% ============================================================================
\section{La Persistance}
\label{sec:c12-persistance}
% ============================================================================

\begin{tableaucours}{L{4.0cm}L{4.4cm}L{3.8cm}}{%
    \tableauHeader{Mécanisme} &
    \tableauHeader{Exemple} &
    \tableauHeader{Trace}
}
Compte créé ou modifié & Compte ajouté aux Admins du domaine & 4720, 4728 \\
Service & Service installé au démarrage & 7045 \\
Tâche planifiée & Tâche lancée chaque nuit & 4698 \\
Clé de démarrage & Programme lancé à l'ouverture de session & Surveillance
du registre \\
Accès distant tiers & Outil de prise de main installé & Inventaire
logiciel \\
\end{tableaucours}
\vspace{-0.8em}
\begin{center}{\small\itshape Tableau~12.3. Mécanismes de persistance}\end{center}

% ============================================================================
\section{Les Limites du Testeur}
\label{sec:c12-limites}
% ============================================================================

\begin{boxalert}
\textbf{Violation typique.} Installer une persistance « pour montrer
qu'on pourrait revenir ».

\smallskip
\textbf{Règle~:} le testeur décrit la persistance possible, il ne l'installe
pas, sauf exercice d'équipe rouge expressément convenu et tracé
(Chapitre~\ref{chap:rouge}). Tout compte, fichier ou service créé est listé
et supprimé avant la fin de la mission, et la suppression est vérifiée.
\end{boxalert}

\begin{boxprocedure}
\textbf{Le nettoyage de fin de mission}
\begin{enumerate}
    \item Relire le journal de mission~: tout ce qui a été créé ou modifié.
    \item Supprimer comptes, fichiers, services, tâches créés.
    \item Restaurer les configurations modifiées.
    \item Faire vérifier la suppression par le client.
    \item Consigner le nettoyage dans le rapport.
\end{enumerate}
\end{boxprocedure}

% ============================================================================
\section{Ce que Voient les Journaux}
\label{sec:c12-journaux}
% ============================================================================

\begin{boxscenario}{Opération LOGONE. Ce que le SOC a vu, et pas vu}
\textbf{Constatations.} Le SOC de nuit a reçu les journaux des contrôleurs
de domaine et des serveurs de Garoua. Il a signalé le balayage. Il n'a pas
corrélé la connexion du compte de la facturation, inhabituelle, à un poste
ordinaire, ni sa nouvelle appartenance au groupe visé.

\medskip
\textbf{Analyse} (confiance moyenne). Les événements utiles étaient
présents~; il manquait les règles. Après l'Opération GAROUA, c'est la
seconde fois que les journaux contiennent l'attaque sans que personne ne la
lise.
\end{boxscenario}

\begin{boxobjet}{T-08 · Comptes d'administration utilisés sur des postes ordinaires}
\textbf{Description~:} des comptes membres de groupes à privilèges ouvrent
des sessions sur des postes de travail. \textbf{Preuve~:} sessions
relevées sur 7~postes de Garoua (26~octobre). \textbf{Gravité~:} moyenne.
\textbf{Impact~:} exposition de secrets à privilèges en cas de
compromission d'un poste. \textbf{Recommandation~:} appliquer le modèle par
niveaux, administrer par le bastion. \textbf{Détection~:} ouverture de
session d'un compte à privilèges sur un poste du niveau inférieur.
\end{boxobjet}

\begin{boxmetier}{Analyste SOC et équipe rouge}
\textbf{Écrire une règle par technique jouée.} Après chaque test, l'analyste
reprend les techniques ATT\&CK jouées et vérifie, une par une, qu'une règle
les aurait vues~; il écrit celles qui manquent.
\end{boxmetier}

\separateur

\begin{boxresume}
\textbf{Ce qu'il faut retenir de ce chapitre~:}
\begin{itemize}
    \item \textbf{Comprendre, s'étendre, durer, atteindre, agir}~: les
          objectifs après l'élévation.
    \item \textbf{Le mouvement latéral} utilise des identifiants valides~;
          il laisse des traces d'ouverture de session et de services.
    \item \textbf{La persistance}~: comptes, services, tâches, clés, outils
          d'accès.
    \item \textbf{Le testeur décrit, n'installe pas}~; il nettoie et fait
          vérifier.
    \item \textbf{Les journaux contiennent l'attaque}~; il faut les règles
          pour la lire.
\end{itemize}
\end{boxresume}

\begin{boxinfo}
\textbf{Points de contrôle de l'élève.} Je sais~: relier les objectifs aux
tactiques ATT\&CK~; associer une technique latérale à sa trace~; citer les
mécanismes de persistance et leurs événements~; dérouler le nettoyage de fin
de mission.
\end{boxinfo}

\bigskip

\begin{boxexo}[title={\ding{45}\quad Exercice~12.1. Tactiques \niveaubase}]
Rangez dans les tactiques d'ATT\&CK les actions de l'attaquant de mars entre
le 3 et le 5~mars.
\end{boxexo}

\begin{boxexo}[title={\ding{45}\quad Exercice~12.2. Événements \niveaubase}]
Quel événement Windows signale~: la création d'un compte~; l'installation
d'un service~; la création d'une tâche planifiée~?
\end{boxexo}

\begin{boxexo}[title={\ding{45}\quad Exercice~12.3. Nettoyage \niveaubase}]
Pourquoi faire vérifier le nettoyage par le client~?
\end{boxexo}

\begin{boxexo}[title={\ding{45}\quad Exercice~12.4. Réutilisation d'empreinte \niveaumoyen}]
Pourquoi la réutilisation d'empreinte rend-elle inutile la longueur du mot
de passe~? Quelle mesure la bloque~?
\end{boxexo}

\begin{boxexo}[title={\ding{45}\quad Exercice~12.5. Une règle \niveaumoyen}]
Écrivez en français la règle qui aurait vu la connexion du compte de la
facturation sur un poste ordinaire.
\end{boxexo}

\begin{boxexo}[title={\ding{45}\quad Exercice~12.6. Persistance tracée \niveaumoyen}]
Dans un exercice d'équipe rouge, une persistance est convenue. Quelles
règles de traçabilité imposez-vous~?
\end{boxexo}

\begin{boxexo}[title={\ding{45}\quad Exercice~12.7. T-08 \niveaumoyen}]
Expliquez le lien entre T-08 et le modèle d'administration par niveaux du
Chapitre~\ref{chap:ad}.
\end{boxexo}

\begin{boxexo}[title={\ding{45}\quad Exercice~12.8. GAROUA rejoué \niveaumoyen}]
Pour chaque mécanisme de persistance de l'attaquant de mars, donnez
l'événement qui l'aurait signalé.
\end{boxexo}

\begin{boxexo}[title={\ding{45}\quad Exercice~12.9. Couverture \niveauexpert}]
Construisez la matrice « technique ATT\&CK jouée / règle de détection /
détecté ou non » de l'Opération LOGONE.
\end{boxexo}

\begin{boxexo}[title={\ding{45}\quad Exercice~12.10. Chasse \niveauexpert}]
Proposez une recherche proactive de persistance à mener sur le parc de
LiZbiZ après l'Opération LOGONE.
\end{boxexo}

\bigskip

\begin{boxinfo}
\textbf{Pour aller plus loin.}
\begin{itemize}
    \item MITRE ATT\&CK~\cite{mitreattack}.
\end{itemize}
\end{boxinfo}

\begin{boxinfo}
\textbf{Transition.} Presque tous les chemins passent par un mot de passe~:
le Chapitre~\ref{chap:mdp} traite de leurs attaques et de leurs défenses.
\end{boxinfo}
}{}
        \IfFileExists{partie3/P03_C13_mots_de_passe.tex}{% ============================================================================
% Fichier  : partie3/P03_C13_mots_de_passe.tex
% Titre    : Mots de Passe : Attaques et Défenses
% Auteur   : MINKA MI NGUIDJOÏ Thierry Emmanuel
% Version  : 1.0 -- novembre 2026
% ============================================================================
% Sources : NIST SP 800-63B ; MITRE ATT&CK (T1110 Brute Force et ses
% sous-techniques). Audit des mots de passe de l'Opération LOGONE : cas
% fictif, chiffres cohérents entre eux.
% ============================================================================

\chapter{Mots de Passe : Attaques et Défenses}
\label{chap:mdp}
\chaptermark{Mots de passe}

\epigraph{%
    « Un mot de passe que l'on retient facilement\\
    se devine facilement. »%
}{Principe du cours}

\niveauChapitre{Intermédiaire}{Chapitre~\ref{chap:ad}~; cours SEC3266,
chapitre sur les identités.}

\begin{boxobjectifs}
\begin{itemize}
    \item Distinguer les attaques en ligne et hors ligne.
    \item Expliquer la force brute, le dictionnaire, la pulvérisation et le
          bourrage d'identifiants.
    \item Comprendre pourquoi le hachage et sa lenteur comptent.
    \item Conduire un audit de mots de passe autorisé et en rendre compte
          sans exposer personne.
    \item Proposer une politique conforme au NIST SP~800-63B.
\end{itemize}
\end{boxobjectifs}

\bigskip

% ============================================================================
\section*{Question directrice : combien de mots de passe de LiZbiZ un attaquant retrouverait-il en un jour~?}
% ============================================================================

Après T-01, M.~ETEME demande un audit des mots de passe de l'annuaire, avec
un accord écrit et sous contrôle de Mme~NKENG. NYONG SÉCURITÉ obtient une
copie des empreintes et tente, pendant 24~heures, de retrouver les mots de
passe par des listes de mots courants et des règles de transformation. Le
résultat est un chiffre, et une politique à revoir.

\medskip
\textbf{Plan du chapitre.}
\begin{enumerate}
    \item En ligne ou hors ligne (section~\ref{sec:c13-ligne}).
    \item Les attaques (section~\ref{sec:c13-attaques}).
    \item Le hachage (section~\ref{sec:c13-hachage}).
    \item L'audit de mots de passe (section~\ref{sec:c13-audit}).
    \item La politique (section~\ref{sec:c13-politique}).
\end{enumerate}

\separateur

% ============================================================================
\section{En Ligne ou Hors Ligne}
\label{sec:c13-ligne}
% ============================================================================

\begin{tableaucours}{L{2.6cm}L{4.8cm}L{4.8cm}}{%
    \tableauHeader{Attaque} &
    \tableauHeader{Principe} &
    \tableauHeader{Limite}
}
En ligne & Essais contre le service (portail, messagerie) & Lente,
visible, freinée par le verrouillage et le second facteur \\
Hors ligne & Essais contre une empreinte volée (base, ticket Kerberos) &
Limitée seulement par la puissance de calcul et la lenteur du hachage \\
\end{tableaucours}
\vspace{-0.8em}
\begin{center}{\small\itshape Tableau~13.1. Deux familles d'attaques}\end{center}

% ============================================================================
\section{Les Attaques}
\label{sec:c13-attaques}
% ============================================================================

\begin{tableaucours}{L{3.4cm}L{8.8cm}}{%
    \tableauHeader{Attaque (ATT\&CK T1110)} &
    \tableauHeader{Principe}
}
Force brute & Toutes les combinaisons jusqu'à une longueur donnée \\
Dictionnaire et règles & Mots courants, noms, années, transformés
(majuscule, chiffre, symbole) \\
Pulvérisation & Un ou deux mots de passe courants essayés sur de nombreux
comptes, pour éviter le verrouillage \\
Bourrage d'identifiants & Couples identifiant et mot de passe volés
ailleurs, réutilisés \\
\end{tableaucours}
\vspace{-0.8em}
\begin{center}{\small\itshape Tableau~13.2. Les attaques sur les mots de passe~\cite{mitreattack}}\end{center}

\begin{boxalert}
\textbf{Violation typique.} Croire qu'une règle de complexité (majuscule,
chiffre, symbole) rend un mot de passe fort.

\smallskip
\textbf{Règle~:} « Lizbiz2026! » respecte la règle et figure dans tout
dictionnaire bien conçu. La longueur et l'imprévisibilité comptent~; les
règles de transformation des attaquants connaissent les habitudes humaines.
\end{boxalert}

% ============================================================================
\section{Le Hachage}
\label{sec:c13-hachage}
% ============================================================================

Un système stocke une \emph{empreinte} du mot de passe, pas le mot de passe.
La vitesse à laquelle l'attaquant peut essayer dépend de la fonction de
hachage~: une fonction rapide (comme l'empreinte NT de Windows) permet des
milliards d'essais par seconde sur du matériel courant~; une fonction lente
et salée, conçue pour les mots de passe (bcrypt, scrypt, Argon2), en permet
des milliers. Le sel empêche de casser tous les comptes à la fois.

% ============================================================================
\section{L'Audit de Mots de Passe}
\label{sec:c13-audit}
% ============================================================================

\begin{boxprocedure}
\textbf{Conduire un audit de mots de passe}
\begin{enumerate}
    \item Accord écrit de la direction~; information du délégué du personnel
          selon les règles internes.
    \item Extraction des empreintes par l'administrateur, sur un poste isolé~;
          copie chiffrée.
    \item Cassage hors ligne, durée et méthodes fixées à l'avance.
    \item Résultats agrégés~; aucun mot de passe en clair dans le rapport~;
          liste nominative remise au seul responsable, pour réinitialisation.
    \item Destruction des empreintes et des résultats à la fin.
\end{enumerate}
\end{boxprocedure}

\begin{boxscenario}{Opération LOGONE. Vingt-quatre heures de cassage}
\textbf{Constatations.} 160~comptes actifs après la revue de mars. En
24~heures~: 38~mots de passe retrouvés (23,75~\%), dont 31 dans la
première heure. 12~comptes partagent un même mot de passe avec au moins un
autre compte. 9~mots de passe retrouvés contiennent « lizbiz ». Aucun compte
d'administration n'est retrouvé, sauf celui de T-01, déjà connu.

\medskip
\textbf{Analyse} (confiance élevée). Un attaquant qui volerait les
empreintes ouvrirait près d'un compte sur quatre en une journée. Les
mots de passe retrouvés respectent tous la règle de complexité en vigueur.
\end{boxscenario}

\begin{boxobjet}{A-01 · Mots de passe prévisibles malgré la règle de complexité}
\textbf{Mission complémentaire~:} ce constat relève de l'audit des mots de
passe, rapporté à part des quatorze constats du test (Chapitre~\ref{chap:rapport}).
\textbf{Description~:} 38~mots de passe sur 160 retrouvés en 24~heures.
\textbf{Preuve~:} statistiques agrégées~; liste nominative remise à
M.~KOUNOU. \textbf{Gravité~:} moyenne (élevée en cas de vol des empreintes).
\textbf{Recommandation~:} longueur minimale de 12 à 15~caractères, liste de
mots interdits (nom de l'entreprise, années), contrôle contre les listes
divulguées, second facteur, réinitialisation des 38~comptes.
\textbf{Détection~:} échecs de connexion répartis sur de nombreux comptes
(pulvérisation).
\end{boxobjet}

% ============================================================================
\section{La Politique}
\label{sec:c13-politique}
% ============================================================================

\begin{tableaucours}{L{5.0cm}L{7.2cm}}{%
    \tableauHeader{Recommandation (NIST SP~800-63B)} &
    \tableauHeader{Application chez LiZbiZ}
}
Longueur plutôt que complexité & 12~caractères au moins, phrases de passe
acceptées \\
Comparer aux listes de mots de passe compromis & Contrôle à chaque
changement \\
Pas de changement périodique imposé sans motif & Changement en cas de
soupçon de compromission \\
Limiter les essais & Verrouillage progressif, surveillance \\
Second facteur & Obligatoire pour les accès distants et les comptes à
privilèges \\
\end{tableaucours}
\vspace{-0.8em}
\begin{center}{\small\itshape Tableau~13.3. Une politique conforme au NIST SP~800-63B~\cite{nist80063b}}\end{center}

\begin{boxmetier}{RSSI}
\textbf{Rendre compte sans désigner.} Le RSSI présente l'audit en chiffres
agrégés au comité~; les personnes concernées reçoivent une demande de
changement, accompagnée de conseils, et non un reproche.
\end{boxmetier}

\separateur

\begin{boxresume}
\textbf{Ce qu'il faut retenir de ce chapitre~:}
\begin{itemize}
    \item \textbf{En ligne}~: lent et visible~; \textbf{hors ligne}~: limité
          par le calcul et le hachage.
    \item \textbf{Force brute, dictionnaire, pulvérisation, bourrage}~: quatre
          attaques.
    \item \textbf{Un hachage lent et salé} freine le cassage.
    \item \textbf{Un audit de mots de passe} est autorisé, agrégé, détruit
          à la fin.
    \item \textbf{Longueur, listes de compromis, second facteur}~: la
          politique du NIST.
\end{itemize}
\end{boxresume}

\begin{boxinfo}
\textbf{Points de contrôle de l'élève.} Je sais~: distinguer attaques en
ligne et hors ligne~; nommer les quatre attaques~; expliquer le rôle du
hachage et du sel~; conduire un audit de mots de passe~; rédiger une
politique conforme.
\end{boxinfo}

\bigskip

\begin{boxexo}[title={\ding{45}\quad Exercice~13.1. Le taux \niveaubase}]
Calculez le taux de mots de passe retrouvés en une heure et en 24~heures.
\end{boxexo}

\begin{boxexo}[title={\ding{45}\quad Exercice~13.2. Quelle attaque~? \niveaubase}]
Nommez l'attaque~: « Ete2026! » essayé sur 160~comptes~; les identifiants
d'un site piraté essayés sur la messagerie de LiZbiZ.
\end{boxexo}

\begin{boxexo}[title={\ding{45}\quad Exercice~13.3. Le sel \niveaubase}]
À quoi sert le sel dans le hachage d'un mot de passe~?
\end{boxexo}

\begin{boxexo}[title={\ding{45}\quad Exercice~13.4. Complexité \niveaumoyen}]
Pourquoi « Lizbiz2026! » respecte-t-il la règle de complexité et se
casse-t-il vite~?
\end{boxexo}

\begin{boxexo}[title={\ding{45}\quad Exercice~13.5. Pulvérisation \niveaumoyen}]
Pourquoi la pulvérisation échappe-t-elle au verrouillage de compte~?
Comment la détecter~?
\end{boxexo}

\begin{boxexo}[title={\ding{45}\quad Exercice~13.6. Rapport \niveaumoyen}]
Rédigez la section du rapport sur l'audit de mots de passe, sans exposer
personne.
\end{boxexo}

\begin{boxexo}[title={\ding{45}\quad Exercice~13.7. Réutilisation \niveaumoyen}]
12~comptes partagent un mot de passe. Quel risque~? Que faire~?
\end{boxexo}

\begin{boxexo}[title={\ding{45}\quad Exercice~13.8. Hachage \niveaumoyen}]
Comparez l'empreinte NT et Argon2 du point de vue d'un attaquant hors
ligne.
\end{boxexo}

\begin{boxexo}[title={\ding{45}\quad Exercice~13.9. L'accord \niveauexpert}]
Rédigez l'accord écrit d'audit des mots de passe~: objet, méthodes, durée,
données, destruction, information des salariés.
\end{boxexo}

\begin{boxexo}[title={\ding{45}\quad Exercice~13.10. Politique \niveauexpert}]
Rédigez la nouvelle politique de mots de passe de LiZbiZ et le plan de sa
mise en œuvre.
\end{boxexo}

\bigskip

\begin{boxinfo}
\textbf{Pour aller plus loin.}
\begin{itemize}
    \item NIST SP~800-63B~\cite{nist80063b}~; MITRE ATT\&CK~\cite{mitreattack}.
\end{itemize}
\end{boxinfo}

\begin{boxinfo}
\textbf{Transition.} La Partie~IV passe des systèmes aux applications~: le
Chapitre~\ref{chap:web} apprend à tester une application web.
\end{boxinfo}
}{}

    \part{Applications et Nuage}
    \label{part:applications}

        \IfFileExists{partie4/P04_C14_web.tex}{% ============================================================================
% Fichier  : partie4/P04_C14_web.tex
% Titre    : Tester une Application Web
% Auteur   : MINKA MI NGUIDJOÏ Thierry Emmanuel
% Version  : 1.0 -- novembre 2026
% ============================================================================
% Sources : OWASP Web Security Testing Guide (WSTG) ; OWASP Top 10:2021 ;
% RFC 9110 (sémantique HTTP) ; RFC 6265 (cookies). Portail de suivi de
% LiZbiZ et Opération LOGONE : cas fictifs. Deux séances (S17 et S18).
% ============================================================================

\chapter{Tester une Application Web}
\label{chap:web}
\chaptermark{Tester une application web}

\epigraph{%
    « Une application web fait confiance à ce que le navigateur lui envoie~;\\
    le testeur, lui, n'utilise pas forcément un navigateur. »%
}{Principe du cours}

\niveauChapitre{Intermédiaire}{Chapitre~\ref{chap:methodes}~; cours SEC3266,
chapitre sur la sécurité applicative.}

\begin{boxobjectifs}
\begin{itemize}
    \item Décrire une application web du point de vue du testeur~: requêtes,
          réponses, état, session.
    \item Utiliser le guide de test de l'OWASP (WSTG) pour structurer une
          campagne.
    \item Cartographier une application~: pages, paramètres, rôles.
    \item Tester l'authentification et la gestion de session.
    \item Lire un historique HTTP pour y repérer une faiblesse.
    \item Classer les constats selon le Top~10 de l'OWASP.
\end{itemize}
\end{boxobjectifs}

\bigskip

% ============================================================================
\section*{Question directrice : par où commencer le test d'un portail que l'on ne connaît pas~?}
% ============================================================================

Le portail de suivi des colis de LiZbiZ permet aux clients de voir leurs
expéditions et à 40~agents de les mettre à jour. M.~TAMO reçoit deux comptes
de test, un client et un agent, et l'adresse du portail de pré-production.
Il dispose de 48~heures. Ce chapitre montre comment cartographier, puis
tester méthodiquement, plutôt que d'essayer au hasard.

\medskip
\textbf{Plan du chapitre.}
\begin{enumerate}
    \item L'application vue par le testeur (section~\ref{sec:c14-modele}).
    \item Le guide de test de l'OWASP (section~\ref{sec:c14-wstg}).
    \item Cartographier (section~\ref{sec:c14-carto}).
    \item Authentification et session (section~\ref{sec:c14-auth}).
    \item Lire un historique HTTP (section~\ref{sec:c14-historique}).
    \item Classer selon le Top~10 (section~\ref{sec:c14-top10}).
\end{enumerate}

\separateur

% ============================================================================
\section{L'Application vue par le Testeur}
\label{sec:c14-modele}
% ============================================================================

\begin{boxdef}
\termecle{Requête HTTP}~: message envoyé par le client au serveur (méthode,
chemin, en-têtes, corps). \termecle{Réponse HTTP}~: message du serveur
(code, en-têtes, corps). \termecle{Session}~: état conservé entre requêtes,
le plus souvent par un cookie (RFC~6265).
\end{boxdef}

Le testeur observe et modifie chaque requête avant son envoi. Un serveur
qui se fie à un contrôle fait côté navigateur (champ masqué, menu grisé,
vérification en JavaScript) accepte une requête forgée qui ignore ce
contrôle. La règle fondatrice du test web~: \emph{tout ce qui vient du
client est contrôlable par le client}.

% ============================================================================
\section{Le Guide de Test de l'OWASP}
\label{sec:c14-wstg}
% ============================================================================

Le \termecle{Web Security Testing Guide} (WSTG) de l'OWASP organise le test
d'une application web en catégories~\cite{owaspwstg}. On ne les parcourt pas
toutes à chaque mission~; on choisit selon le périmètre et le temps.

\begin{tableaucours}{L{4.4cm}L{7.8cm}}{%
    \tableauHeader{Catégorie WSTG} &
    \tableauHeader{Ce qu'on y vérifie}
}
Collecte d'informations & Technologies, points d'entrée, pages cachées \\
Configuration & En-têtes, méthodes autorisées, messages d'erreur \\
Gestion d'identité et authentification & Création de compte, connexion,
réinitialisation \\
Gestion de session & Cookies, expiration, déconnexion \\
Contrôle d'accès & Un rôle peut-il atteindre les données d'un autre~? \\
Validation des entrées & Injections, scripts, chemins
(Chapitre~\ref{chap:injections}) \\
Logique métier & Enchaînements non prévus, contournement d'étapes \\
\end{tableaucours}
\vspace{-0.8em}
\begin{center}{\small\itshape Tableau~14.1. Catégories du WSTG retenues pour LiZbiZ}\end{center}

% ============================================================================
\section{Cartographier}
\label{sec:c14-carto}
% ============================================================================

\begin{boxprocedure}
\textbf{Cartographier une application web}
\begin{enumerate}
    \item Parcourir l'application avec chaque rôle, en enregistrant les
          requêtes dans un relais local autorisé.
    \item Lister les pages, les points d'entrée et leurs paramètres.
    \item Noter les en-têtes, les cookies, les technologies détectées.
    \item Repérer les paramètres qui désignent un objet (numéro de colis,
          identifiant de compte).
    \item Dresser la matrice rôle par page, pour préparer les tests de
          contrôle d'accès.
\end{enumerate}
\end{boxprocedure}

\begin{boxalert}
\textbf{Violation typique.} Lancer un robot d'exploration agressif sur la
pré-production sans prévenir.

\smallskip
\textbf{Règle~:} l'exploration automatique crée des données, envoie des
courriels, déclenche des traitements. On fixe la cadence avec le client, on
évite les formulaires d'envoi, on prévient l'équipe d'exploitation.
\end{boxalert}

% ============================================================================
\section{Authentification et Session}
\label{sec:c14-auth}
% ============================================================================

\begin{tableaucours}{L{4.4cm}L{7.8cm}}{%
    \tableauHeader{Point testé} &
    \tableauHeader{Faiblesse recherchée}
}
Robustesse de la connexion & Verrouillage absent, énumération des comptes
par les messages \\
Réinitialisation de mot de passe & Jeton prévisible, lien qui n'expire
pas \\
Cookie de session & Absence des attributs de sécurité (indisponible au
script, réservé au canal chiffré) \\
Expiration & Session valide trop longtemps, pas d'invalidation à la
déconnexion \\
Second facteur & Absent sur les comptes sensibles \\
\end{tableaucours}
\vspace{-0.8em}
\begin{center}{\small\itshape Tableau~14.2. Tests d'authentification et de session}\end{center}

\begin{boxscenario}{Opération LOGONE. Le portail qui ne ferme pas la session}
\textbf{Constatations.} Sur le portail de suivi, le cookie de session n'a
pas l'attribut qui le réserve au canal chiffré~; à la déconnexion, le jeton
reste valide 8~heures. Un message de connexion distingue « compte inconnu »
de « mot de passe erroné », ce qui permet d'énumérer les comptes.

\medskip
\textbf{Analyse} (confiance élevée). Aucune de ces faiblesses n'est critique
seule~; ensemble, elles facilitent le vol et la réutilisation de session.
Elles alimentent le constat T-14 (délai de session trop long) du rapport.
\end{boxscenario}

% ============================================================================
\section{Lire un Historique HTTP}
\label{sec:c14-historique}
% ============================================================================

L'essentiel du test web se lit dans l'historique des requêtes et des
réponses. Le lab HE2-LAB-08 fournit un historique exporté~; l'élève y
cherche les points d'entrée, les paramètres d'objet, les en-têtes
manquants et les codes de réponse révélateurs.

\begin{tableaucours}{L{3.0cm}L{9.2cm}}{%
    \tableauHeader{Indice} &
    \tableauHeader{Ce qu'il peut révéler}
}
Paramètre numérique d'objet & Référence directe, à tester pour l'accès
indirect (Chapitre~\ref{chap:injections}) \\
Code 500 & Erreur serveur, message parfois bavard \\
En-tête de version & Technologie et version, utile à l'attaquant \\
Absence d'en-têtes de sécurité & Protection du navigateur non activée \\
Réponse différente selon le compte & Énumération possible \\
\end{tableaucours}
\vspace{-0.8em}
\begin{center}{\small\itshape Tableau~14.3. Lire un historique HTTP}\end{center}

% ============================================================================
\section{Classer selon le Top 10}
\label{sec:c14-top10}
% ============================================================================

Le \termecle{Top~10} de l'OWASP classe les grandes familles de risques des
applications web~\cite{owasptop10}. Il sert à ranger les constats et à
parler un langage commun avec les équipes de développement.

\begin{boxobjet}{T-12 · En-têtes de sécurité HTTP absents}
\textbf{Description~:} le portail de suivi ne renvoie pas les en-têtes de
sécurité recommandés (politique de contenu, type de contenu strict,
transport strict). \textbf{Preuve~:} réponses relevées dans l'historique.
\textbf{Gravité~:} faible. \textbf{Impact~:} protections du navigateur non
activées~; aggrave d'autres failles. \textbf{Recommandation~:} ajouter les
en-têtes au niveau du serveur ou de l'application. \textbf{Détection~:}
contrôle périodique des en-têtes de réponse.
\end{boxobjet}

\begin{boxmetier}{Pentesteur et développeur}
\textbf{Rapprocher le constat du code.} Le pentesteur décrit la requête qui
démontre la faille~; le développeur retrouve, à partir d'elle, la ligne de
code responsable. Un bon constat web contient toujours la requête exacte.
\end{boxmetier}

\separateur

\begin{boxresume}
\textbf{Ce qu'il faut retenir de ce chapitre~:}
\begin{itemize}
    \item \textbf{Tout ce qui vient du client} est contrôlable par le
          client.
    \item \textbf{Le WSTG} structure la campagne en catégories.
    \item \textbf{Cartographier d'abord}~: pages, paramètres, rôles.
    \item \textbf{Authentification et session}~: verrouillage, cookies,
          expiration.
    \item \textbf{L'historique HTTP} contient l'essentiel des indices.
    \item \textbf{Le Top~10} range les constats et parle aux développeurs.
\end{itemize}
\end{boxresume}

\begin{boxinfo}
\textbf{Points de contrôle de l'élève.} Je sais~: décrire requête, réponse
et session~; choisir des catégories du WSTG~; cartographier une
application~; tester authentification et session~; lire un historique
HTTP~; classer un constat dans le Top~10.
\end{boxinfo}

\bigskip

\begin{boxexo}[title={\ding{45}\quad Exercice~14.1. Confiance au client \niveaubase}]
Pourquoi un contrôle fait uniquement en JavaScript ne protège-t-il pas le
serveur~?
\end{boxexo}

\begin{boxexo}[title={\ding{45}\quad Exercice~14.2. Catégories \niveaubase}]
Citez trois catégories du WSTG et un test pour chacune.
\end{boxexo}

\begin{boxexo}[title={\ding{45}\quad Exercice~14.3. Cookie \niveaubase}]
À quoi servent les attributs « réservé au canal chiffré » et « indisponible
au script » d'un cookie de session~?
\end{boxexo}

\begin{boxexo}[title={\ding{45}\quad Exercice~14.4. Énumération \niveaumoyen}]
Pourquoi un message de connexion qui distingue « compte inconnu » de « mot
de passe erroné » est-il une faiblesse~? Quel message proposer~?
\end{boxexo}

\begin{boxexo}[title={\ding{45}\quad Exercice~14.5. Cartographie \niveaumoyen}]
Décrivez la matrice rôle par page du portail de suivi, avec les deux comptes
de test.
\end{boxexo}

\begin{boxexo}[title={\ding{45}\quad Exercice~14.6. Exploration prudente \niveaumoyen}]
Quelles précautions prenez-vous avant de lancer une exploration automatique
sur la pré-production~?
\end{boxexo}

\begin{boxexo}[title={\ding{45}\quad Exercice~14.7. Historique \niveaumoyen}]
Dans un historique, quels éléments trahissent un paramètre d'objet à
tester~?
\end{boxexo}

\begin{boxexo}[title={\ding{45}\quad Exercice~14.8. Session \niveaumoyen}]
Rédigez le constat du cookie de session sans attribut de canal chiffré et
sans invalidation à la déconnexion.
\end{boxexo}

\begin{boxexo}[title={\ding{45}\quad Exercice~14.9. Plan WSTG \niveauexpert}]
Établissez le plan de test, en catégories du WSTG, pour les 48~heures
accordées à M.~TAMO.
\end{boxexo}

\begin{boxexo}[title={\ding{45}\quad Exercice~14.10. Top 10 \niveauexpert}]
Rattachez chacun des constats web de l'Opération LOGONE à une catégorie du
Top~10, en justifiant.
\end{boxexo}

\bigskip

\begin{boxinfo}
\textbf{Pour aller plus loin.}
\begin{itemize}
    \item OWASP WSTG~\cite{owaspwstg}~; OWASP Top~10~\cite{owasptop10}.
\end{itemize}
\end{boxinfo}

\begin{boxinfo}
\textbf{Transition.} La catégorie la plus riche du test web est celle des
entrées mal validées~: le Chapitre~\ref{chap:injections} traite les
injections et les défauts de contrôle d'accès.
\end{boxinfo}
}{}
        \IfFileExists{partie4/P04_C15_injections.tex}{% ============================================================================
% Fichier  : partie4/P04_C15_injections.tex
% Titre    : Injections et Contrôle d'Accès
% Auteur   : MINKA MI NGUIDJOÏ Thierry Emmanuel
% Version  : 1.0 -- novembre 2026
% ============================================================================
% Sources : OWASP Top 10:2021 (A01 Broken Access Control, A03 Injection) ;
% OWASP WSTG. Portail de suivi de LiZbiZ (constat T-05) : cas fictif.
% ============================================================================

\chapter{Injections et Contrôle d'Accès}
\label{chap:injections}
\chaptermark{Injections et contrôle d'accès}

\epigraph{%
    « L'injection confond la donnée et la commande~;\\
    le défaut d'accès confond l'identité et le droit. »%
}{Principe du cours}

\niveauChapitre{Intermédiaire à expert}{Chapitre~\ref{chap:web}~; cours
SEC3266, chapitre sur la sécurité applicative.}

\begin{boxobjectifs}
\begin{itemize}
    \item Expliquer le mécanisme commun à toutes les injections.
    \item Reconnaître l'injection de requête de base de données et l'exécution
          de script dans le navigateur.
    \item Comprendre la référence directe non contrôlée à un objet.
    \item Distinguer contrôle d'accès horizontal et vertical.
    \item Tester sans endommager ni divulguer de données réelles.
    \item Rédiger les constats et leurs corrections.
\end{itemize}
\end{boxobjectifs}

\bigskip

% ============================================================================
\section*{Question directrice : pourquoi changer un seul chiffre dans une adresse donne-t-il accès au colis d'un autre client~?}
% ============================================================================

Sur le portail de suivi, M.~TAMO, connecté avec le compte client de test,
remarque que l'adresse d'une page contient le numéro de son colis. Il
remplace ce numéro par un autre~: le portail affiche le colis d'un autre
client, nom et adresse compris. Aucune injection, aucun outil~: un simple
défaut de contrôle d'accès. C'est le constat T-05.

\medskip
\textbf{Plan du chapitre.}
\begin{enumerate}
    \item Le mécanisme de l'injection (section~\ref{sec:c15-mecanisme}).
    \item Injection de requête et script dans le navigateur
          (section~\ref{sec:c15-familles}).
    \item La référence directe à un objet (section~\ref{sec:c15-idor}).
    \item Horizontal et vertical (section~\ref{sec:c15-acces}).
    \item Tester sans nuire (section~\ref{sec:c15-prudence}).
\end{enumerate}

\separateur

% ============================================================================
\section{Le Mécanisme de l'Injection}
\label{sec:c15-mecanisme}
% ============================================================================

\begin{boxdef}
\termecle{Injection}~: faille par laquelle une donnée fournie par
l'utilisateur est interprétée comme une commande par un système en aval
(base de données, interpréteur, navigateur), faute de séparation entre la
donnée et la commande.
\end{boxdef}

La parade est toujours la même~: séparer la donnée de la commande
(requêtes préparées, échappement selon le contexte, validation stricte des
entrées). Le testeur, lui, cherche les endroits où cette séparation
manque.

% ============================================================================
\section{Injection de Requête et Script dans le Navigateur}
\label{sec:c15-familles}
% ============================================================================

\begin{tableaucours}{L{3.6cm}L{4.4cm}L{4.2cm}}{%
    \tableauHeader{Famille} &
    \tableauHeader{Principe} &
    \tableauHeader{Parade}
}
Injection de requête de base de données & Une entrée modifie la requête
exécutée & Requêtes préparées, paramètres liés \\
Script dans le navigateur & Une entrée renvoyée par la page s'exécute chez
un autre visiteur & Échappement à la sortie, politique de contenu \\
Injection de commande système & Une entrée atteint un interpréteur du
serveur & Éviter l'appel système, listes blanches \\
\end{tableaucours}
\vspace{-0.8em}
\begin{center}{\small\itshape Tableau~15.1. Trois familles d'injection}\end{center}

\begin{boxalert}
\textbf{Violation typique.} Prouver une injection de base de données par
une commande destructrice.

\smallskip
\textbf{Règle~:} on démontre la faille par une preuve inoffensive (délai,
valeur calculée, erreur contrôlée), jamais par une modification ou une
suppression de données. On ne lit pas non plus les données réelles au-delà
de ce qui prouve la faille.
\end{boxalert}

% ============================================================================
\section{La Référence Directe à un Objet}
\label{sec:c15-idor}
% ============================================================================

\begin{boxdef}
\termecle{Référence directe non contrôlée à un objet}~: l'application
désigne un objet (colis, compte, facture) par un identifiant visible, et ne
vérifie pas que l'utilisateur a le droit d'y accéder. Changer l'identifiant
donne accès à l'objet d'autrui.
\end{boxdef}

\begin{boxscenario}{Opération LOGONE. Le colis du voisin}
\textbf{Constatation.} Le portail affiche un colis à partir de son numéro,
passé dans l'adresse, sans vérifier qu'il appartient au client connecté. En
incrémentant le numéro, on accède aux colis d'autres clients~: nom,
téléphone, adresse de livraison.

\medskip
\textbf{Analyse} (confiance élevée). La preuve est faite sur trois colis du
compte de test et un colis voisin, puis arrêtée. L'impact théorique porte
sur l'ensemble des expéditions. C'est le constat T-05, de gravité élevée.
\end{boxscenario}

\begin{boxobjet}{T-05 · Accès indirect aux colis d'autrui}
\textbf{Description~:} le portail ne vérifie pas la propriété du colis
désigné par son numéro. \textbf{Preuve~:} requêtes sur quatre numéros,
captures masquées. \textbf{Gravité~:} élevée. \textbf{Impact~:} divulgation
des données de livraison de tous les clients. \textbf{Recommandation~:}
vérifier, à chaque accès, que l'objet appartient à l'utilisateur
(contrôle côté serveur)~; identifiants non devinables.
\textbf{Détection~:} un compte qui consulte de nombreux colis qui ne sont
pas les siens.
\end{boxobjet}

% ============================================================================
\section{Horizontal et Vertical}
\label{sec:c15-acces}
% ============================================================================

\begin{tableaucours}{L{3.6cm}L{8.8cm}}{%
    \tableauHeader{Défaut d'accès} &
    \tableauHeader{Exemple au portail}
}
Horizontal & Un client accède aux données d'un autre client (T-05) \\
Vertical & Un client accède à une fonction réservée aux agents (mise à jour
d'un colis) \\
\end{tableaucours}
\vspace{-0.8em}
\begin{center}{\small\itshape Tableau~15.2. Deux défauts de contrôle d'accès}\end{center}

% ============================================================================
\section{Tester sans Nuire}
\label{sec:c15-prudence}
% ============================================================================

\begin{boxprocedure}
\textbf{Tester une référence directe à un objet, proprement}
\begin{enumerate}
    \item Repérer l'identifiant d'objet dans l'adresse ou le corps.
    \item Avec un second compte de test, vérifier d'abord sur ses propres
          objets.
    \item Tenter l'accès à un objet voisin~; s'arrêter à la première preuve.
    \item Ne pas parcourir l'ensemble~: la preuve suffit.
    \item Masquer les données réelles dans le rapport.
\end{enumerate}
\end{boxprocedure}

\begin{boxmetier}{Pentesteur et auditeur de sécurité}
\textbf{La preuve minimale.} Une faille d'accès se prouve avec deux comptes
et un objet voisin. L'auditeur qui télécharge toute la base pour « mesurer
l'impact » dépasse son mandat et crée un risque~; l'impact se décrit, il ne
se collecte pas.
\end{boxmetier}

\separateur

\begin{boxresume}
\textbf{Ce qu'il faut retenir de ce chapitre~:}
\begin{itemize}
    \item \textbf{L'injection} confond donnée et commande~; la parade est de
          les séparer.
    \item \textbf{Requête, script, commande}~: trois familles.
    \item \textbf{La référence directe à un objet} donne accès à l'objet
          d'autrui.
    \item \textbf{Horizontal} (autre utilisateur) et \textbf{vertical}
          (fonction réservée).
    \item \textbf{La preuve minimale}~: deux comptes, un objet voisin, puis
          on s'arrête.
\end{itemize}
\end{boxresume}

\begin{boxinfo}
\textbf{Points de contrôle de l'élève.} Je sais~: expliquer le mécanisme
d'injection~; nommer les trois familles~; définir la référence directe à un
objet~; distinguer accès horizontal et vertical~; tester sans nuire.
\end{boxinfo}

\bigskip

\begin{boxexo}[title={\ding{45}\quad Exercice~15.1. Mécanisme \niveaubase}]
En une phrase, qu'ont en commun toutes les injections~?
\end{boxexo}

\begin{boxexo}[title={\ding{45}\quad Exercice~15.2. Familles \niveaubase}]
Associez à chaque famille sa parade principale.
\end{boxexo}

\begin{boxexo}[title={\ding{45}\quad Exercice~15.3. T-05 \niveaubase}]
Pourquoi T-05 n'est-il pas une injection~?
\end{boxexo}

\begin{boxexo}[title={\ding{45}\quad Exercice~15.4. Preuve inoffensive \niveaumoyen}]
Comment prouver une injection de base de données sans modifier ni lire de
données réelles~?
\end{boxexo}

\begin{boxexo}[title={\ding{45}\quad Exercice~15.5. Horizontal ou vertical \niveaumoyen}]
Classez~: un client lit le colis d'un autre~; un client met à jour un colis
alors que seuls les agents le peuvent.
\end{boxexo}

\begin{boxexo}[title={\ding{45}\quad Exercice~15.6. Correction \niveaumoyen}]
Rédigez la correction de T-05 pour l'équipe de développement.
\end{boxexo}

\begin{boxexo}[title={\ding{45}\quad Exercice~15.7. Détection \niveaumoyen}]
Quelle règle de détection signalerait l'exploitation de T-05 en
production~?
\end{boxexo}

\begin{boxexo}[title={\ding{45}\quad Exercice~15.8. Script \niveaumoyen}]
Pourquoi un script renvoyé par la page met-il en danger un autre visiteur,
et non le serveur~?
\end{boxexo}

\begin{boxexo}[title={\ding{45}\quad Exercice~15.9. Procédure \niveauexpert}]
Adaptez la procédure de test de référence directe au cas des factures d'un
portail comptable.
\end{boxexo}

\begin{boxexo}[title={\ding{45}\quad Exercice~15.10. Mandat \niveauexpert}]
Un client demande « la preuve que toutes les données sont exposées ».
Rédigez une réponse qui respecte le principe de la preuve minimale.
\end{boxexo}

\bigskip

\begin{boxinfo}
\textbf{Pour aller plus loin.}
\begin{itemize}
    \item OWASP Top~10~\cite{owasptop10}~; OWASP WSTG~\cite{owaspwstg}.
\end{itemize}
\end{boxinfo}

\begin{boxinfo}
\textbf{Transition.} Les applications modernes exposent des interfaces de
programmation et des applications mobiles~: le Chapitre~\ref{chap:api} leur
est consacré.
\end{boxinfo}
}{}
        \IfFileExists{partie4/P04_C16_api_mobile.tex}{% ============================================================================
% Fichier  : partie4/P04_C16_api_mobile.tex
% Titre    : Interfaces de Programmation et Applications Mobiles
% Auteur   : MINKA MI NGUIDJOÏ Thierry Emmanuel
% Version  : 1.0 -- novembre 2026
% ============================================================================
% Sources : OWASP API Security Top 10:2023 ; OWASP Mobile Application
% Security Verification Standard (MASVS). Opération MEFOU (juillet 2026) :
% cas fictif.
% ============================================================================

\chapter{Interfaces de Programmation et Applications Mobiles}
\label{chap:api}
\chaptermark{Interfaces de programmation et mobiles}

\epigraph{%
    « L'application mobile n'est qu'une façade~;\\
    le vrai test porte sur l'interface qu'elle appelle. »%
}{Principe du cours}

\niveauChapitre{Intermédiaire à expert}{Chapitre~\ref{chap:injections}.}

\begin{boxobjectifs}
\begin{itemize}
    \item Décrire une interface de programmation (API) du point de vue du
          testeur.
    \item Reconnaître les risques du Top~10 des API de l'OWASP, à commencer
          par le défaut d'accès au niveau de l'objet.
    \item Comprendre pourquoi l'application mobile ne protège pas le serveur.
    \item Tester une application mobile selon le standard MASVS.
    \item Rédiger les constats d'une mission mobile.
\end{itemize}
\end{boxobjectifs}

\bigskip

% ============================================================================
\section*{Question directrice : l'application mobile de LiZbiZ est-elle plus sûre que son portail web~?}
% ============================================================================

En juillet, avant le test principal, LiZbiZ confie à NYONG SÉCURITÉ un test
de son application mobile de suivi, l'Opération MEFOU. L'équipe découvre
vite que l'application n'est qu'une fenêtre~: toute la logique est dans
l'interface de programmation qu'elle appelle, et c'est là que se trouvent
les faiblesses. Six constats en ressortent.

\medskip
\textbf{Plan du chapitre.}
\begin{enumerate}
    \item L'interface de programmation vue par le testeur
          (section~\ref{sec:c16-api}).
    \item Le Top~10 des API (section~\ref{sec:c16-top10}).
    \item La façade mobile (section~\ref{sec:c16-facade}).
    \item Tester le mobile avec le MASVS (section~\ref{sec:c16-masvs}).
    \item L'Opération MEFOU (section~\ref{sec:c16-mefou}).
\end{enumerate}

\separateur

% ============================================================================
\section{L'Interface de Programmation vue par le Testeur}
\label{sec:c16-api}
% ============================================================================

\begin{boxdef}
\termecle{Interface de programmation} (API)~: ensemble de points d'entrée
par lesquels des programmes échangent des données, le plus souvent en HTTP,
avec des messages structurés. Le testeur l'interroge directement, sans
passer par l'application qui s'en sert.
\end{boxdef}

Une interface de programmation reçoit souvent moins d'attention que le site
web~: pas d'écran, pas de visiteur humain, l'illusion qu'elle est « interne ».
C'est une erreur~: elle est aussi exposée, et porte la logique métier.

% ============================================================================
\section{Le Top 10 des API}
\label{sec:c16-top10}
% ============================================================================

\begin{tableaucours}{L{4.6cm}L{7.6cm}}{%
    \tableauHeader{Risque (OWASP API 2023)} &
    \tableauHeader{Question du testeur}
}
Accès au niveau de l'objet & Puis-je lire ou modifier l'objet d'un autre en
changeant son identifiant~? \\
Authentification défaillante & Les jetons sont-ils solides, expirent-ils~? \\
Accès au niveau d'une propriété & Puis-je lire ou écrire un champ qui ne
m'est pas destiné~? \\
Consommation de ressources & Puis-je épuiser le service par des appels
massifs~? \\
Accès au niveau de la fonction & Puis-je appeler une fonction réservée à
un administrateur~? \\
\end{tableaucours}
\vspace{-0.8em}
\begin{center}{\small\itshape Tableau~16.1. Principaux risques des API~\cite{owaspapitop10}}\end{center}

Le premier risque, le défaut d'accès au niveau de l'objet, est l'équivalent
de la référence directe du Chapitre~\ref{chap:injections}, appliqué aux
interfaces de programmation~: c'est le plus fréquent et le plus grave.

% ============================================================================
\section{La Façade Mobile}
\label{sec:c16-facade}
% ============================================================================

\begin{boxalert}
\textbf{Violation typique.} Croire qu'un contrôle présent dans
l'application mobile protège le service.

\smallskip
\textbf{Règle~:} tout ce que fait l'application mobile, un programme peut le
refaire en appelant directement l'interface, sans ses contrôles. Une clé
cachée dans l'application n'est pas secrète~: elle se lit dans le paquet
installé.
\end{boxalert}

% ============================================================================
\section{Tester le Mobile avec le MASVS}
\label{sec:c16-masvs}
% ============================================================================

Le \termecle{standard de vérification de la sécurité des applications
mobiles} (MASVS) de l'OWASP fixe les exigences à vérifier~\cite{owaspmasvs}.

\begin{tableaucours}{L{4.2cm}L{8.0cm}}{%
    \tableauHeader{Domaine MASVS} &
    \tableauHeader{Ce qu'on vérifie}
}
Stockage des données & Secrets, jetons, données personnelles stockés en
clair sur l'appareil \\
Communication réseau & Chiffrement du canal, contrôle du certificat \\
Authentification & Jetons, biométrie, verrouillage \\
Code et résistance & Secrets codés en dur, journaux bavards \\
\end{tableaucours}
\vspace{-0.8em}
\begin{center}{\small\itshape Tableau~16.2. Domaines du MASVS}\end{center}

% ============================================================================
\section{L'Opération MEFOU}
\label{sec:c16-mefou}
% ============================================================================

\begin{boxscenario}{Opération MEFOU. Six constats sur le mobile}
\textbf{Constatations} (juillet 2026). (1)~L'application stocke le jeton de
session en clair sur l'appareil. (2)~Le service ne vérifie pas la propriété
de l'objet demandé~: un colis d'autrui est accessible par son numéro.
(3)~Une fonction d'administration est appelable par un compte ordinaire.
(4)~Le canal accepte un certificat non vérifié. (5)~Une clé de service est
codée en dur dans l'application. (6)~Les journaux de l'application contiennent
des données personnelles.

\medskip
\textbf{Analyse} (confiance élevée). Les constats (2) et (3), des défauts
d'accès, sont les plus graves~; ils préfigurent le constat T-05 du test
principal. L'application a été corrigée avant octobre, sauf le service, dont
le défaut d'accès subsiste et devient T-05.
\end{boxscenario}

\begin{boxmetier}{Pentesteur et RSSI}
\textbf{Tester le service, pas l'écran.} Pour une application mobile, le
pentesteur intercepte le trafic et rejoue les appels sans l'application~; le
RSSI retient que corriger l'application ne corrige pas le service, et que le
service doit refaire tous les contrôles.
\end{boxmetier}

\separateur

\begin{boxresume}
\textbf{Ce qu'il faut retenir de ce chapitre~:}
\begin{itemize}
    \item \textbf{L'interface de programmation} porte la logique~; on
          l'interroge directement.
    \item \textbf{Le défaut d'accès au niveau de l'objet} est le premier
          risque des API.
    \item \textbf{L'application mobile est une façade}~; ses contrôles ne
          protègent pas le service.
    \item \textbf{Le MASVS} structure le test mobile.
    \item \textbf{MEFOU} a montré, dès juillet, le défaut d'accès devenu
          T-05.
\end{itemize}
\end{boxresume}

\begin{boxinfo}
\textbf{Points de contrôle de l'élève.} Je sais~: décrire une API pour le
test~; citer les risques du Top~10 des API~; expliquer pourquoi le mobile ne
protège pas le service~; nommer les domaines du MASVS~; relier MEFOU à T-05.
\end{boxinfo}

\bigskip

\begin{boxexo}[title={\ding{45}\quad Exercice~16.1. API \niveaubase}]
Pourquoi teste-t-on une interface de programmation directement, sans
l'application qui l'utilise~?
\end{boxexo}

\begin{boxexo}[title={\ding{45}\quad Exercice~16.2. Premier risque \niveaubase}]
Quel est le premier risque du Top~10 des API~? À quoi correspond-il dans le
chapitre précédent~?
\end{boxexo}

\begin{boxexo}[title={\ding{45}\quad Exercice~16.3. Clé cachée \niveaubase}]
Pourquoi une clé « cachée » dans une application mobile n'est-elle pas
secrète~?
\end{boxexo}

\begin{boxexo}[title={\ding{45}\quad Exercice~16.4. Façade \niveaumoyen}]
Un développeur affirme que l'application empêche l'accès au colis d'autrui.
Que lui répondez-vous~?
\end{boxexo}

\begin{boxexo}[title={\ding{45}\quad Exercice~16.5. MASVS \niveaumoyen}]
Pour chacun des six constats de MEFOU, donnez le domaine MASVS concerné.
\end{boxexo}

\begin{boxexo}[title={\ding{45}\quad Exercice~16.6. MEFOU vers T-05 \niveaumoyen}]
Expliquez pourquoi le défaut d'accès de MEFOU se retrouve en octobre comme
T-05.
\end{boxexo}

\begin{boxexo}[title={\ding{45}\quad Exercice~16.7. Interception \niveaumoyen}]
Décrivez, sans outil précis, comment on observe les appels d'une application
mobile lors d'un test autorisé.
\end{boxexo}

\begin{boxexo}[title={\ding{45}\quad Exercice~16.8. Correction \niveaumoyen}]
Rédigez la correction du constat « fonction d'administration appelable par
un compte ordinaire ».
\end{boxexo}

\begin{boxexo}[title={\ding{45}\quad Exercice~16.9. Plan API \niveauexpert}]
Établissez le plan de test de l'interface de programmation de LiZbiZ selon
le Top~10 des API.
\end{boxexo}

\begin{boxexo}[title={\ding{45}\quad Exercice~16.10. Note au comité \niveauexpert}]
Rédigez une note qui explique au comité pourquoi « l'application a été
corrigée » ne suffit pas.
\end{boxexo}

\bigskip

\begin{boxinfo}
\textbf{Pour aller plus loin.}
\begin{itemize}
    \item OWASP API Security Top~10~\cite{owaspapitop10}~; OWASP
          MASVS~\cite{owaspmasvs}.
\end{itemize}
\end{boxinfo}

\begin{boxinfo}
\textbf{Transition.} Les services de LiZbiZ migrent vers le nuage~: le
Chapitre~\ref{chap:nuage} traite du test des environnements en nuage et des
conteneurs.
\end{boxinfo}
}{}
        \IfFileExists{partie4/P04_C17_nuage.tex}{% ============================================================================
% Fichier  : partie4/P04_C17_nuage.tex
% Titre    : Tester le Nuage et les Conteneurs
% Auteur   : MINKA MI NGUIDJOÏ Thierry Emmanuel
% Version  : 1.0 -- novembre 2026
% ============================================================================
% Sources : CIS Benchmarks (fournisseurs de nuage, Docker, Kubernetes) ;
% modèle de responsabilité partagée ; MITRE ATT&CK (matrices nuage et
% conteneurs). Services en nuage de LiZbiZ : cas fictif.
% ============================================================================

\chapter{Tester le Nuage et les Conteneurs}
\label{chap:nuage}
\chaptermark{Tester le nuage et les conteneurs}

\epigraph{%
    « Dans le nuage, la faille n'est presque jamais celle du fournisseur~;\\
    c'est celle de la configuration du client. »%
}{Principe du cours}

\niveauChapitre{Expert}{Chapitre~\ref{chap:api}~; notions
d'administration système.}

\begin{boxobjectifs}
\begin{itemize}
    \item Expliquer le modèle de responsabilité partagée.
    \item Reconnaître les erreurs de configuration les plus courantes dans
          le nuage.
    \item Comprendre le risque des secrets et des droits trop larges.
    \item Décrire les faiblesses des conteneurs et de leur orchestration.
    \item Tester un environnement en nuage dans les limites du contrat avec
          le fournisseur.
\end{itemize}
\end{boxobjectifs}

\bigskip

% ============================================================================
\section*{Question directrice : à qui appartient la faille quand le service est hébergé chez un fournisseur~?}
% ============================================================================

LiZbiZ héberge son nouveau portail et sa base chez un fournisseur de nuage.
M.~KOUNOU pense que la sécurité « est comprise dans l'abonnement ». Le test
montre le contraire~: le fournisseur sécurise son socle, mais la
configuration, les droits et les secrets restent la responsabilité de
LiZbiZ, et c'est là que sont les faiblesses.

\medskip
\textbf{Plan du chapitre.}
\begin{enumerate}
    \item La responsabilité partagée (section~\ref{sec:c17-partage}).
    \item Les erreurs de configuration (section~\ref{sec:c17-config}).
    \item Secrets et droits (section~\ref{sec:c17-droits}).
    \item Conteneurs et orchestration (section~\ref{sec:c17-conteneurs}).
    \item Tester dans les limites du contrat (section~\ref{sec:c17-limites}).
\end{enumerate}

\separateur

% ============================================================================
\section{La Responsabilité Partagée}
\label{sec:c17-partage}
% ============================================================================

\begin{boxdef}
\termecle{Modèle de responsabilité partagée}~: répartition des tâches de
sécurité entre le fournisseur de nuage (socle matériel, virtualisation,
services de base) et le client (configuration, droits, données,
applications). La frontière dépend du service souscrit.
\end{boxdef}

% ============================================================================
\section{Les Erreurs de Configuration}
\label{sec:c17-config}
% ============================================================================

\begin{tableaucours}{L{4.4cm}L{7.8cm}}{%
    \tableauHeader{Erreur fréquente} &
    \tableauHeader{Conséquence}
}
Espace de stockage ouvert au public & Fichiers et sauvegardes lisibles par
tous \\
Service d'administration exposé & Console ou base accessible depuis
l'Internet \\
Journalisation désactivée & Aucune trace en cas d'incident \\
Chiffrement non activé & Données au repos en clair \\
Réseau trop ouvert & Règles de filtrage permissives \\
\end{tableaucours}
\vspace{-0.8em}
\begin{center}{\small\itshape Tableau~17.1. Erreurs de configuration courantes}\end{center}

% ============================================================================
\section{Secrets et Droits}
\label{sec:c17-droits}
% ============================================================================

Deux risques dominent le nuage~: les \termecle{secrets} (clés d'accès, jetons)
laissés dans le code, les images ou les dépôts~; et les \termecle{droits}
trop larges accordés à un compte ou à une machine, qui permettent, depuis
une ressource compromise, d'atteindre tout l'environnement.

\begin{boxalert}
\textbf{Violation typique.} Utiliser une clé d'accès trouvée pour agir
au-delà du périmètre.

\smallskip
\textbf{Règle~:} une clé découverte se signale et se démontre dans la
limite stricte du périmètre convenu. On ne crée pas de ressources, on ne
touche pas aux données des autres clients du fournisseur, on respecte les
règles de test de celui-ci.
\end{boxalert}

% ============================================================================
\section{Conteneurs et Orchestration}
\label{sec:c17-conteneurs}
% ============================================================================

\begin{tableaucours}{L{4.4cm}L{7.8cm}}{%
    \tableauHeader{Faiblesse} &
    \tableauHeader{Risque}
}
Image non à jour & Vulnérabilités connues embarquées \\
Conteneur privilégié & Sortie du conteneur vers l'hôte \\
Secret dans l'image & Clé lisible par qui tire l'image \\
Tableau de bord d'orchestration exposé & Contrôle du groupe de conteneurs \\
Absence de limites & Un conteneur épuise l'hôte \\
\end{tableaucours}
\vspace{-0.8em}
\begin{center}{\small\itshape Tableau~17.2. Faiblesses des conteneurs}\end{center}

% ============================================================================
\section{Tester dans les Limites du Contrat}
\label{sec:c17-limites}
% ============================================================================

\begin{boxprocedure}
\textbf{Avant un test en nuage}
\begin{enumerate}
    \item Lire les règles de test du fournisseur~; certaines actions sont
          interdites ou exigent une déclaration.
    \item Délimiter les comptes et ressources du périmètre.
    \item Convenir des actions permises (lecture de configuration, oui~;
          création de ressources, non).
    \item Prévoir la conduite en cas de découverte d'un secret.
    \item Vérifier que l'environnement testé n'est pas partagé avec d'autres
          clients.
\end{enumerate}
\end{boxprocedure}

\begin{boxobjet}{T-10 · Service de test exposé sur l'Internet}
\textbf{Description~:} une interface de programmation de test, hébergée dans
le nuage, est accessible publiquement sans authentification.
\textbf{Preuve~:} adresse et réponse obtenue (données masquées).
\textbf{Gravité~:} moyenne. \textbf{Impact~:} divulgation de données de test,
parfois proches des données réelles~; point d'entrée. \textbf{Recommandation~:}
restreindre l'accès, authentifier, retirer les environnements de test
inutiles. \textbf{Détection~:} inventaire périodique des services exposés.
\end{boxobjet}

\begin{boxmetier}{Auditeur de sécurité et RSSI}
\textbf{Auditer la configuration, pas le fournisseur.} L'audit du nuage
porte sur ce que le client configure~: droits, secrets, exposition,
journalisation. Les référentiels de configuration sécurisée (CIS) donnent la
liste des points à vérifier.
\end{boxmetier}

\separateur

\begin{boxresume}
\textbf{Ce qu'il faut retenir de ce chapitre~:}
\begin{itemize}
    \item \textbf{La responsabilité est partagée}~; la faille est presque
          toujours côté client.
    \item \textbf{Les erreurs de configuration} (stockage ouvert, service
          exposé) dominent.
    \item \textbf{Secrets et droits trop larges} sont les deux grands
          risques.
    \item \textbf{Les conteneurs} ajoutent leurs faiblesses propres.
    \item \textbf{Le test en nuage} respecte les règles du fournisseur.
\end{itemize}
\end{boxresume}

\begin{boxinfo}
\textbf{Points de contrôle de l'élève.} Je sais~: expliquer la
responsabilité partagée~; citer les erreurs de configuration~; expliquer le
risque des secrets et des droits~; décrire les faiblesses des conteneurs~;
préparer un test en nuage conforme.
\end{boxinfo}

\bigskip

\begin{boxexo}[title={\ding{45}\quad Exercice~17.1. Frontière \niveaubase}]
Qui, du fournisseur ou du client, est responsable de la configuration des
droits d'accès~?
\end{boxexo}

\begin{boxexo}[title={\ding{45}\quad Exercice~17.2. Stockage ouvert \niveaubase}]
Pourquoi un espace de stockage ouvert au public est-il dangereux~?
\end{boxexo}

\begin{boxexo}[title={\ding{45}\quad Exercice~17.3. Secret \niveaubase}]
Citez trois endroits où un secret peut traîner dans un environnement en
nuage.
\end{boxexo}

\begin{boxexo}[title={\ding{45}\quad Exercice~17.4. Droits larges \niveaumoyen}]
Expliquez comment des droits trop larges transforment une ressource
compromise en compromission de tout l'environnement.
\end{boxexo}

\begin{boxexo}[title={\ding{45}\quad Exercice~17.5. Conteneur privilégié \niveaumoyen}]
Pourquoi un conteneur privilégié met-il en danger l'hôte~?
\end{boxexo}

\begin{boxexo}[title={\ding{45}\quad Exercice~17.6. Clé trouvée \niveaumoyen}]
Vous découvrez une clé d'accès valide. Que faites-vous, et que ne faites-vous
pas~?
\end{boxexo}

\begin{boxexo}[title={\ding{45}\quad Exercice~17.7. T-10 \niveaumoyen}]
Rédigez le plan de correction de T-10.
\end{boxexo}

\begin{boxexo}[title={\ding{45}\quad Exercice~17.8. Règles du fournisseur \niveaumoyen}]
Pourquoi faut-il lire les règles de test du fournisseur avant un test en
nuage~?
\end{boxexo}

\begin{boxexo}[title={\ding{45}\quad Exercice~17.9. Audit de configuration \niveauexpert}]
Construisez une liste de contrôle de dix points pour auditer la
configuration en nuage de LiZbiZ.
\end{boxexo}

\begin{boxexo}[title={\ding{45}\quad Exercice~17.10. Responsabilité \niveauexpert}]
Rédigez une note qui clarifie, service par service, la frontière de
responsabilité entre LiZbiZ et son fournisseur.
\end{boxexo}

\bigskip

\begin{boxinfo}
\textbf{Pour aller plus loin.}
\begin{itemize}
    \item CIS Benchmarks~\cite{cis2021controls}~; MITRE ATT\&CK~\cite{mitreattack}.
\end{itemize}
\end{boxinfo}

\begin{boxinfo}
\textbf{Transition.} La Partie~V ouvre sur les tests au-delà du système
d'information classique~: le Chapitre~\ref{chap:sansfil} traite du sans-fil
et des objets connectés.
\end{boxinfo}
}{}

    \part{Au-delà des Systèmes}
    \label{part:audela}

        \IfFileExists{partie5/P05_C18_sans_fil.tex}{% ============================================================================
% Fichier  : partie5/P05_C18_sans_fil.tex
% Titre    : Sans-Fil et Objets Connectés
% Auteur   : MINKA MI NGUIDJOÏ Thierry Emmanuel
% Version  : 1.0 -- novembre 2026
% ============================================================================
% Sources : documentation publique Wi-Fi (WPA2/WPA3) ; MITRE ATT&CK ;
% bonnes pratiques objets connectés. Entrepôt de Douala de LiZbiZ : cas
% fictif.
% ============================================================================

\chapter{Sans-Fil et Objets Connectés}
\label{chap:sansfil}
\chaptermark{Sans-fil et objets connectés}

\epigraph{%
    « Le réseau filaire s'arrête aux murs~;\\
    le réseau sans-fil franchit la rue. »%
}{Principe du cours}

\niveauChapitre{Intermédiaire}{Chapitre~\ref{chap:reseau}~; notions de
réseaux radio.}

\begin{boxobjectifs}
\begin{itemize}
    \item Décrire les faiblesses des réseaux sans-fil d'entreprise.
    \item Comprendre le point d'accès pirate et le portail jumeau.
    \item Reconnaître les risques des objets connectés.
    \item Tester un réseau sans-fil dans un cadre légal et autorisé.
    \item Rédiger les constats et les mesures de cloisonnement.
\end{itemize}
\end{boxobjectifs}

\bigskip

% ============================================================================
\section*{Question directrice : jusqu'où porte le réseau de l'entrepôt, et qui peut l'écouter~?}
% ============================================================================

L'entrepôt de Douala utilise un réseau sans-fil pour ses lecteurs de codes
et ses imprimantes d'étiquettes. Depuis le parking, Mme~NGO~MBOCK capte le
signal. Ce chapitre explique les risques du sans-fil et des objets
connectés, et comment les tester sans intercepter les communications de
tiers.

\medskip
\textbf{Plan du chapitre.}
\begin{enumerate}
    \item Les faiblesses du sans-fil (section~\ref{sec:c18-faiblesses}).
    \item Point d'accès pirate et portail jumeau (section~\ref{sec:c18-jumeau}).
    \item Les objets connectés (section~\ref{sec:c18-objets}).
    \item Tester légalement (section~\ref{sec:c18-legal}).
\end{enumerate}

\separateur

% ============================================================================
\section{Les Faiblesses du Sans-Fil}
\label{sec:c18-faiblesses}
% ============================================================================

\begin{tableaucours}{L{4.4cm}L{7.8cm}}{%
    \tableauHeader{Faiblesse} &
    \tableauHeader{Risque}
}
Clé partagée faible & Déchiffrement du trafic après capture \\
Protocole ancien & Attaques connues sur le chiffrement \\
Réseau invité relié au réseau interne & Un visiteur atteint les serveurs \\
Absence de cloisonnement & Un lecteur compromis atteint la bureautique \\
Identifiants par défaut sur les bornes & Prise de contrôle de la borne \\
\end{tableaucours}
\vspace{-0.8em}
\begin{center}{\small\itshape Tableau~18.1. Faiblesses des réseaux sans-fil}\end{center}

% ============================================================================
\section{Point d'Accès Pirate et Portail Jumeau}
\label{sec:c18-jumeau}
% ============================================================================

\begin{boxdef}
\termecle{Point d'accès pirate}~: borne non autorisée branchée sur le
réseau de l'entreprise, qui ouvre une porte dérobée.
\termecle{Portail jumeau}~: fausse borne qui imite le nom du réseau légitime
pour capter les connexions et les identifiants des usagers.
\end{boxdef}

\begin{boxalert}
\textbf{Violation typique.} Monter un portail jumeau sur le lieu de travail
et capter les connexions des salariés.

\smallskip
\textbf{Règle~:} capter les communications d'autrui est une infraction
pénale (Chapitre~\ref{chap:droit}). Un test de portail jumeau n'est permis
que sur un réseau et des comptes dédiés au test, avec accord écrit, sans
toucher aux usagers réels.
\end{boxalert}

% ============================================================================
\section{Les Objets Connectés}
\label{sec:c18-objets}
% ============================================================================

Les objets connectés (lecteurs, imprimantes, caméras, capteurs) partagent
des faiblesses récurrentes~: identifiants par défaut, micrologiciel non mis
à jour, services d'administration ouverts, absence de chiffrement. Branchés
sur le même réseau que la bureautique, ils deviennent un point d'entrée.

\begin{boxscenario}{Opération LOGONE. L'imprimante d'étiquettes}
\textbf{Constatation.} À Douala, les imprimantes d'étiquettes conservent le
mot de passe d'administration par défaut et sont sur le même réseau que les
postes. L'une d'elles expose une interface web d'administration.

\medskip
\textbf{Analyse} (confiance moyenne). Une imprimante compromise sert de
point d'appui discret~; le constat rejoint T-04 (identifiants par défaut).
La correction tient surtout au cloisonnement du réseau des objets.
\end{boxscenario}

\begin{boxobjet}{T-04 · Identifiants par défaut sur des équipements}
\textbf{Description~:} imprimantes et une borne gardent leur mot de passe
d'administration d'usine. \textbf{Preuve~:} accès aux interfaces
d'administration (sans modification). \textbf{Gravité~:} moyenne.
\textbf{Impact~:} prise de contrôle d'équipements, point d'appui sur le
réseau. \textbf{Recommandation~:} changer les mots de passe, cloisonner le
réseau des objets, inventorier. \textbf{Détection~:} connexions
d'administration depuis des adresses inhabituelles.
\end{boxobjet}

% ============================================================================
\section{Tester Légalement}
\label{sec:c18-legal}
% ============================================================================

\begin{boxprocedure}
\textbf{Encadrer un test sans-fil}
\begin{enumerate}
    \item Autorisation écrite précisant les réseaux visés et le lieu.
    \item Réseau et comptes dédiés au test pour tout essai de capture.
    \item Pas d'interception des communications de tiers.
    \item Inventaire des objets connectés avant les essais.
    \item Restitution des équipements dans leur état initial.
\end{enumerate}
\end{boxprocedure}

\begin{boxmetier}{Pentesteur et administrateur réseau}
\textbf{Cloisonner d'abord.} La meilleure réponse au risque sans-fil et
objets n'est pas l'audit répété, mais le cloisonnement~: un réseau séparé
pour les objets, un réseau invité coupé du réseau interne.
\end{boxmetier}

\separateur

\begin{boxresume}
\textbf{Ce qu'il faut retenir de ce chapitre~:}
\begin{itemize}
    \item \textbf{Le sans-fil franchit les murs}~; sa portée dépasse le
          bâtiment.
    \item \textbf{Point d'accès pirate} et \textbf{portail jumeau} sont les
          deux menaces majeures.
    \item \textbf{Les objets connectés} cumulent identifiants par défaut et
          absence de mise à jour.
    \item \textbf{Capter les communications d'autrui} est une infraction~;
          on teste sur des réseaux dédiés.
    \item \textbf{Le cloisonnement} est la première parade.
\end{itemize}
\end{boxresume}

\begin{boxinfo}
\textbf{Points de contrôle de l'élève.} Je sais~: citer les faiblesses du
sans-fil~; définir point d'accès pirate et portail jumeau~; décrire les
risques des objets connectés~; encadrer légalement un test sans-fil.
\end{boxinfo}

\bigskip

\begin{boxexo}[title={\ding{45}\quad Exercice~18.1. Portée \niveaubase}]
Pourquoi le réseau sans-fil de l'entrepôt est-il captable depuis le
parking~?
\end{boxexo}

\begin{boxexo}[title={\ding{45}\quad Exercice~18.2. Jumeau \niveaubase}]
Quelle est la différence entre un point d'accès pirate et un portail
jumeau~?
\end{boxexo}

\begin{boxexo}[title={\ding{45}\quad Exercice~18.3. Objets \niveaubase}]
Citez trois faiblesses récurrentes des objets connectés.
\end{boxexo}

\begin{boxexo}[title={\ding{45}\quad Exercice~18.4. Légalité \niveaumoyen}]
Pourquoi un test de portail jumeau sur les salariés réels est-il interdit~?
Comment le mener légalement~?
\end{boxexo}

\begin{boxexo}[title={\ding{45}\quad Exercice~18.5. T-04 \niveaumoyen}]
Rédigez la correction de T-04 pour l'entrepôt de Douala.
\end{boxexo}

\begin{boxexo}[title={\ding{45}\quad Exercice~18.6. Cloisonnement \niveaumoyen}]
Décrivez le cloisonnement réseau à mettre en place pour les objets
connectés.
\end{boxexo}

\begin{boxexo}[title={\ding{45}\quad Exercice~18.7. Réseau invité \niveaumoyen}]
Pourquoi un réseau invité relié au réseau interne est-il dangereux~?
\end{boxexo}

\begin{boxexo}[title={\ding{45}\quad Exercice~18.8. Inventaire \niveaumoyen}]
Pourquoi commence-t-on un test d'objets connectés par un inventaire~?
\end{boxexo}

\begin{boxexo}[title={\ding{45}\quad Exercice~18.9. Plan \niveauexpert}]
Rédigez le plan d'un test sans-fil autorisé pour l'entrepôt, de
l'autorisation à la restitution.
\end{boxexo}

\begin{boxexo}[title={\ding{45}\quad Exercice~18.10. Architecture \niveauexpert}]
Proposez une architecture réseau qui sépare bureautique, objets et invités.
\end{boxexo}

\bigskip

\begin{boxinfo}
\textbf{Pour aller plus loin.}
\begin{itemize}
    \item MITRE ATT\&CK~\cite{mitreattack}.
\end{itemize}
\end{boxinfo}

\begin{boxinfo}
\textbf{Transition.} La faiblesse la plus universelle n'est pas technique~:
le Chapitre~\ref{chap:social} traite de l'ingénierie sociale.
\end{boxinfo}
}{}
        \IfFileExists{partie5/P05_C19_ingenierie_sociale.tex}{% ============================================================================
% Fichier  : partie5/P05_C19_ingenierie_sociale.tex
% Titre    : Ingénierie Sociale
% Auteur   : MINKA MI NGUIDJOÏ Thierry Emmanuel
% Version  : 1.0 -- novembre 2026
% ============================================================================
% Sources : MITRE ATT&CK (T1566 Phishing) ; principes d'éthique des tests
% d'hameçonnage. Campagne autorisée de l'Opération LOGONE : cas fictif,
% chiffres cohérents.
% ============================================================================

\chapter{Ingénierie Sociale}
\label{chap:social}
\chaptermark{Ingénierie sociale}

\epigraph{%
    « On ne teste pas les personnes pour les prendre en défaut~;\\
    on teste un dispositif pour mieux les protéger. »%
}{Principe du cours}

\niveauChapitre{Intermédiaire}{Chapitre~\ref{chap:droit}~; cours SEC3266,
chapitre sur le facteur humain.}

\begin{boxobjectifs}
\begin{itemize}
    \item Définir l'ingénierie sociale et ses principaux leviers.
    \item Décrire les formes d'hameçonnage.
    \item Conduire une campagne d'hameçonnage autorisée et éthique.
    \item Mesurer et interpréter les résultats sans stigmatiser.
    \item Proposer des mesures techniques et de sensibilisation.
\end{itemize}
\end{boxobjectifs}

\bigskip

% ============================================================================
\section*{Question directrice : combien de salariés de LiZbiZ cliquent sur un lien piégé, et qu'en conclure~?}
% ============================================================================

Avec l'accord de la direction et l'information du personnel sur le principe
d'un test, NYONG SÉCURITÉ envoie un courriel d'hameçonnage simulé à
60~salariés. Le but n'est pas de piéger des personnes, mais de mesurer un
risque et d'améliorer un dispositif. Les chiffres orientent la
sensibilisation, non les sanctions.

\medskip
\textbf{Plan du chapitre.}
\begin{enumerate}
    \item Les leviers de l'ingénierie sociale (section~\ref{sec:c19-leviers}).
    \item Les formes d'hameçonnage (section~\ref{sec:c19-formes}).
    \item Une campagne éthique (section~\ref{sec:c19-campagne}).
    \item Mesurer sans stigmatiser (section~\ref{sec:c19-mesure}).
    \item Les parades (section~\ref{sec:c19-parades}).
\end{enumerate}

\separateur

% ============================================================================
\section{Les Leviers de l'Ingénierie Sociale}
\label{sec:c19-leviers}
% ============================================================================

\begin{boxdef}
\termecle{Ingénierie sociale}~: manipulation d'une personne pour obtenir une
information, un accès ou une action, en exploitant la confiance, l'autorité,
l'urgence, la peur ou la serviabilité.
\end{boxdef}

\begin{tableaucours}{L{3.0cm}L{9.2cm}}{%
    \tableauHeader{Levier} &
    \tableauHeader{Exemple}
}
Autorité & Un courriel qui se présente comme la direction \\
Urgence & « Votre compte sera fermé dans une heure » \\
Peur & « Activité suspecte détectée, cliquez pour vérifier » \\
Serviabilité & Un faux collègue demande de l'aide \\
Appât & Une pièce jointe « bulletin de paie » \\
\end{tableaucours}
\vspace{-0.8em}
\begin{center}{\small\itshape Tableau~19.1. Leviers de l'ingénierie sociale}\end{center}

% ============================================================================
\section{Les Formes d'Hameçonnage}
\label{sec:c19-formes}
% ============================================================================

\begin{tableaucours}{L{3.6cm}L{8.6cm}}{%
    \tableauHeader{Forme} &
    \tableauHeader{Description}
}
Hameçonnage de masse & Courriel identique envoyé à un grand nombre \\
Hameçonnage ciblé & Courriel personnalisé visant une personne précise \\
Hameçonnage par message court & Par message téléphonique \\
Hameçonnage vocal & Par appel téléphonique \\
\end{tableaucours}
\vspace{-0.8em}
\begin{center}{\small\itshape Tableau~19.2. Formes d'hameçonnage (ATT\&CK T1566)~\cite{mitreattack}}\end{center}

% ============================================================================
\section{Une Campagne Éthique}
\label{sec:c19-campagne}
% ============================================================================

\begin{boxprocedure}
\textbf{Conduire une campagne d'hameçonnage autorisée}
\begin{enumerate}
    \item Accord écrit de la direction~; information préalable du personnel
          sur le principe (pas sur la date).
    \item Scénario réaliste mais sans contenu humiliant ni thème sensible.
    \item Page d'atterrissage pédagogique~: la personne qui clique apprend
          aussitôt qu'il s'agissait d'un test.
    \item Collecte limitée~: taux de clic et de saisie, jamais le mot de
          passe réel.
    \item Résultats agrégés~; aucune liste nominative de « fautifs ».
\end{enumerate}
\end{boxprocedure}

\begin{boxalert}
\textbf{Violation typique.} Diffuser la liste des personnes qui ont cliqué.

\smallskip
\textbf{Règle~:} une campagne mesure un dispositif, pas des individus. On ne
collecte jamais les mots de passe réels saisis~; on ne nomme personne~; on
ne tourne pas le test en sanction, sous peine de détruire la confiance et le
signalement.
\end{boxalert}

% ============================================================================
\section{Mesurer sans Stigmatiser}
\label{sec:c19-mesure}
% ============================================================================

\begin{boxscenario}{Opération LOGONE. La campagne de 60 courriels}
\textbf{Constatations.} 60~courriels envoyés. 14~personnes ont cliqué sur le
lien (23,3~\%). 5~ont commencé à saisir des identifiants sur la page de test
(la saisie a été interrompue, aucun mot de passe réel n'a été conservé).
9~personnes ont signalé le courriel à l'informatique.

\medskip
\textbf{Analyse} (confiance élevée). Le taux de clic (23,3~\%) rejoint le
constat T-11. Le taux de signalement (15~\%) est encourageant et constitue
le meilleur indicateur à faire progresser. Aucune conclusion individuelle
n'est tirée.
\end{boxscenario}

\begin{boxobjet}{T-11 · Taux de clic élevé à l'hameçonnage simulé}
\textbf{Description~:} 14~clics sur 60 (23,3~\%) lors de la campagne
autorisée~; 5~débuts de saisie. \textbf{Preuve~:} statistiques agrégées de
la plateforme de test. \textbf{Gravité~:} moyenne. \textbf{Impact~:} un
courriel piégé réel ouvrirait un premier accès. \textbf{Recommandation~:}
sensibilisation régulière, bouton de signalement, second facteur,
marquage des courriels externes. \textbf{Détection~:} suivi du taux de
signalement et du taux de clic dans le temps.
\end{boxobjet}

% ============================================================================
\section{Les Parades}
\label{sec:c19-parades}
% ============================================================================

\begin{tableaucours}{L{4.4cm}L{7.8cm}}{%
    \tableauHeader{Parade} &
    \tableauHeader{Effet}
}
Sensibilisation régulière & Réduit le taux de clic, augmente le signalement \\
Bouton de signalement & Transforme chaque salarié en capteur \\
Second facteur & Un mot de passe saisi par erreur ne suffit plus \\
Marquage des courriels externes & Aide à repérer l'usurpation \\
Filtrage amont & Bloque une partie des courriels piégés \\
\end{tableaucours}
\vspace{-0.8em}
\begin{center}{\small\itshape Tableau~19.3. Parades à l'hameçonnage}\end{center}

\begin{boxmetier}{RSSI et analyste SOC}
\textbf{Faire du salarié un capteur.} L'objectif d'une campagne n'est pas de
faire baisser le taux de clic à zéro, impossible, mais de faire monter le
taux de signalement~: un salarié qui signale vite permet au SOC de réagir
avant la propagation.
\end{boxmetier}

\separateur

\begin{boxresume}
\textbf{Ce qu'il faut retenir de ce chapitre~:}
\begin{itemize}
    \item \textbf{L'ingénierie sociale} exploite autorité, urgence, peur,
          serviabilité, appât.
    \item \textbf{L'hameçonnage} prend plusieurs formes, de la masse au
          ciblé.
    \item \textbf{Une campagne éthique} informe, ne collecte pas les mots de
          passe, ne nomme personne.
    \item \textbf{On mesure un dispositif}, pas des individus.
    \item \textbf{Le taux de signalement} est le meilleur indicateur à faire
          progresser.
\end{itemize}
\end{boxresume}

\begin{boxinfo}
\textbf{Points de contrôle de l'élève.} Je sais~: définir l'ingénierie
sociale et ses leviers~; citer les formes d'hameçonnage~; conduire une
campagne éthique~; interpréter les résultats sans stigmatiser~; proposer des
parades.
\end{boxinfo}

\bigskip

\begin{boxexo}[title={\ding{45}\quad Exercice~19.1. Leviers \niveaubase}]
Associez à chaque courriel le levier exploité~: « la direction exige »~;
« compte fermé dans une heure ».
\end{boxexo}

\begin{boxexo}[title={\ding{45}\quad Exercice~19.2. Formes \niveaubase}]
Distinguez hameçonnage de masse et hameçonnage ciblé.
\end{boxexo}

\begin{boxexo}[title={\ding{45}\quad Exercice~19.3. Taux \niveaubase}]
Calculez le taux de clic et le taux de signalement de la campagne.
\end{boxexo}

\begin{boxexo}[title={\ding{45}\quad Exercice~19.4. Page pédagogique \niveaumoyen}]
Pourquoi la page d'atterrissage doit-elle être pédagogique et non
culpabilisante~?
\end{boxexo}

\begin{boxexo}[title={\ding{45}\quad Exercice~19.5. Collecte \niveaumoyen}]
Pourquoi ne conserve-t-on jamais les mots de passe réels saisis~?
\end{boxexo}

\begin{boxexo}[title={\ding{45}\quad Exercice~19.6. Signalement \niveaumoyen}]
Pourquoi le taux de signalement est-il un meilleur objectif que le taux de
clic~?
\end{boxexo}

\begin{boxexo}[title={\ding{45}\quad Exercice~19.7. T-11 \niveaumoyen}]
Rédigez le plan d'action lié à T-11.
\end{boxexo}

\begin{boxexo}[title={\ding{45}\quad Exercice~19.8. Thèmes sensibles \niveaumoyen}]
Pourquoi évite-t-on les thèmes sensibles (salaire, santé) dans un scénario
de test~?
\end{boxexo}

\begin{boxexo}[title={\ding{45}\quad Exercice~19.9. Campagne \niveauexpert}]
Rédigez la note d'information préalable du personnel annonçant le principe
d'une campagne, sans en révéler la date.
\end{boxexo}

\begin{boxexo}[title={\ding{45}\quad Exercice~19.10. Programme \niveauexpert}]
Proposez un programme annuel de sensibilisation et de mesure pour LiZbiZ.
\end{boxexo}

\bigskip

\begin{boxinfo}
\textbf{Pour aller plus loin.}
\begin{itemize}
    \item MITRE ATT\&CK~\cite{mitreattack}.
\end{itemize}
\end{boxinfo}

\begin{boxinfo}
\textbf{Transition.} L'ensemble de ces techniques se combine dans un
exercice réaliste~: le Chapitre~\ref{chap:rouge} traite de l'équipe rouge et
de l'équipe violette.
\end{boxinfo}
}{}
        \IfFileExists{partie5/P05_C20_equipe_rouge.tex}{% ============================================================================
% Fichier  : partie5/P05_C20_equipe_rouge.tex
% Titre    : Équipe Rouge et Équipe Violette
% Auteur   : MINKA MI NGUIDJOÏ Thierry Emmanuel
% Version  : 1.0 -- novembre 2026
% ============================================================================
% Sources : MITRE ATT&CK ; cadre de l'évaluation fondée sur les menaces ;
% distinction test d'intrusion / équipe rouge. Opération MBAM (septembre
% 2026, BANQUE SANAGA) : cas fictif, chiffres cohérents. Deux séances
% (S23 et S24).
% ============================================================================

\chapter{Équipe Rouge et Équipe Violette}
\label{chap:rouge}
\chaptermark{Équipe rouge et équipe violette}

\epigraph{%
    « L'équipe rouge ne cherche pas toutes les failles~;\\
    elle cherche à atteindre un objectif sans se faire voir. »%
}{Principe du cours}

\niveauChapitre{Expert}{Toute la Partie~III et la
Partie~IV~; Chapitre~\ref{chap:social}.}

\begin{boxobjectifs}
\begin{itemize}
    \item Distinguer test d'intrusion, équipe rouge et équipe violette.
    \item Décrire le déroulement d'un exercice d'équipe rouge par objectifs.
    \item Comprendre la notion de scénario fondé sur une menace.
    \item Mesurer la détection et le temps de réaction du défenseur.
    \item Conduire une restitution conjointe qui améliore la défense.
\end{itemize}
\end{boxobjectifs}

\bigskip

% ============================================================================
\section*{Question directrice : la BANQUE SANAGA voit-elle une attaque discrète, et en combien de temps~?}
% ============================================================================

En septembre, avant le test de LiZbiZ, NYONG SÉCURITÉ mène pour la BANQUE
SANAGA l'Opération MBAM, un exercice d'équipe rouge par objectifs. La
question n'est plus « combien de failles~? » mais « l'attaquant atteint-il
l'objectif, et le défenseur le voit-il~? ». Cet exercice sert de référence
au cours.

\medskip
\textbf{Plan du chapitre.}
\begin{enumerate}
    \item Trois exercices distincts (section~\ref{sec:c20-distinction}).
    \item L'exercice par objectifs (section~\ref{sec:c20-objectifs}).
    \item Le scénario fondé sur une menace (section~\ref{sec:c20-menace}).
    \item Mesurer la détection (section~\ref{sec:c20-detection}).
    \item L'équipe violette (section~\ref{sec:c20-violette}).
    \item L'Opération MBAM (section~\ref{sec:c20-mbam}).
\end{enumerate}

\separateur

% ============================================================================
\section{Trois Exercices Distincts}
\label{sec:c20-distinction}
% ============================================================================

\begin{tableaucours}{L{3.2cm}L{4.6cm}L{4.4cm}}{%
    \tableauHeader{Exercice} &
    \tableauHeader{But} &
    \tableauHeader{Discrétion}
}
Test d'intrusion & Trouver le plus de failles exploitables & Peu importe \\
Équipe rouge & Atteindre un objectif précis comme le ferait un attaquant &
Essentielle \\
Équipe violette & Rouge et bleue travaillent ensemble pour améliorer la
détection & Transparente \\
\end{tableaucours}
\vspace{-0.8em}
\begin{center}{\small\itshape Tableau~20.1. Trois types d'exercice offensif}\end{center}

% ============================================================================
\section{L'Exercice par Objectifs}
\label{sec:c20-objectifs}
% ============================================================================

\begin{boxdef}
\termecle{Équipe rouge}~: exercice où une équipe offensive cherche à
atteindre un objectif convenu (un « drapeau »), par tous les moyens du
périmètre, en restant discrète, pour évaluer la capacité de détection et de
réaction de l'organisation.
\end{boxdef}

Les objectifs, ou \termecle{drapeaux}, sont définis à l'avance avec un cercle
restreint de responsables~: atteindre un système critique, obtenir un droit
d'administration, accéder à une donnée précise. Le reste de l'organisation
n'est pas prévenu, pour mesurer une détection réelle.

% ============================================================================
\section{Le Scénario Fondé sur une Menace}
\label{sec:c20-menace}
% ============================================================================

Un bon exercice d'équipe rouge imite un attaquant plausible pour
l'organisation~: ses techniques, sa patience, ses objectifs. On s'appuie sur
la connaissance des menaces et sur les techniques d'ATT\&CK pour construire
un scénario réaliste, plutôt que de dérouler des outils au hasard.

\begin{boxalert}
\textbf{Violation typique.} Transformer l'exercice en démonstration de
force qui perturbe la production.

\smallskip
\textbf{Règle~:} l'équipe rouge garde un cercle de confiance qui peut tout
arrêter, respecte des lignes à ne pas franchir (pas d'arrêt de service, pas
de données réelles exfiltrées), et consigne chaque action avec son horaire
pour la restitution.
\end{boxalert}

% ============================================================================
\section{Mesurer la Détection}
\label{sec:c20-detection}
% ============================================================================

\begin{tableaucours}{L{4.8cm}L{7.4cm}}{%
    \tableauHeader{Indicateur} &
    \tableauHeader{Ce qu'il mesure}
}
Temps jusqu'au premier objectif & Rapidité de l'attaquant \\
Temps jusqu'à la détection & Rapidité du défenseur à voir \\
Temps jusqu'à la réaction & Rapidité à contenir \\
Techniques détectées sur techniques jouées & Couverture de la détection \\
Techniques bloquées & Efficacité des parades \\
\end{tableaucours}
\vspace{-0.8em}
\begin{center}{\small\itshape Tableau~20.2. Indicateurs d'un exercice d'équipe rouge}\end{center}

% ============================================================================
\section{L'Équipe Violette}
\label{sec:c20-violette}
% ============================================================================

\begin{boxdef}
\termecle{Équipe violette}~: démarche où l'équipe offensive (rouge) et
l'équipe de défense (bleue) travaillent ensemble, en transparence~: la rouge
joue une technique, la bleue vérifie si elle la détecte, et l'on corrige
aussitôt la détection manquante.
\end{boxdef}

C'est souvent la suite la plus utile d'un exercice rouge~: on rejoue les
techniques non détectées, on écrit les règles manquantes, on vérifie qu'elles
fonctionnent. Le lab HE2-LAB-09 fait construire cette matrice de couverture.

% ============================================================================
\section{L'Opération MBAM}
\label{sec:c20-mbam}
% ============================================================================

\begin{boxscenario}{Opération MBAM. Huit jours à la BANQUE SANAGA}
\textbf{Constatations} (septembre 2026). Objectifs~: F1 (accès au réseau
interne), F2 (administration du domaine), F3 (accès au système de paiement).
\begin{itemize}
    \item F1 atteint au 2\textsuperscript{e} jour, après un hameçonnage ciblé de
          12~employés (2~clics).
    \item F2 atteint au 8\textsuperscript{e} jour (administration du domaine).
    \item F3 non atteint~: le système de paiement, cloisonné, a résisté.
    \item Détection par le SOC au 9\textsuperscript{e} jour~; confinement au 10\textsuperscript{e}.
    \item 30~techniques d'ATT\&CK jouées~: 18~détectées (dont 11 après coup
          dans les journaux), 7~bloquées.
\end{itemize}

\medskip
\textbf{Analyse} (confiance élevée). La banque a résisté là où elle avait
cloisonné (F3)~; elle a vu tard (9~jours). La couverture de détection est de
60~\% (18 sur 30), mais une part n'a été vue qu'a posteriori. La valeur de
l'exercice tient à la phase violette qui a suivi~: les 12~techniques non
détectées en direct ont donné 12~règles nouvelles.
\end{boxscenario}

\begin{boxmetier}{Analyste SOC / Red Team et RSSI}
\textbf{La vraie livraison, c'est la détection.} Un exercice d'équipe rouge
réussi ne se mesure pas au nombre de drapeaux pris, mais aux règles de
détection écrites ensuite. Le RSSI valorise la phase violette autant que
l'intrusion.
\end{boxmetier}

\separateur

\begin{boxresume}
\textbf{Ce qu'il faut retenir de ce chapitre~:}
\begin{itemize}
    \item \textbf{Test d'intrusion} (toutes les failles), \textbf{équipe
          rouge} (un objectif, discrètement), \textbf{équipe violette}
          (ensemble).
    \item \textbf{Les drapeaux} sont définis à l'avance avec un cercle
          restreint.
    \item \textbf{Le scénario} imite une menace plausible.
    \item \textbf{On mesure} temps jusqu'à l'objectif, jusqu'à la détection,
          couverture.
    \item \textbf{La phase violette} transforme l'exercice en règles de
          détection.
\end{itemize}
\end{boxresume}

\begin{boxinfo}
\textbf{Points de contrôle de l'élève.} Je sais~: distinguer les trois
exercices~; décrire un exercice par objectifs~; expliquer le scénario fondé
sur une menace~; citer les indicateurs~; expliquer la valeur de l'équipe
violette.
\end{boxinfo}

\bigskip

\begin{boxexo}[title={\ding{45}\quad Exercice~20.1. Distinction \niveaubase}]
En quoi un exercice d'équipe rouge diffère-t-il d'un test d'intrusion~?
\end{boxexo}

\begin{boxexo}[title={\ding{45}\quad Exercice~20.2. Drapeaux \niveaubase}]
Qu'est-ce qu'un drapeau~? Qui le définit~?
\end{boxexo}

\begin{boxexo}[title={\ding{45}\quad Exercice~20.3. Couverture \niveaubase}]
Calculez la couverture de détection de l'Opération MBAM.
\end{boxexo}

\begin{boxexo}[title={\ding{45}\quad Exercice~20.4. F3 \niveaumoyen}]
Pourquoi le drapeau F3 n'a-t-il pas été atteint~? Quelle leçon en tirer~?
\end{boxexo}

\begin{boxexo}[title={\ding{45}\quad Exercice~20.5. Détecté après coup \niveaumoyen}]
Quelle différence entre une technique détectée en direct et une technique
retrouvée après coup dans les journaux~?
\end{boxexo}

\begin{boxexo}[title={\ding{45}\quad Exercice~20.6. Cercle de confiance \niveaumoyen}]
Pourquoi un cercle restreint doit-il pouvoir arrêter l'exercice à tout
moment~?
\end{boxexo}

\begin{boxexo}[title={\ding{45}\quad Exercice~20.7. Lignes rouges \niveaumoyen}]
Citez trois lignes à ne pas franchir pendant un exercice d'équipe rouge.
\end{boxexo}

\begin{boxexo}[title={\ding{45}\quad Exercice~20.8. Violette \niveaumoyen}]
Décrivez le déroulement d'une séance d'équipe violette à partir des
techniques non détectées de MBAM.
\end{boxexo}

\begin{boxexo}[title={\ding{45}\quad Exercice~20.9. Scénario \niveauexpert}]
Construisez un scénario d'équipe rouge fondé sur une menace plausible pour
LiZbiZ, avec trois drapeaux.
\end{boxexo}

\begin{boxexo}[title={\ding{45}\quad Exercice~20.10. Tableau de bord \niveauexpert}]
Proposez le tableau de bord d'indicateurs que la BANQUE SANAGA devrait
suivre d'un exercice à l'autre.
\end{boxexo}

\bigskip

\begin{boxinfo}
\textbf{Pour aller plus loin.}
\begin{itemize}
    \item MITRE ATT\&CK~\cite{mitreattack}.
\end{itemize}
\end{boxinfo}

\begin{boxinfo}
\textbf{Transition.} L'exercice rouge n'a de valeur que par la défense qu'il
nourrit~: le Chapitre~\ref{chap:detection} prend le point de vue du
défenseur.
\end{boxinfo}
}{}
        \IfFileExists{partie5/P05_C21_detection.tex}{% ============================================================================
% Fichier  : partie5/P05_C21_detection.tex
% Titre    : Se Faire Détecter : le Point de Vue du Défenseur
% Auteur   : MINKA MI NGUIDJOÏ Thierry Emmanuel
% Version  : 1.0 -- novembre 2026
% ============================================================================
% Sources : MITRE ATT&CK ; journaux de sécurité Windows ; principes de
% détection. Ce chapitre prend le point de vue défensif : il n'enseigne
% aucune technique d'évasion. Opérations LOGONE et MBAM : cas fictifs.
% ============================================================================

\chapter{Se Faire Détecter : le Point de Vue du Défenseur}
\label{chap:detection}
\chaptermark{Le point de vue du défenseur}

\epigraph{%
    « Le meilleur testeur offensif est celui qui apprend au défenseur\\
    à le voir la prochaine fois. »%
}{Principe du cours}

\niveauChapitre{Intermédiaire à expert}{Chapitre~\ref{chap:rouge}~;
Chapitre~\ref{chap:post}.}

\begin{boxobjectifs}
\begin{itemize}
    \item Comprendre comment un défenseur détecte chaque étape d'une
          attaque.
    \item Relier les techniques offensives étudiées à leurs traces.
    \item Construire une matrice de couverture de détection.
    \item Écrire une règle de détection à partir d'une technique jouée.
    \item Faire du testeur offensif un contributeur à la défense.
\end{itemize}
\end{boxobjectifs}

\bigskip

\begin{boiteRemarque}{Orientation défensive}
Ce chapitre prend délibérément le point de vue du défenseur. Il n'enseigne
pas à contourner la détection~: il enseigne à la mesurer et à l'améliorer.
Le testeur éthique joue ses techniques \emph{au vu et au su} du défenseur,
pour renforcer ce dernier.
\end{boiteRemarque}

% ============================================================================
\section*{Question directrice : pour chaque technique jouée, le défenseur disposait-il d'une trace et d'une règle~?}
% ============================================================================

Dans l'Opération MBAM, 30~techniques ont été jouées, 18~détectées, dont la
moitié seulement en direct. La question de ce chapitre est systématique~:
pour chacune des techniques étudiées dans ce cours, quelle est sa trace, et
quelle règle la transforme en alerte~?

\medskip
\textbf{Plan du chapitre.}
\begin{enumerate}
    \item Voir chaque étape (section~\ref{sec:c21-etapes}).
    \item Des techniques aux traces (section~\ref{sec:c21-traces}).
    \item La matrice de couverture (section~\ref{sec:c21-matrice}).
    \item Écrire une règle (section~\ref{sec:c21-regle}).
    \item Le testeur, allié du défenseur (section~\ref{sec:c21-allie}).
\end{enumerate}

\separateur

% ============================================================================
\section{Voir Chaque Étape}
\label{sec:c21-etapes}
% ============================================================================

\begin{tableaucours}{L{3.6cm}L{4.4cm}L{4.2cm}}{%
    \tableauHeader{Étape} &
    \tableauHeader{Trace typique} &
    \tableauHeader{Source}
}
Balayage & Nombreuses connexions refusées & Pare-feu, flux réseau \\
Hameçonnage & Courriel signalé, clic sur lien connu & Messagerie, mandataire
web \\
Authentification suspecte & Connexion inhabituelle, NTLM depuis un poste
anormal & Journaux du domaine \\
Mouvement latéral & Service distant installé, session réseau & Événements
7045, 4624 \\
Persistance & Compte créé, tâche planifiée & Événements 4720, 4698 \\
\end{tableaucours}
\vspace{-0.8em}
\begin{center}{\small\itshape Tableau~21.1. Voir chaque étape d'une attaque}\end{center}

% ============================================================================
\section{Des Techniques aux Traces}
\label{sec:c21-traces}
% ============================================================================

Chaque technique des Parties~III à~V laisse une trace, déjà signalée dans
son chapitre. Le défenseur regroupe ces traces par source de journal et
s'assure que chacune est bien collectée~: une trace non journalisée est une
technique invisible.

% ============================================================================
\section{La Matrice de Couverture}
\label{sec:c21-matrice}
% ============================================================================

\begin{boxdef}
\termecle{Matrice de couverture de détection}~: tableau qui croise les
techniques (par exemple celles d'ATT\&CK) avec l'état de la détection~: trace
collectée, règle existante, règle testée. Elle montre les trous.
\end{boxdef}

\begin{boxprocedure}
\textbf{Construire une matrice de couverture}
\begin{enumerate}
    \item Lister les techniques jouées ou attendues.
    \item Pour chacune, noter la trace et sa source.
    \item Indiquer si la trace est collectée.
    \item Indiquer si une règle existe, et si elle a été testée.
    \item Prioriser les trous selon la gravité de la technique.
\end{enumerate}
\end{boxprocedure}

% ============================================================================
\section{Écrire une Règle}
\label{sec:c21-regle}
% ============================================================================

Une règle de détection s'exprime d'abord en français~: \emph{quelle
condition, sur quelle source, déclenche quelle alerte}. On la teste en
rejouant la technique (phase violette), on mesure les fausses alertes, puis
on l'ajuste. Le lab HE2-LAB-09 fait écrire et tester de telles règles.

\begin{boxscenario}{Opération LOGONE. Les règles qui manquaient}
\textbf{Constatation.} Les traces de la connexion anormale du compte de la
facturation et de son ajout à un groupe à privilèges existaient, mais aucune
règle ne les transformait en alerte.

\medskip
\textbf{Analyse} (confiance élevée). Après restitution, deux règles ont été
écrites et testées~: « compte à privilèges se connectant depuis un poste
ordinaire » et « ajout d'un compte à un groupe sensible hors fenêtre de
maintenance ». La détection ne coûtait pas d'outil nouveau, seulement des
règles.
\end{boxscenario}

% ============================================================================
\section{Le Testeur, Allié du Défenseur}
\label{sec:c21-allie}
% ============================================================================

\begin{boxmetier}{Analyste SOC / Red Team}
\textbf{Livrer les traces, pas seulement les failles.} À la fin d'une
mission, l'analyste remet au défenseur la liste horodatée de ses actions~:
c'est le meilleur jeu de test pour vérifier et compléter la détection.
\end{boxmetier}

\begin{boxalert}
\textbf{Violation typique.} Se vanter d'être « resté invisible » sans aider
le défenseur à le devenir moins.

\smallskip
\textbf{Règle~:} rester discret est un moyen de mesurer la détection, pas une
fin. Le testeur éthique livre toujours ce qu'il faut pour que le défenseur
le détecte la fois suivante.
\end{boxalert}

\separateur

\begin{boxresume}
\textbf{Ce qu'il faut retenir de ce chapitre~:}
\begin{itemize}
    \item \textbf{Chaque étape} d'une attaque laisse une trace, dans une
          source précise.
    \item \textbf{Une trace non collectée} rend la technique invisible.
    \item \textbf{La matrice de couverture} montre les trous de détection.
    \item \textbf{Une règle} s'énonce en français, se teste, s'ajuste.
    \item \textbf{Le testeur} livre ses actions horodatées pour nourrir la
          défense.
\end{itemize}
\end{boxresume}

\begin{boxinfo}
\textbf{Points de contrôle de l'élève.} Je sais~: relier une étape à sa
trace et à sa source~; expliquer pourquoi une trace doit être collectée~;
construire une matrice de couverture~; énoncer une règle de détection~;
décrire l'apport du testeur à la défense.
\end{boxinfo}

\bigskip

\begin{boxexo}[title={\ding{45}\quad Exercice~21.1. Traces \niveaubase}]
Donnez la trace et la source de~: un balayage~; la création d'un compte.
\end{boxexo}

\begin{boxexo}[title={\ding{45}\quad Exercice~21.2. Invisible \niveaubase}]
Pourquoi une technique dont la trace n'est pas collectée est-elle
invisible~?
\end{boxexo}

\begin{boxexo}[title={\ding{45}\quad Exercice~21.3. Matrice \niveaubase}]
Quelles colonnes comporte une matrice de couverture de détection~?
\end{boxexo}

\begin{boxexo}[title={\ding{45}\quad Exercice~21.4. Règle en français \niveaumoyen}]
Énoncez en français la règle qui détecte un compte à privilèges se
connectant depuis un poste ordinaire.
\end{boxexo}

\begin{boxexo}[title={\ding{45}\quad Exercice~21.5. Fausses alertes \niveaumoyen}]
Pourquoi faut-il mesurer les fausses alertes d'une règle avant de la
déployer~?
\end{boxexo}

\begin{boxexo}[title={\ding{45}\quad Exercice~21.6. En direct ou après coup \niveaumoyen}]
Que faut-il changer pour détecter en direct une technique aujourd'hui vue
seulement après coup~?
\end{boxexo}

\begin{boxexo}[title={\ding{45}\quad Exercice~21.7. LOGONE \niveaumoyen}]
À partir du scénario, proposez une troisième règle utile pour LiZbiZ.
\end{boxexo}

\begin{boxexo}[title={\ding{45}\quad Exercice~21.8. Livraison \niveaumoyen}]
Pourquoi la liste horodatée des actions du testeur est-elle un bon jeu de
test pour le SOC~?
\end{boxexo}

\begin{boxexo}[title={\ding{45}\quad Exercice~21.9. Couverture complète \niveauexpert}]
Construisez la matrice de couverture des techniques des Chapitres~\ref{chap:reseau}
à~\ref{chap:post}.
\end{boxexo}

\begin{boxexo}[title={\ding{45}\quad Exercice~21.10. Programme \niveauexpert}]
Proposez un programme d'amélioration continue de la détection fondé sur des
exercices violets trimestriels.
\end{boxexo}

\bigskip

\begin{boxinfo}
\textbf{Pour aller plus loin.}
\begin{itemize}
    \item MITRE ATT\&CK~\cite{mitreattack}.
\end{itemize}
\end{boxinfo}

\begin{boxinfo}
\textbf{Transition.} Toutes ces conclusions doivent être restituées~: la
Partie~VI commence par le Chapitre~\ref{chap:rapport}, consacré au rapport.
\end{boxinfo}
}{}

    \part{Restituer et Progresser}
    \label{part:restituer}

        \IfFileExists{partie6/P06_C22_rapport.tex}{% ============================================================================
% Fichier  : partie6/P06_C22_rapport.tex
% Titre    : Le Rapport de Test d'Intrusion
% Auteur   : MINKA MI NGUIDJOÏ Thierry Emmanuel
% Version  : 1.0 -- novembre 2026
% ============================================================================
% Sources : PTES (phase de restitution) ; CVSS 3.1 pour la cotation ;
% pratiques de rapport. Les 14 constats de l'Opération LOGONE : cas fictif,
% cohérents avec les chapitres précédents.
% ============================================================================

\chapter{Le Rapport de Test d'Intrusion}
\label{chap:rapport}
\chaptermark{Le rapport de test d'intrusion}

\epigraph{%
    « Un test d'intrusion ne vaut que par son rapport~;\\
    le reste n'est qu'une collection de preuves. »%
}{Principe du cours}

\niveauChapitre{Intermédiaire}{Toutes les parties
précédentes~; Chapitre~\ref{chap:methodes}.}

\begin{boxobjectifs}
\begin{itemize}
    \item Décrire la structure d'un rapport de test d'intrusion.
    \item Rédiger une synthèse pour la direction et un détail technique.
    \item Coter chaque constat et le hiérarchiser.
    \item Rédiger un constat complet~: description, preuve, impact,
          recommandation.
    \item Présenter les 14~constats de l'Opération LOGONE comme modèle.
\end{itemize}
\end{boxobjectifs}

\bigskip

% ============================================================================
\section*{Question directrice : comment transformer dix jours de test en un document utile à la direction et aux techniciens~?}
% ============================================================================

Le test de LiZbiZ est terminé~: 14~constats, des preuves, des notes. Il faut
maintenant un document que M.~ETEME lira en dix minutes et que M.~KOUNOU
suivra ligne à ligne pour corriger. Ce chapitre construit ce document.

\medskip
\textbf{Plan du chapitre.}
\begin{enumerate}
    \item La structure du rapport (section~\ref{sec:c22-structure}).
    \item Deux lecteurs, deux niveaux (section~\ref{sec:c22-lecteurs}).
    \item Coter et hiérarchiser (section~\ref{sec:c22-cotation}).
    \item L'anatomie d'un constat (section~\ref{sec:c22-constat}).
    \item Les 14 constats de LOGONE (section~\ref{sec:c22-logone}).
\end{enumerate}

\separateur

% ============================================================================
\section{La Structure du Rapport}
\label{sec:c22-structure}
% ============================================================================

\begin{tableaucours}{L{4.0cm}L{8.2cm}}{%
    \tableauHeader{Partie} &
    \tableauHeader{Contenu}
}
Synthèse pour la direction & Verdict, risques majeurs, recommandations
prioritaires, sans jargon \\
Contexte et périmètre & Objectifs, limites, dates, règles d'engagement \\
Démarche & Méthode suivie, outils, degré de couverture \\
Constats détaillés & Un par faiblesse, coté et documenté \\
Plan de correction & Actions priorisées, responsables, échéances \\
Annexes & Preuves, liste des actions horodatées, périmètre technique \\
\end{tableaucours}
\vspace{-0.8em}
\begin{center}{\small\itshape Tableau~22.1. Structure d'un rapport de test d'intrusion}\end{center}

% ============================================================================
\section{Deux Lecteurs, Deux Niveaux}
\label{sec:c22-lecteurs}
% ============================================================================

\begin{boxalert}
\textbf{Violation typique.} Un rapport uniquement technique, illisible pour
la direction, ou une synthèse vague, inutilisable par les techniciens.

\smallskip
\textbf{Règle~:} la synthèse parle risque et décision~; le détail parle
technique et correction. Les deux disent la même vérité, à deux niveaux.
\end{boxalert}

% ============================================================================
\section{Coter et Hiérarchiser}
\label{sec:c22-cotation}
% ============================================================================

Chaque constat reçoit une gravité. On peut s'appuyer sur le système commun
de notation des vulnérabilités (CVSS~3.1, Chapitre~\ref{chap:vulnerabilites})
pour les failles techniques, et sur une appréciation du risque métier pour
le reste. La gravité combine la facilité d'exploitation et l'impact.

\begin{tableaucours}{L{2.4cm}L{3.0cm}L{6.6cm}}{%
    \tableauHeader{Gravité} &
    \tableauHeader{Nombre} &
    \tableauHeader{Exigence de traitement}
}
Critique & 2 & Correction immédiate \\
Élevée & 4 & Sous quelques jours \\
Moyenne & 5 & Planifiée au trimestre \\
Faible & 3 & Au fil de l'eau \\
\end{tableaucours}
\vspace{-0.8em}
\begin{center}{\small\itshape Tableau~22.2. Répartition des 14 constats de LOGONE}\end{center}

% ============================================================================
\section{L'Anatomie d'un Constat}
\label{sec:c22-constat}
% ============================================================================

\begin{boxprocedure}
\textbf{Rédiger un constat}
\begin{enumerate}
    \item Identifiant et titre court.
    \item Description~: la faiblesse, en termes clairs.
    \item Preuve~: ce qui l'établit (requête, capture masquée, relevé).
    \item Impact~: ce qu'un attaquant en ferait.
    \item Gravité~: cotation justifiée.
    \item Recommandation~: la correction, concrète et vérifiable.
    \item Détection~: comment voir l'exploitation en production.
\end{enumerate}
\end{boxprocedure}

Tous les encadrés « Fiche de constat » des chapitres précédents suivent
exactement cette anatomie.

% ============================================================================
\section{Les 14 Constats de LOGONE}
\label{sec:c22-logone}
% ============================================================================

\begin{tableaucours}{L{1.2cm}L{6.4cm}L{2.0cm}L{2.2cm}}{%
    \tableauHeader{Id} &
    \tableauHeader{Constat} &
    \tableauHeader{Gravité} &
    \tableauHeader{Chapitre}
}
T-01 & Compte de service vulnérable au vol de ticket (facturation) &
Critique & \ref{chap:ad} \\
T-02 & Mot de passe administrateur local en clair dans un script partagé &
Critique & \ref{chap:privileges} \\
T-03 & Accès distant sans second facteur & Élevée & \ref{chap:vulnerabilites} \\
T-05 & Accès indirect aux colis d'autrui (portail) & Élevée &
\ref{chap:injections} \\
T-06 & Protocole de partage ancien et vulnérable & Élevée &
\ref{chap:exploitation} \\
T-09 & Exécution de code sur l'ancien portail & Élevée &
\ref{chap:exploitation} \\
T-04 & Identifiants par défaut sur des équipements & Moyenne &
\ref{chap:sansfil} \\
T-07 & Signature de partage non exigée et résolutions locales actives &
Moyenne & \ref{chap:reseau} \\
T-08 & Comptes d'administration utilisés sur des postes ordinaires & Moyenne
& \ref{chap:post} \\
T-10 & Service de test exposé sur l'Internet & Moyenne & \ref{chap:nuage} \\
T-11 & Taux de clic élevé à l'hameçonnage simulé & Moyenne &
\ref{chap:social} \\
T-12 & En-têtes de sécurité HTTP absents & Faible & \ref{chap:web} \\
T-13 & Bannières de version exposées & Faible & \ref{chap:balayage} \\
T-14 & Délai d'expiration de session trop long & Faible & \ref{chap:web} \\
\end{tableaucours}
\vspace{-0.8em}
\begin{center}{\small\itshape Tableau~22.3. Les 14 constats de l'Opération LOGONE}\end{center}

L'audit des mots de passe (constat A-01, Chapitre~\ref{chap:mdp}) est
rapporté séparément, comme mission complémentaire.

\begin{boxmetier}{Pentesteur et RSSI}
\textbf{Écrire pour l'action.} Un constat réussi est celui que le
destinataire sait corriger sans rappeler l'auteur. La recommandation est
concrète, vérifiable, et distingue la correction de fond du palliatif
immédiat.
\end{boxmetier}

\separateur

\begin{boxresume}
\textbf{Ce qu'il faut retenir de ce chapitre~:}
\begin{itemize}
    \item \textbf{Le rapport} est le livrable~; tout le reste le prépare.
    \item \textbf{Deux niveaux}~: synthèse pour la direction, détail pour les
          techniciens.
    \item \textbf{Chaque constat} est coté et hiérarchisé.
    \item \textbf{L'anatomie}~: description, preuve, impact, gravité,
          recommandation, détection.
    \item \textbf{Les 14 constats de LOGONE} suivent ce modèle.
\end{itemize}
\end{boxresume}

\begin{boxinfo}
\textbf{Points de contrôle de l'élève.} Je sais~: donner la structure d'un
rapport~; écrire à deux niveaux~; coter et hiérarchiser~; rédiger un constat
complet~; situer chaque constat de LOGONE.
\end{boxinfo}

\bigskip

\begin{boxexo}[title={\ding{45}\quad Exercice~22.1. Structure \niveaubase}]
Citez les six parties d'un rapport de test d'intrusion.
\end{boxexo}

\begin{boxexo}[title={\ding{45}\quad Exercice~22.2. Deux lecteurs \niveaubase}]
Qu'attend la direction du rapport~? Qu'attend le technicien~?
\end{boxexo}

\begin{boxexo}[title={\ding{45}\quad Exercice~22.3. Répartition \niveaubase}]
Combien de constats critiques, élevés, moyens et faibles compte LOGONE~?
\end{boxexo}

\begin{boxexo}[title={\ding{45}\quad Exercice~22.4. Anatomie \niveaumoyen}]
Rédigez le constat T-06 selon l'anatomie en sept points.
\end{boxexo}

\begin{boxexo}[title={\ding{45}\quad Exercice~22.5. Synthèse \niveaumoyen}]
Rédigez la synthèse pour la direction de l'Opération LOGONE, en dix lignes.
\end{boxexo}

\begin{boxexo}[title={\ding{45}\quad Exercice~22.6. Priorité \niveaumoyen}]
Justifiez pourquoi T-01 et T-02 passent avant T-12 et T-13.
\end{boxexo}

\begin{boxexo}[title={\ding{45}\quad Exercice~22.7. Mission complémentaire \niveaumoyen}]
Pourquoi l'audit des mots de passe (A-01) est-il rapporté à part des
14~constats~?
\end{boxexo}

\begin{boxexo}[title={\ding{45}\quad Exercice~22.8. Recommandation \niveaumoyen}]
Réécrivez une recommandation vague (« améliorer la sécurité ») en
recommandation concrète et vérifiable.
\end{boxexo}

\begin{boxexo}[title={\ding{45}\quad Exercice~22.9. Plan de correction \niveauexpert}]
Construisez le plan de correction de LOGONE~: ordre, responsables, échéances.
\end{boxexo}

\begin{boxexo}[title={\ding{45}\quad Exercice~22.10. Rapport complet \niveauexpert}]
Rédigez le plan complet du rapport de LOGONE, partie par partie, avec le
nombre de pages visé.
\end{boxexo}

\bigskip

\begin{boxinfo}
\textbf{Pour aller plus loin.}
\begin{itemize}
    \item PTES~\cite{ptes}~; CVSS~3.1~\cite{first2019cvss31}.
\end{itemize}
\end{boxinfo}

\begin{boxinfo}
\textbf{Transition.} Un rapport sans vérification des corrections reste
théorique~: le Chapitre~\ref{chap:retest} traite du contre-test.
\end{boxinfo}
}{}
        \IfFileExists{partie6/P06_C23_contre_test.tex}{% ============================================================================
% Fichier  : partie6/P06_C23_contre_test.tex
% Titre    : Le Contre-Test
% Auteur   : MINKA MI NGUIDJOÏ Thierry Emmanuel
% Version  : 1.0 -- novembre 2026
% ============================================================================
% Sources : pratiques de retest ; gestion des vulnérabilités. Contre-test de
% l'Opération LOGONE (novembre 2026) : cas fictif.
% ============================================================================

\chapter{Le Contre-Test}
\label{chap:retest}
\chaptermark{Le contre-test}

\epigraph{%
    « Une faille n'est corrigée que lorsqu'un second test le prouve. »%
}{Principe du cours}

\niveauChapitre{Intermédiaire}{Chapitre~\ref{chap:rapport}.}

\begin{boxobjectifs}
\begin{itemize}
    \item Définir le contre-test et sa place dans le cycle.
    \item Vérifier qu'une correction ferme réellement la faille.
    \item Distinguer correction de fond et palliatif.
    \item Documenter l'état de chaque constat après correction.
    \item Clore une mission proprement.
\end{itemize}
\end{boxobjectifs}

\bigskip

% ============================================================================
\section*{Question directrice : les corrections de LiZbiZ ont-elles vraiment fermé les failles~?}
% ============================================================================

Trois semaines après le rapport, M.~KOUNOU annonce que « tout est corrigé ».
NYONG SÉCURITÉ revient vérifier, constat par constat. Certaines failles sont
bien fermées, d'autres seulement masquées, une a été rouverte par une
nouvelle configuration. Le contre-test établit la vérité.

\medskip
\textbf{Plan du chapitre.}
\begin{enumerate}
    \item Qu'est-ce qu'un contre-test (section~\ref{sec:c23-definition}).
    \item Fond ou palliatif (section~\ref{sec:c23-fond}).
    \item L'état des constats (section~\ref{sec:c23-etats}).
    \item Clore la mission (section~\ref{sec:c23-cloture}).
\end{enumerate}

\separateur

% ============================================================================
\section{Qu'est-ce qu'un Contre-Test}
\label{sec:c23-definition}
% ============================================================================

\begin{boxdef}
\termecle{Contre-test}~: vérification, après correction, que chaque faiblesse
rapportée a bien été traitée, en rejouant le test qui l'avait révélée, dans
les mêmes conditions.
\end{boxdef}

% ============================================================================
\section{Fond ou Palliatif}
\label{sec:c23-fond}
% ============================================================================

\begin{boxalert}
\textbf{Violation typique.} Déclarer une faille corrigée parce qu'un symptôme
a disparu.

\smallskip
\textbf{Règle~:} on vérifie la cause, pas le symptôme. Masquer une bannière
de version ne corrige pas la version~; bloquer une adresse ne corrige pas le
défaut d'accès. On distingue clairement la correction de fond du palliatif
temporaire.
\end{boxalert}

% ============================================================================
\section{L'État des Constats}
\label{sec:c23-etats}
% ============================================================================

\begin{boxscenario}{Opération LOGONE. Le contre-test de novembre}
\textbf{Constatations.} Sur les 14~constats~: 9~fermés (dont T-01 et T-02,
critiques), 3~partiellement traités (T-07 corrigé sur les serveurs mais pas
les postes~; T-08 en cours~; T-11, sensibilisation engagée), 1~palliatif
(T-13, bannières masquées mais versions inchangées), 1~rouvert (T-10, un
nouveau service de test exposé).

\medskip
\textbf{Analyse} (confiance élevée). Les deux failles critiques sont bien
fermées. Le constat rouvert montre l'intérêt d'un contre-test~: sans lui, la
réexposition serait passée inaperçue.
\end{boxscenario}

\begin{tableaucours}{L{3.2cm}L{9.0cm}}{%
    \tableauHeader{État} &
    \tableauHeader{Signification}
}
Fermé & La faille ne se reproduit plus~; correction de fond vérifiée \\
Partiel & Une partie du périmètre reste exposée \\
Palliatif & Symptôme masqué, cause non traitée \\
Rouvert & La faille, fermée, est réapparue \\
Non traité & Aucune correction \\
\end{tableaucours}
\vspace{-0.8em}
\begin{center}{\small\itshape Tableau~23.1. États d'un constat au contre-test}\end{center}

% ============================================================================
\section{Clore la Mission}
\label{sec:c23-cloture}
% ============================================================================

\begin{boxprocedure}
\textbf{Clore une mission}
\begin{enumerate}
    \item Rejouer chaque test et noter l'état du constat.
    \item Rédiger un addendum au rapport avec les états.
    \item Distinguer fond et palliatif pour chaque correction.
    \item Convenir d'un suivi pour les constats partiels.
    \item Détruire les données de test et restituer les accès.
\end{enumerate}
\end{boxprocedure}

\begin{boxmetier}{Auditeur de sécurité et RSSI}
\textbf{Le contre-test ferme la boucle.} Sans contre-test, un rapport reste
une liste de promesses. Le RSSI inscrit le contre-test au contrat et en fait
la condition de clôture.
\end{boxmetier}

\separateur

\begin{boxresume}
\textbf{Ce qu'il faut retenir de ce chapitre~:}
\begin{itemize}
    \item \textbf{Le contre-test} rejoue le test après correction.
    \item \textbf{On vérifie la cause}, pas le symptôme.
    \item \textbf{Chaque constat} reçoit un état~: fermé, partiel, palliatif,
          rouvert, non traité.
    \item \textbf{Un constat rouvert} justifie à lui seul le contre-test.
    \item \textbf{La clôture} documente les états et détruit les données.
\end{itemize}
\end{boxresume}

\begin{boxinfo}
\textbf{Points de contrôle de l'élève.} Je sais~: définir le contre-test~;
distinguer fond et palliatif~; attribuer un état à un constat~; rédiger un
addendum~; clore une mission.
\end{boxinfo}

\bigskip

\begin{boxexo}[title={\ding{45}\quad Exercice~23.1. Définition \niveaubase}]
Qu'est-ce qu'un contre-test~?
\end{boxexo}

\begin{boxexo}[title={\ding{45}\quad Exercice~23.2. Fond ou palliatif \niveaubase}]
Masquer une bannière de version est-il une correction de fond ou un
palliatif~? Pourquoi~?
\end{boxexo}

\begin{boxexo}[title={\ding{45}\quad Exercice~23.3. États \niveaubase}]
Donnez les cinq états possibles d'un constat au contre-test.
\end{boxexo}

\begin{boxexo}[title={\ding{45}\quad Exercice~23.4. T-10 rouvert \niveaumoyen}]
Comment un constat fermé peut-il être rouvert~? Que recommander~?
\end{boxexo}

\begin{boxexo}[title={\ding{45}\quad Exercice~23.5. Partiel \niveaumoyen}]
T-07 est corrigé sur les serveurs mais pas les postes. Quel état, et quel
suivi~?
\end{boxexo}

\begin{boxexo}[title={\ding{45}\quad Exercice~23.6. Addendum \niveaumoyen}]
Rédigez l'addendum de contre-test pour T-01 et T-13.
\end{boxexo}

\begin{boxexo}[title={\ding{45}\quad Exercice~23.7. Condition de clôture \niveaumoyen}]
Pourquoi faire du contre-test une condition de clôture au contrat~?
\end{boxexo}

\begin{boxexo}[title={\ding{45}\quad Exercice~23.8. Bilan \niveaumoyen}]
Dressez le bilan chiffré du contre-test de LOGONE.
\end{boxexo}

\begin{boxexo}[title={\ding{45}\quad Exercice~23.9. Procédure \niveauexpert}]
Proposez une procédure de contre-test réutilisable pour les futures missions
de NYONG SÉCURITÉ.
\end{boxexo}

\begin{boxexo}[title={\ding{45}\quad Exercice~23.10. Indicateur \niveauexpert}]
Définissez un indicateur de qualité de correction, calculé à partir des
états de contre-test.
\end{boxexo}

\bigskip

\begin{boxinfo}
\textbf{Pour aller plus loin.}
\begin{itemize}
    \item PTES~\cite{ptes}.
\end{itemize}
\end{boxinfo}

\begin{boxinfo}
\textbf{Transition.} Lorsqu'une faille touche un produit utilisé par
d'autres, sa correction dépasse le client~: le Chapitre~\ref{chap:divulgation}
traite de la divulgation coordonnée.
\end{boxinfo}
}{}
        \IfFileExists{partie6/P06_C24_divulgation.tex}{% ============================================================================
% Fichier  : partie6/P06_C24_divulgation.tex
% Titre    : La Divulgation Coordonnée
% Auteur   : MINKA MI NGUIDJOÏ Thierry Emmanuel
% Version  : 1.0 -- novembre 2026
% ============================================================================
% Sources : ISO/IEC 29147 (divulgation de vulnérabilités) ; ISO/IEC 30111
% (traitement des vulnérabilités) ; loi camerounaise 2010/012. Cas d'une
% faille découverte dans un produit tiers : fictif.
% ============================================================================

\chapter{La Divulgation Coordonnée}
\label{chap:divulgation}
\chaptermark{La divulgation coordonnée}

\epigraph{%
    « Trouver une faille crée un devoir~:\\
    celui de la faire corriger sans nuire. »%
}{Principe du cours}

\niveauChapitre{Intermédiaire}{Chapitre~\ref{chap:droit}~;
Chapitre~\ref{chap:rapport}.}

\begin{boxobjectifs}
\begin{itemize}
    \item Définir la divulgation coordonnée de vulnérabilités.
    \item Distinguer divulgation privée, coordonnée et complète.
    \item Décrire le processus des normes ISO/IEC 29147 et 30111.
    \item Situer la divulgation dans le cadre juridique camerounais.
    \item Rédiger un signalement responsable.
\end{itemize}
\end{boxobjectifs}

\bigskip

% ============================================================================
\section*{Question directrice : que faire d'une faille découverte dans un logiciel utilisé par beaucoup d'autres~?}
% ============================================================================

En testant LiZbiZ, M.~TAMO découvre que la faille de l'ancien portail (T-09)
vient d'un composant répandu, utilisé par de nombreuses organisations. La
corriger chez LiZbiZ ne protège pas les autres. Comment signaler la faille à
l'éditeur sans exposer tout le monde~? C'est l'objet de la divulgation
coordonnée.

\medskip
\textbf{Plan du chapitre.}
\begin{enumerate}
    \item Trois façons de divulguer (section~\ref{sec:c24-trois}).
    \item Le processus coordonné (section~\ref{sec:c24-processus}).
    \item Les normes 29147 et 30111 (section~\ref{sec:c24-normes}).
    \item Le cadre juridique (section~\ref{sec:c24-droit}).
    \item Rédiger un signalement (section~\ref{sec:c24-signalement}).
\end{enumerate}

\separateur

% ============================================================================
\section{Trois Façons de Divulguer}
\label{sec:c24-trois}
% ============================================================================

\begin{tableaucours}{L{3.4cm}L{8.6cm}}{%
    \tableauHeader{Façon} &
    \tableauHeader{Principe et limite}
}
Privée & On informe le seul éditeur~; sûr, mais sans pression si l'éditeur
n'agit pas \\
Coordonnée & On informe l'éditeur, puis on publie après correction ou
délai convenu \\
Complète et immédiate & On publie tout aussitôt~; expose les utilisateurs
avant correctif \\
\end{tableaucours}
\vspace{-0.8em}
\begin{center}{\small\itshape Tableau~24.1. Trois modes de divulgation}\end{center}

La divulgation \termecle{coordonnée} est le compromis responsable~: elle
donne à l'éditeur le temps de corriger, puis informe les utilisateurs pour
qu'ils se protègent.

% ============================================================================
\section{Le Processus Coordonné}
\label{sec:c24-processus}
% ============================================================================

\begin{boxprocedure}
\textbf{Divulguer de façon coordonnée}
\begin{enumerate}
    \item Trouver le contact de sécurité de l'éditeur.
    \item Signaler la faille par un canal sûr, avec les éléments de
          reproduction.
    \item Convenir d'un délai de correction raisonnable.
    \item Laisser l'éditeur corriger et, le cas échéant, obtenir un
          identifiant public de vulnérabilité.
    \item Publier une description après le correctif, sans mode d'emploi
          d'attaque inutile.
\end{enumerate}
\end{boxprocedure}

\begin{boxalert}
\textbf{Violation typique.} Publier la faille avec son code d'exploitation
avant tout correctif, « pour forcer l'éditeur ».

\smallskip
\textbf{Règle~:} la divulgation vise à protéger, pas à nuire ni à se mettre
en avant. On laisse un délai raisonnable, on ne publie pas d'arme prête à
l'emploi, on pense d'abord aux utilisateurs exposés.
\end{boxalert}

% ============================================================================
\section{Les Normes 29147 et 30111}
\label{sec:c24-normes}
% ============================================================================

\begin{tableaucours}{L{4.0cm}L{8.2cm}}{%
    \tableauHeader{Norme} &
    \tableauHeader{Objet}
}
ISO/IEC 29147 & Comment \emph{recevoir} et publier les signalements de
vulnérabilités (côté celui qui divulgue et celui qui reçoit) \\
ISO/IEC 30111 & Comment \emph{traiter} en interne une vulnérabilité signalée
jusqu'à sa correction \\
\end{tableaucours}
\vspace{-0.8em}
\begin{center}{\small\itshape Tableau~24.2. Deux normes complémentaires~\cite{iso29147,iso30111}}\end{center}

Les deux normes sont complémentaires~: 29147 organise l'échange avec le
monde extérieur, 30111 organise le travail interne de correction.

% ============================================================================
\section{Le Cadre Juridique}
\label{sec:c24-droit}
% ============================================================================

\begin{boxjuris}
La découverte d'une faille ne vaut pas autorisation d'y accéder. Au
Cameroun, la loi 2010/012 sur la cybersécurité réprime l'accès et le
maintien frauduleux dans un système, même sans dommage
(Chapitre~\ref{chap:droit}). Un chercheur qui découvre une faille hors d'un
mandat doit s'en tenir à ce qu'il a constaté sans intrusion supplémentaire,
et signaler.
\end{boxjuris}

% ============================================================================
\section{Rédiger un Signalement}
\label{sec:c24-signalement}
% ============================================================================

\begin{boxmetier}{Pentesteur et RSSI}
\textbf{Signaler avec mesure.} Un bon signalement est factuel, reproductible
et courtois~: il décrit la faille, son impact, les conditions de
reproduction, et propose un délai. Il n'exige rien, il coopère.
\end{boxmetier}

\separateur

\begin{boxresume}
\textbf{Ce qu'il faut retenir de ce chapitre~:}
\begin{itemize}
    \item \textbf{Trouver une faille} crée le devoir de la faire corriger
          sans nuire.
    \item \textbf{La divulgation coordonnée} donne le temps de corriger, puis
          informe.
    \item \textbf{29147} organise l'échange, \textbf{30111} le traitement
          interne.
    \item \textbf{Découvrir n'autorise pas à accéder}~: la loi 2010/012
          s'applique.
    \item \textbf{Un signalement} est factuel, reproductible, coopératif.
\end{itemize}
\end{boxresume}

\begin{boxinfo}
\textbf{Points de contrôle de l'élève.} Je sais~: distinguer les trois modes
de divulgation~; dérouler le processus coordonné~; distinguer 29147 et
30111~; rappeler le cadre juridique~; rédiger un signalement.
\end{boxinfo}

\bigskip

\begin{boxexo}[title={\ding{45}\quad Exercice~24.1. Trois modes \niveaubase}]
Distinguez divulgation privée, coordonnée et complète immédiate.
\end{boxexo}

\begin{boxexo}[title={\ding{45}\quad Exercice~24.2. Normes \niveaubase}]
Quelle norme organise la réception des signalements~? Laquelle organise leur
traitement interne~?
\end{boxexo}

\begin{boxexo}[title={\ding{45}\quad Exercice~24.3. Devoir \niveaubase}]
Pourquoi corriger T-09 chez LiZbiZ ne suffit-il pas~?
\end{boxexo}

\begin{boxexo}[title={\ding{45}\quad Exercice~24.4. Délai \niveaumoyen}]
Pourquoi laisse-t-on un délai à l'éditeur avant publication~? Qu'est-ce qu'un
délai raisonnable~?
\end{boxexo}

\begin{boxexo}[title={\ding{45}\quad Exercice~24.5. Pas d'arme \niveaumoyen}]
Pourquoi ne publie-t-on pas le code d'exploitation prêt à l'emploi~?
\end{boxexo}

\begin{boxexo}[title={\ding{45}\quad Exercice~24.6. Droit \niveaumoyen}]
Un chercheur découvre une faille sur un site public. Que peut-il faire, que
ne peut-il pas faire, selon la loi 2010/012~?
\end{boxexo}

\begin{boxexo}[title={\ding{45}\quad Exercice~24.7. Contact \niveaumoyen}]
Comment trouve-t-on le contact de sécurité d'un éditeur~?
\end{boxexo}

\begin{boxexo}[title={\ding{45}\quad Exercice~24.8. Signalement \niveaumoyen}]
Rédigez le signalement de la faille de composant à l'éditeur.
\end{boxexo}

\begin{boxexo}[title={\ding{45}\quad Exercice~24.9. Politique \niveauexpert}]
Rédigez une politique de divulgation que LiZbiZ pourrait publier pour
recevoir des signalements, inspirée de 29147.
\end{boxexo}

\begin{boxexo}[title={\ding{45}\quad Exercice~24.10. Processus interne \niveauexpert}]
Décrivez le processus interne de traitement d'une vulnérabilité signalée,
inspiré de 30111.
\end{boxexo}

\bigskip

\begin{boxinfo}
\textbf{Pour aller plus loin.}
\begin{itemize}
    \item ISO/IEC 29147~\cite{iso29147}~; ISO/IEC 30111~\cite{iso30111}.
    \item Loi 2010/012~\cite{loi2010012}.
\end{itemize}
\end{boxinfo}

\begin{boxinfo}
\textbf{Transition.} Toutes ces compétences se cultivent dans un cadre
d'entraînement~: le Chapitre~\ref{chap:labo} traite du laboratoire et de la
pratique continue.
\end{boxinfo}
}{}
        \IfFileExists{partie6/P06_C25_laboratoire.tex}{% ============================================================================
% Fichier  : partie6/P06_C25_laboratoire.tex
% Titre    : Le Laboratoire et la Pratique Continue
% Auteur   : MINKA MI NGUIDJOÏ Thierry Emmanuel
% Version  : 1.0 -- novembre 2026
% ============================================================================
% Sources : pratiques d'entraînement en sécurité offensive ; cadre légal des
% laboratoires. NDAP SUKULU : plateforme d'entraînement du cours.
% ============================================================================

\chapter{Le Laboratoire et la Pratique Continue}
\label{chap:labo}
\chaptermark{Le laboratoire et la pratique continue}

\epigraph{%
    « On n'apprend pas à tester en lisant~;\\
    on apprend dans un laboratoire que l'on a le droit d'attaquer. »%
}{Principe du cours}

\niveauChapitre{Tous niveaux}{L'ensemble du cours.}

\begin{boxobjectifs}
\begin{itemize}
    \item Comprendre pourquoi l'entraînement se fait en laboratoire fermé.
    \item Monter un laboratoire personnel légal.
    \item Utiliser les plateformes d'entraînement et les épreuves.
    \item Entretenir ses compétences dans la durée et dans la légalité.
    \item Situer la plateforme NDAP SUKULU du cours.
\end{itemize}
\end{boxobjectifs}

\bigskip

% ============================================================================
\section*{Question directrice : où s'entraîner sans jamais enfreindre la loi~?}
% ============================================================================

Les techniques de ce cours ne se pratiquent que sur des systèmes que l'on a
le droit d'attaquer. La tentation de « s'exercer » sur un site réel est une
faute et une infraction. Ce chapitre montre où et comment s'entraîner
légalement, et présente la plateforme du cours.

\medskip
\textbf{Plan du chapitre.}
\begin{enumerate}
    \item Pourquoi un laboratoire fermé (section~\ref{sec:c25-pourquoi}).
    \item Monter son laboratoire (section~\ref{sec:c25-monter}).
    \item Plateformes et épreuves (section~\ref{sec:c25-plateformes}).
    \item Entretenir ses compétences (section~\ref{sec:c25-entretien}).
    \item La plateforme NDAP SUKULU (section~\ref{sec:c25-ndap}).
\end{enumerate}

\separateur

% ============================================================================
\section{Pourquoi un Laboratoire Fermé}
\label{sec:c25-pourquoi}
% ============================================================================

\begin{boxalert}
\textbf{Violation typique.} S'entraîner sur un site réel « juste pour voir ».

\smallskip
\textbf{Règle~:} tout essai sur un système dont on n'a pas l'autorisation
écrite est une infraction (Chapitre~\ref{chap:droit}), quelles que soient
l'intention et l'absence de dommage. L'entraînement se fait en laboratoire
fermé, isolé du réseau réel.
\end{boxalert}

% ============================================================================
\section{Monter son Laboratoire}
\label{sec:c25-monter}
% ============================================================================

\begin{boxprocedure}
\textbf{Monter un laboratoire personnel légal}
\begin{enumerate}
    \item Un hyperviseur sur sa propre machine~; un réseau virtuel isolé.
    \item Des machines cibles prévues pour l'entraînement (volontairement
          vulnérables, fournies à cet effet).
    \item Aucune connexion de ce réseau au réseau réel ni à l'Internet lors
          des essais offensifs.
    \item Des instantanés pour revenir à l'état initial.
    \item Un journal d'entraînement pour suivre ses progrès.
\end{enumerate}
\end{boxprocedure}

% ============================================================================
\section{Plateformes et Épreuves}
\label{sec:c25-plateformes}
% ============================================================================

Des plateformes en ligne proposent des cibles légales à attaquer et des
épreuves de capture de drapeau. Elles offrent un cadre autorisé~: on y a le
droit explicite de tester. Elles complètent le laboratoire personnel et
permettent de se mesurer.

\begin{tableaucours}{L{4.0cm}L{8.2cm}}{%
    \tableauHeader{Cadre} &
    \tableauHeader{Usage}
}
Laboratoire personnel & Apprendre à son rythme, tout casser sans risque \\
Plateforme d'entraînement en ligne & Cibles variées, légales, mises à jour \\
Épreuve de capture de drapeau & Se mesurer, travailler en équipe, sous
pression \\
Plateforme du cours (NDAP SUKULU) & Fiches, labs et exercices alignés sur ce
cours \\
\end{tableaucours}
\vspace{-0.8em}
\begin{center}{\small\itshape Tableau~25.1. Cadres d'entraînement légaux}\end{center}

% ============================================================================
\section{Entretenir ses Compétences}
\label{sec:c25-entretien}
% ============================================================================

La sécurité offensive évolue vite. Les techniques de fond (reconnaissance,
accès, élévation, latéralité, restitution) restent~; les détails changent.
On entretient ses compétences par une pratique régulière, une veille
mesurée, et le partage au sein d'une communauté.

\begin{boxmetier}{Pentesteur, analyste SOC / Red Team, RSSI}
\textbf{Pratiquer, dans la durée et la légalité.} Le professionnel réserve
un temps régulier au laboratoire, documente ses essais, et ne cède jamais à
la facilité d'un essai « réel » non autorisé~: sa réputation et sa liberté
en dépendent.
\end{boxmetier}

% ============================================================================
\section{La Plateforme NDAP SUKULU}
\label{sec:c25-ndap}
% ============================================================================

Le cours s'appuie sur la plateforme \sigle{NDAP SUKULU}~: pour chaque
chapitre, une fiche d'auto-évaluation~; pour les chapitres pratiques, un lab
qui fait manipuler des pièces (relevés, historiques, journaux) dans un cadre
entièrement fermé. Les dix labs du cours (HE2-LAB-01 à~10) couvrent le
périmètre, l'OSINT, le balayage, la priorisation, l'annuaire, l'élévation,
l'audit de mots de passe, l'historique web, la couverture de détection et le
calcul de gravité.

\separateur

\begin{boxresume}
\textbf{Ce qu'il faut retenir de ce chapitre~:}
\begin{itemize}
    \item \textbf{On s'entraîne} seulement sur ce qu'on a le droit
          d'attaquer.
    \item \textbf{Le laboratoire fermé} est isolé du réseau réel.
    \item \textbf{Plateformes et épreuves} offrent des cibles légales.
    \item \textbf{On entretient} ses compétences par une pratique régulière.
    \item \textbf{NDAP SUKULU} aligne fiches et labs sur le cours.
\end{itemize}
\end{boxresume}

\begin{boxinfo}
\textbf{Points de contrôle de l'élève.} Je sais~: expliquer pourquoi
l'entraînement est en laboratoire fermé~; monter un laboratoire légal~;
citer les cadres d'entraînement~; organiser l'entretien de mes
compétences~; décrire NDAP SUKULU.
\end{boxinfo}

\bigskip

\begin{boxexo}[title={\ding{45}\quad Exercice~25.1. Légalité \niveaubase}]
Pourquoi ne s'entraîne-t-on jamais sur un site réel sans autorisation~?
\end{boxexo}

\begin{boxexo}[title={\ding{45}\quad Exercice~25.2. Isolation \niveaubase}]
Pourquoi le réseau du laboratoire doit-il être isolé du réseau réel~?
\end{boxexo}

\begin{boxexo}[title={\ding{45}\quad Exercice~25.3. Cadres \niveaubase}]
Citez trois cadres d'entraînement légaux.
\end{boxexo}

\begin{boxexo}[title={\ding{45}\quad Exercice~25.4. Instantanés \niveaumoyen}]
À quoi servent les instantanés dans un laboratoire~?
\end{boxexo}

\begin{boxexo}[title={\ding{45}\quad Exercice~25.5. Capture de drapeau \niveaumoyen}]
Qu'apporte une épreuve de capture de drapeau par rapport au laboratoire
personnel~?
\end{boxexo}

\begin{boxexo}[title={\ding{45}\quad Exercice~25.6. Veille \niveaumoyen}]
Comment organiser une veille mesurée sans s'y noyer~?
\end{boxexo}

\begin{boxexo}[title={\ding{45}\quad Exercice~25.7. Journal \niveaumoyen}]
Que consigne-t-on dans un journal d'entraînement~?
\end{boxexo}

\begin{boxexo}[title={\ding{45}\quad Exercice~25.8. NDAP SUKULU \niveaumoyen}]
Décrivez ce que la plateforme du cours apporte à chaque chapitre.
\end{boxexo}

\begin{boxexo}[title={\ding{45}\quad Exercice~25.9. Montage \niveauexpert}]
Décrivez le montage complet d'un laboratoire personnel pour pratiquer les
Parties~III et~IV.
\end{boxexo}

\begin{boxexo}[title={\ding{45}\quad Exercice~25.10. Plan annuel \niveauexpert}]
Proposez votre plan d'entraînement sur un an pour entretenir les compétences
du cours.
\end{boxexo}

\bigskip

\begin{boxinfo}
\textbf{Pour aller plus loin.}
\begin{itemize}
    \item Les dix labs de NDAP SUKULU pour ce cours.
\end{itemize}
\end{boxinfo}

\begin{boxinfo}
\textbf{Transition.} Reste à savoir mesurer la valeur de tout cela~: le
Chapitre~\ref{chap:mesurer} clôt le cours sur la mesure de la sécurité
offensive.
\end{boxinfo}
}{}
        \IfFileExists{partie6/P06_C26_mesurer.tex}{% ============================================================================
% Fichier  : partie6/P06_C26_mesurer.tex
% Titre    : Mesurer la Sécurité Offensive
% Auteur   : MINKA MI NGUIDJOÏ Thierry Emmanuel
% Version  : 1.0 -- novembre 2026
% ============================================================================
% Sources : cadre de verdict OUI / NON / NEK de l'auditologie (MINKA) ;
% indicateurs de sécurité offensive. Opérations LOGONE et MBAM : cas fictifs.
% ============================================================================

\chapter{Mesurer la Sécurité Offensive}
\label{chap:mesurer}
\chaptermark{Mesurer la sécurité offensive}

\epigraph{%
    « Un test qui ne se mesure pas ne s'améliore pas~;\\
    et ce qui ne peut être tranché doit être dit indéterminé. »%
}{Principe du cours}

\niveauChapitre{Avancé}{L'ensemble du
cours~; notions d'auditologie (cours d'audit).}

\begin{boxobjectifs}
\begin{itemize}
    \item Définir ce que mesure un test d'intrusion.
    \item Construire des indicateurs utiles et honnêtes.
    \item Appliquer le verdict à trois valeurs \sigle{OUI}, \sigle{NON},
          \sigle{NEK}.
    \item Éviter les indicateurs trompeurs.
    \item Inscrire le test dans un cycle d'amélioration.
\end{itemize}
\end{boxobjectifs}

\bigskip

% ============================================================================
\section*{Question directrice : LiZbiZ est-elle « sûre » après le test~?}
% ============================================================================

À la fin de la mission, M.~ETEME pose la question la plus simple et la plus
piégée~: « sommes-nous sûrs maintenant~? ». La réponse honnête n'est ni oui
ni non en général~: elle est un ensemble de verdicts précis, et, pour ce qui
n'a pu être tranché, un franc « indéterminé ». Ce chapitre apprend à
mesurer sans tromper.

\medskip
\textbf{Plan du chapitre.}
\begin{enumerate}
    \item Ce que mesure un test (section~\ref{sec:c26-quoi}).
    \item Des indicateurs honnêtes (section~\ref{sec:c26-indicateurs}).
    \item Le verdict à trois valeurs (section~\ref{sec:c26-verdict}).
    \item Les indicateurs trompeurs (section~\ref{sec:c26-trompeurs}).
    \item Le cycle d'amélioration (section~\ref{sec:c26-cycle}).
\end{enumerate}

\separateur

% ============================================================================
\section{Ce que Mesure un Test}
\label{sec:c26-quoi}
% ============================================================================

Un test d'intrusion ne mesure pas « la sécurité » en général~: il mesure la
résistance d'un périmètre donné, à un moment donné, face à un effort donné.
Il dit ce qu'un attaquant de niveau comparable aurait pu faire, pas ce qu'un
attaquant quelconque fera demain.

% ============================================================================
\section{Des Indicateurs Honnêtes}
\label{sec:c26-indicateurs}
% ============================================================================

\begin{tableaucours}{L{4.8cm}L{7.4cm}}{%
    \tableauHeader{Indicateur} &
    \tableauHeader{Ce qu'il dit}
}
Objectifs atteints sur objectifs fixés & Résistance aux scénarios visés \\
Constats par gravité & Ampleur et priorité des faiblesses \\
Couverture de détection & Part des techniques vues par le défenseur \\
Temps jusqu'à la détection & Rapidité du défenseur \\
Taux de correction au contre-test & Qualité du traitement \\
Taux de signalement d'hameçonnage & Vigilance humaine \\
\end{tableaucours}
\vspace{-0.8em}
\begin{center}{\small\itshape Tableau~26.1. Indicateurs d'un test d'intrusion}\end{center}

% ============================================================================
\section{Le Verdict à Trois Valeurs}
\label{sec:c26-verdict}
% ============================================================================

\begin{boxdef}
\termecle{Verdict à trois valeurs}~: pour chaque question d'audit, la réponse
est \sigle{OUI} (la propriété est vérifiée, avec preuve), \sigle{NON} (elle
est réfutée, avec preuve), ou \sigle{NEK} (indéterminé~: les preuves
manquent pour trancher). Le \sigle{NEK} n'est pas un aveu de faiblesse~:
c'est une exigence d'honnêteté.
\end{boxdef}

\begin{tableaucours}{L{2.0cm}L{4.6cm}L{5.6cm}}{%
    \tableauHeader{Verdict} &
    \tableauHeader{Sens} &
    \tableauHeader{Exemple (LOGONE)}
}
OUI & Propriété démontrée & « Le système de paiement cloisonné a résisté »
(F3 non atteint, MBAM) \\
NON & Propriété réfutée & « Un attaquant interne devient administrateur du
domaine » (T-01) \\
NEK & Indéterminé & « La résistance du nuage au-delà du périmètre testé » (hors
mandat) \\
\end{tableaucours}
\vspace{-0.8em}
\begin{center}{\small\itshape Tableau~26.2. Le verdict OUI / NON / NEK appliqué~\cite{minka2026auditologie}}\end{center}

\begin{boxalert}
\textbf{Violation typique.} Conclure \sigle{OUI} (« c'est sûr ») faute
d'avoir trouvé une faille, alors qu'on n'a pas cherché.

\smallskip
\textbf{Règle~:} l'absence de preuve n'est pas une preuve d'absence. Ce qui
n'a pas été testé, ou n'a pu être tranché, est \sigle{NEK}, pas \sigle{OUI}.
Le \sigle{NEK} dit au client où porter l'effort suivant.
\end{boxalert}

% ============================================================================
\section{Les Indicateurs Trompeurs}
\label{sec:c26-trompeurs}
% ============================================================================

\begin{tableaucours}{L{4.4cm}L{7.8cm}}{%
    \tableauHeader{Indicateur trompeur} &
    \tableauHeader{Pourquoi il trompe}
}
Nombre brut de failles & Dépend de l'outil, pas du risque \\
« Zéro faille trouvée » & Peut signifier « rien cherché » \\
Note globale unique & Masque les constats critiques dans la moyenne \\
Temps passé & Mesure l'effort, pas le résultat \\
\end{tableaucours}
\vspace{-0.8em}
\begin{center}{\small\itshape Tableau~26.3. Indicateurs à manier avec prudence}\end{center}

% ============================================================================
\section{Le Cycle d'Amélioration}
\label{sec:c26-cycle}
% ============================================================================

Un test n'est pas une fin~: c'est une étape d'un cycle. Les constats
alimentent les corrections~; le contre-test les vérifie~; la phase violette
nourrit la détection~; le test suivant repart des \sigle{NEK} et des zones
non couvertes. La sécurité offensive sert la défense, et se mesure à ce
qu'elle lui apporte.

\begin{boxmetier}{RSSI, auditeur de sécurité, pentesteur}
\textbf{Dire le vrai, y compris le NEK.} Le professionnel ne cède pas à la
pression du \sigle{OUI} rassurant. Il distingue ce qu'il a démontré, ce qu'il
a réfuté, et ce qu'il n'a pu trancher~: c'est à ce triple langage qu'on
reconnaît un rapport honnête.
\end{boxmetier}

\separateur

\begin{boxresume}
\textbf{Ce qu'il faut retenir de ce chapitre~:}
\begin{itemize}
    \item \textbf{Un test} mesure une résistance, à un moment, face à un
          effort donné.
    \item \textbf{Les bons indicateurs} relient objectif, gravité, détection,
          correction.
    \item \textbf{Le verdict} est \sigle{OUI}, \sigle{NON} ou \sigle{NEK}.
    \item \textbf{L'absence de preuve} n'est pas preuve d'absence~: c'est
          \sigle{NEK}.
    \item \textbf{Le test} sert un cycle d'amélioration continue de la
          défense.
\end{itemize}
\end{boxresume}

\begin{boxinfo}
\textbf{Points de contrôle de l'élève.} Je sais~: dire ce que mesure un
test~; construire des indicateurs honnêtes~; appliquer OUI / NON / NEK~;
reconnaître les indicateurs trompeurs~; inscrire le test dans un cycle.
\end{boxinfo}

\bigskip

\begin{boxexo}[title={\ding{45}\quad Exercice~26.1. Ce que mesure un test \niveaubase}]
Un test mesure-t-il « la sécurité » en général~? Que mesure-t-il~?
\end{boxexo}

\begin{boxexo}[title={\ding{45}\quad Exercice~26.2. Trois valeurs \niveaubase}]
Donnez le sens de OUI, NON et NEK.
\end{boxexo}

\begin{boxexo}[title={\ding{45}\quad Exercice~26.3. Trompeur \niveaubase}]
Pourquoi « zéro faille trouvée » peut-il tromper~?
\end{boxexo}

\begin{boxexo}[title={\ding{45}\quad Exercice~26.4. Absence de preuve \niveaumoyen}]
Expliquez « l'absence de preuve n'est pas une preuve d'absence » sur un
exemple de LOGONE.
\end{boxexo}

\begin{boxexo}[title={\ding{45}\quad Exercice~26.5. Classer \niveaumoyen}]
Attribuez OUI, NON ou NEK~: le système de paiement de MBAM~; l'accès
administrateur de LiZbiZ~; la sécurité du nuage hors périmètre.
\end{boxexo}

\begin{boxexo}[title={\ding{45}\quad Exercice~26.6. Note unique \niveaumoyen}]
Pourquoi une note globale unique est-elle dangereuse~?
\end{boxexo}

\begin{boxexo}[title={\ding{45}\quad Exercice~26.7. Indicateurs \niveaumoyen}]
Choisissez trois indicateurs pour le tableau de bord de LiZbiZ et justifiez.
\end{boxexo}

\begin{boxexo}[title={\ding{45}\quad Exercice~26.8. Répondre au DG \niveaumoyen}]
Rédigez la réponse honnête à « sommes-nous sûrs maintenant~? », en OUI / NON
/ NEK.
\end{boxexo}

\begin{boxexo}[title={\ding{45}\quad Exercice~26.9. Cycle \niveauexpert}]
Décrivez le cycle d'amélioration reliant test, correction, contre-test,
détection et test suivant.
\end{boxexo}

\begin{boxexo}[title={\ding{45}\quad Exercice~26.10. Tableau de bord \niveauexpert}]
Construisez le tableau de bord de sécurité offensive que LiZbiZ devrait
suivre d'une année sur l'autre, verdicts NEK compris.
\end{boxexo}

\bigskip

\begin{boxinfo}
\textbf{Pour aller plus loin.}
\begin{itemize}
    \item Cadre d'auditologie et verdict à trois valeurs~\cite{minka2026auditologie}.
\end{itemize}
\end{boxinfo}

\begin{boxinfo}
\textbf{Fin du cours.} De la lettre d'autorisation au verdict mesuré, le
test d'intrusion aura montré qu'autoriser, reconnaître, tester et restituer
forment un même service~: défendre, en pensant comme l'attaquant, sans
jamais le devenir.
\end{boxinfo}
}{}

% --------------------------------------------------------------------------
% ANNEXES
% --------------------------------------------------------------------------
\appendix

    %% ============================================================================
%% Fichier  : annexes/Ann_A_glossaire.tex
%% Titre    : Glossaire
%% Version  : 1.0 -- novembre 2026
%% ============================================================================

\chapter{Glossaire}
\label{ann:glossaire}

Ce glossaire reprend, sous une forme abrégée, les notions définies dans les
chapitres. Le renvoi indique le chapitre où la notion est présentée.

\begin{description}
    \item[Active Directory] Annuaire des réseaux Windows~: comptes, groupes, machines et droits d'un domaine (Ch.~\ref{chap:ad}).
    \item[Charge utile] Ce qu'un exploit fait exécuter après avoir déclenché la faille (Ch.~\ref{chap:exploitation}).
    \item[Contre-test] Vérification, après correction, que chaque faiblesse rapportée a bien été traitée (Ch.~\ref{chap:retest}).
    \item[CVSS] Système commun de notation des vulnérabilités, de 0 à 10 (Ch.~\ref{chap:vulnerabilites}).
    \item[Divulgation coordonnée] Signalement d'une faille à l'éditeur, suivi d'une publication après correction ou délai convenu (Ch.~\ref{chap:divulgation}).
    \item[Élévation de privilèges] Obtention de droits supérieurs à ceux du compte de départ (Ch.~\ref{chap:privileges}).
    \item[Empoisonnement de résolution] Réponse frauduleuse à une requête de résolution de nom diffusée, qui détourne la connexion de la victime (Ch.~\ref{chap:reseau}).
    \item[Énumération] Recueil détaillé des services, comptes et partages exposés, après le balayage (Ch.~\ref{chap:balayage}).
    \item[Équipe rouge] Exercice offensif visant un objectif précis, discrètement, pour mesurer la détection (Ch.~\ref{chap:rouge}).
    \item[Équipe violette] Démarche où offensive et défensive travaillent ensemble pour améliorer la détection (Ch.~\ref{chap:rouge}).
    \item[Exploit] Moyen technique de tirer parti d'une vulnérabilité (Ch.~\ref{chap:exploitation}).
    \item[Faux négatif] Vulnérabilité réelle qu'un outil ne signale pas (Ch.~\ref{chap:vulnerabilites}).
    \item[Faux positif] Alerte d'un outil sans vulnérabilité réelle (Ch.~\ref{chap:vulnerabilites}).
    \item[Hacking éthique] Recherche autorisée de faiblesses dans un système, pour le renforcer, dans un cadre légal et sans dommage (Ch.~\ref{chap:ethique}).
    \item[Ingénierie sociale] Manipulation d'une personne pour obtenir une information, un accès ou une action (Ch.~\ref{chap:social}).
    \item[Injection] Faille où une donnée fournie est interprétée comme une commande, faute de séparation (Ch.~\ref{chap:injections}).
    \item[Interface de programmation (API)] Points d'entrée par lesquels des programmes échangent des données (Ch.~\ref{chap:api}).
    \item[Kerberoasting] Vol d'un ticket de service Kerberos, cassé hors ligne si le mot de passe du compte de service est faible (Ch.~\ref{chap:ad}).
    \item[Matrice de couverture de détection] Tableau croisant les techniques avec l'état de leur détection (Ch.~\ref{chap:detection}).
    \item[MITRE ATT\&CK] Base de connaissances des tactiques et techniques des attaquants (Ch.~\ref{chap:methodes}).
    \item[Modèle d'administration par niveaux] Organisation qui interdit d'utiliser un compte à privilèges sur une machine de niveau inférieur (Ch.~\ref{chap:ad}).
    \item[Modèle de responsabilité partagée] Répartition des tâches de sécurité entre le fournisseur de nuage et le client (Ch.~\ref{chap:nuage}).
    \item[Mouvement horizontal] Passage d'un compte à un autre de même niveau (Ch.~\ref{chap:privileges}).
    \item[Mouvement latéral] Prise de contrôle d'autres machines du réseau avec des identifiants valides (Ch.~\ref{chap:post}).
    \item[OSINT] Renseignement tiré de sources ouvertes, sans solliciter la cible (Ch.~\ref{chap:osint}).
    \item[Point d'accès pirate] Borne non autorisée branchée sur le réseau de l'entreprise (Ch.~\ref{chap:sansfil}).
    \item[Portail jumeau] Fausse borne qui imite le réseau légitime pour capter les connexions (Ch.~\ref{chap:sansfil}).
    \item[Preuve de concept] Démonstration minimale et inoffensive qu'une faille est réelle (Ch.~\ref{chap:exploitation}).
    \item[Reconnaissance active] Collecte d'informations en sollicitant directement la cible (Ch.~\ref{chap:osint}).
    \item[Reconnaissance passive] Collecte d'informations sans solliciter la cible (Ch.~\ref{chap:osint}).
    \item[Référence directe non contrôlée à un objet] Accès à l'objet d'autrui en changeant un identifiant visible, faute de contrôle (Ch.~\ref{chap:injections}).
    \item[Règles d'engagement] Cadre convenu d'une mission~: périmètre, horaires, actions permises, contacts, mot d'arrêt (Ch.~\ref{chap:engagement}).
    \item[Requête HTTP] Message du client au serveur~: méthode, chemin, en-têtes, corps (Ch.~\ref{chap:web}).
    \item[Réponse HTTP] Message du serveur au client~: code, en-têtes, corps (Ch.~\ref{chap:web}).
    \item[Session] État conservé entre requêtes, souvent par un cookie (Ch.~\ref{chap:web}).
    \item[Signature SMB] Mécanisme qui lie les échanges de partage et empêche leur relais (Ch.~\ref{chap:reseau}).
    \item[Test d'intrusion] Simulation autorisée d'attaque pour démontrer ce qu'un attaquant ferait (Ch.~\ref{chap:ethique}).
    \item[Top 10 de l'OWASP] Classement des grandes familles de risques des applications web ou des API (Ch.~\ref{chap:web}, \ref{chap:api}).
    \item[Transparence des certificats] Journaux publics des certificats émis, utiles pour cartographier la surface (Ch.~\ref{chap:osint}).
    \item[Verdict à trois valeurs (OUI, NON, NEK)] Réponse d'audit~: vérifiée, réfutée, ou indéterminée faute de preuves (Ch.~\ref{chap:mesurer}).
    \item[Web Security Testing Guide (WSTG)] Guide de l'OWASP organisant le test d'une application web en catégories (Ch.~\ref{chap:web}).
\end{description}

    %% ============================================================================
%% Fichier  : annexes/Ann_B_outils.tex
%% Titre    : Référentiel des outils
%% Version  : 1.0 -- novembre 2026
%% ============================================================================

\chapter{Référentiel des Outils}
\label{ann:outils}

Les outils cités sont des exemples. Le cours insiste sur la méthode et la
compréhension~: l'outil sert le raisonnement, il ne le remplace pas. Tout
usage se fait dans un cadre autorisé ou en laboratoire fermé
(Ch.~\ref{chap:labo}).

\begin{tableaucours}{L{4.2cm}L{5.4cm}L{2.8cm}}{%
    \tableauHeader{Catégorie d'outil} &
    \tableauHeader{Usage} &
    \tableauHeader{Chapitre, lab}
}
Sources ouvertes (OSINT) & Journaux de certificats, moteurs de recherche, registres publics & Ch.~\ref{chap:osint}, HE2-LAB-02 \\
Balayage réseau & Découverte des hôtes et des services & Ch.~\ref{chap:balayage}, HE2-LAB-03 \\
Scanner de vulnérabilités & Détection assistée, à vérifier & Ch.~\ref{chap:vulnerabilites}, HE2-LAB-04 \\
Relais web local & Observation et modification des requêtes & Ch.~\ref{chap:web}, HE2-LAB-08 \\
Analyse d'annuaire & Groupes, comptes, chemins d'attaque & Ch.~\ref{chap:ad}, HE2-LAB-05 \\
Audit de mots de passe & Cassage hors ligne autorisé & Ch.~\ref{chap:mdp}, HE2-LAB-07 \\
Référentiel ATT\&CK & Planifier et mesurer la détection & Ch.~\ref{chap:detection}, HE2-LAB-09 \\
Tableur & Priorisation, calcul de gravité, matrices & Ch.~\ref{chap:vulnerabilites}, HE2-LAB-10 \\
Laboratoire virtuel & Entraînement en environnement fermé & Ch.~\ref{chap:labo} \\
\end{tableaucours}

\begin{boxalert}
\textbf{Rappel.} Aucun de ces outils ne se dirige vers un système réel sans
autorisation écrite. En dehors d'une mission mandatée, l'entraînement se fait
en laboratoire isolé (Ch.~\ref{chap:labo}).
\end{boxalert}

    %% ============================================================================
%% Fichier  : annexes/Ann_C_templates.tex
%% Titre    : Modèles et procédures
%% Version  : 1.0 -- novembre 2026
%% ============================================================================

\chapter{Modèles et Procédures}
\label{ann:templates}

Les modèles ci-dessous reprennent ceux des chapitres, sous une forme prête à
l'emploi.

\section*{Lettre d'autorisation}

\begin{tableaucours}{L{3.4cm}L{8.8cm}}{%
    \tableauHeader{Rubrique} &
    \tableauHeader{Contenu}
}
Parties & Commanditaire et prestataire, signataires habilités \\
Périmètre & Adresses et systèmes visés (notation CIDR), exclusions \\
Dates et horaires & Fenêtres autorisées \\
Règles d'engagement & Actions permises et interdites, mot d'arrêt \\
Contacts d'urgence & Joignables jour et nuit \\
Données et preuves & Minimisation, masquage, destruction en fin \\
Signature & Responsable habilité, date \\
\end{tableaucours}

\section*{Fiche de constat}

\begin{tableaucours}{L{3.4cm}L{8.8cm}}{%
    \tableauHeader{Rubrique} &
    \tableauHeader{Contenu}
}
Identifiant et titre & T-xx, titre court \\
Description & La faiblesse, en termes clairs \\
Preuve & Requête, capture masquée, relevé \\
Impact & Ce qu'un attaquant en ferait \\
Gravité & Cotation justifiée (CVSS ou risque métier) \\
Recommandation & Correction concrète et vérifiable \\
Détection & Comment voir l'exploitation en production \\
\end{tableaucours}

\section*{Procédure de nettoyage de fin de mission}

\begin{enumerate}
    \item Relire le journal de mission~: tout ce qui a été créé ou modifié.
    \item Supprimer comptes, fichiers, services, tâches créés.
    \item Restaurer les configurations modifiées.
    \item Faire vérifier la suppression par le client.
    \item Consigner le nettoyage dans le rapport.
\end{enumerate}

\section*{Plan du rapport}

\begin{enumerate}
    \item Synthèse pour la direction.
    \item Contexte et périmètre.
    \item Démarche et couverture.
    \item Constats détaillés, cotés.
    \item Plan de correction priorisé.
    \item Annexes~: preuves, actions horodatées.
\end{enumerate}

    %% ============================================================================
%% Fichier  : annexes/Ann_D_textes.tex
%% Titre    : Textes juridiques et normatifs
%% Version  : 1.0 -- novembre 2026
%% ============================================================================

\chapter{Textes Juridiques et Normatifs}
\label{ann:textes}

Seuls les articles lus dans le texte officiel sont cités entre guillemets dans
les chapitres. Ce tableau indique où les trouver. Les numéros d'articles non
vérifiés ne sont pas reproduits~; on renvoie au texte publié.

\section*{Droit camerounais}

\begin{tableaucours}{L{4.4cm}L{5.2cm}L{2.6cm}}{%
    \tableauHeader{Texte} &
    \tableauHeader{Objet} &
    \tableauHeader{Chapitre}
}
Loi n° 2010/012 du 21 décembre 2010 & Cybersécurité et cybercriminalité~; accès et maintien frauduleux & \ref{chap:droit} \\
Loi n° 2024/017 du 23 décembre 2024 & Protection des données à caractère personnel & \ref{chap:droit} \\
\end{tableaucours}

\begin{boxinfo}
Les articles de ces lois se lisent dans le texte publié au Journal officiel de
la République du Cameroun. Le cours ne cite que des articles vérifiés dans ce
texte.
\end{boxinfo}

\section*{Normes et guides internationaux}

\begin{tableaucours}{L{4.4cm}L{5.2cm}L{2.6cm}}{%
    \tableauHeader{Référence} &
    \tableauHeader{Objet} &
    \tableauHeader{Chapitre}
}
NIST SP 800-115 & Guide technique du test de sécurité & \ref{chap:methodes} \\
PTES & Standard d'exécution d'un test d'intrusion & \ref{chap:methodes} \\
OSSTMM 3 & Méthodologie de test de sécurité & \ref{chap:methodes} \\
OWASP WSTG & Guide de test des applications web & \ref{chap:web} \\
OWASP Top 10 & Risques des applications web & \ref{chap:web} \\
OWASP API Top 10 & Risques des interfaces de programmation & \ref{chap:api} \\
OWASP MASVS & Vérification de la sécurité des applications mobiles & \ref{chap:api} \\
CIS Controls v8 & Mesures de sécurité prioritaires & \ref{chap:ethique} \\
CVSS 3.1 (FIRST) & Notation des vulnérabilités & \ref{chap:vulnerabilites} \\
MITRE ATT\&CK & Tactiques et techniques des attaquants & \ref{chap:methodes} \\
ISO/IEC 29147 & Divulgation de vulnérabilités & \ref{chap:divulgation} \\
ISO/IEC 30111 & Traitement des vulnérabilités & \ref{chap:divulgation} \\
RFC 4120 & Kerberos version 5 & \ref{chap:reseau} \\
RFC 4795 & Résolution de noms locale (LLMNR) & \ref{chap:reseau} \\
\end{tableaucours}

    %% ============================================================================
%% Fichier  : annexes/Ann_E_bibliographie.tex
%% Titre    : Bibliographie commentée
%% Version  : 1.0 -- novembre 2026
%% ============================================================================

\chapter{Bibliographie Commentée}
\label{ann:bibliographie}

Pour chaque partie, quelques lectures, avec ce qu'elles apportent. Les
références complètes figurent dans la bibliographie du livre.

\section*{Partie I · Le cadre}
\begin{itemize}
    \item CIS Controls v8~\cite{cis2021controls}~: la place du test dans un programme de sécurité.
    \item NIST SP 800-115~\cite{nist800115}~: le cadre d'un test de sécurité.
    \item Lois 2010/012~\cite{loi2010012} et 2024/017~\cite{loi2024017}~: le droit applicable au Cameroun.
\end{itemize}

\section*{Partie II · Reconnaître}
\begin{itemize}
    \item OWASP WSTG~\cite{owaspwstg}~: la collecte d'informations sur une application.
    \item CVSS 3.1~\cite{first2019cvss31}, EPSS~\cite{firstepss}, catalogue KEV~\cite{cisakev}~: prioriser les vulnérabilités.
    \item Programmes CVE~\cite{mitrecve}~: identifier les vulnérabilités publiées.
\end{itemize}

\section*{Partie III · Les systèmes}
\begin{itemize}
    \item MITRE ATT\&CK~\cite{mitreattack}~: le vocabulaire des techniques.
    \item RFC 4120~\cite{rfc4120} et RFC 4795~\cite{rfc4795}~: Kerberos et la résolution locale.
    \item NIST SP 800-63B~\cite{nist80063b}~: les mots de passe.
\end{itemize}

\section*{Partie IV · Les applications}
\begin{itemize}
    \item OWASP Top 10~\cite{owasptop10} et API Top 10~\cite{owaspapitop10}~: les risques web et API.
    \item OWASP MASVS~\cite{owaspmasvs}~: la sécurité mobile.
    \item CIS Controls v8~\cite{cis2021controls}~: la configuration du nuage.
\end{itemize}

\section*{Partie V · Au-delà}
\begin{itemize}
    \item MITRE ATT\&CK~\cite{mitreattack}~: scénarios fondés sur une menace et couverture de détection.
\end{itemize}

\section*{Partie VI · Restituer}
\begin{itemize}
    \item PTES~\cite{ptes}~: la phase de restitution.
    \item ISO/IEC 29147~\cite{iso29147} et 30111~\cite{iso30111}~: la divulgation coordonnée.
    \item T.~E.~Minka Mi Nguidjoï~\cite{minka2026auditologie}~: le verdict à trois valeurs et la mesure.
\end{itemize}

    %% ============================================================================
%% Fichier  : annexes/Ann_F_contacts.tex
%% Titre    : Contacts et ressources
%% Version  : 1.0 -- novembre 2026
%% ============================================================================

\chapter{Contacts et Ressources}
\label{ann:contacts}

Les adresses sont celles des sites institutionnels à la date de rédaction~; à
vérifier.

\begin{tableaucours}{L{4.4cm}L{4.8cm}L{3.2cm}}{%
    \tableauHeader{Organisation} &
    \tableauHeader{Rôle} &
    \tableauHeader{Site}
}
ANTIC & Régulation de la sécurité électronique au Cameroun & \texttt{www.antic.cm} \\
MINPOSTEL & Tutelle des communications électroniques & \texttt{www.minpostel.gov.cm} \\
OWASP & Guides de test et Top 10 & \texttt{owasp.org} \\
MITRE & ATT\&CK, CVE & \texttt{attack.mitre.org} \\
FIRST & CVSS, EPSS & \texttt{www.first.org} \\
NIST & Publications SP 800 & \texttt{www.nist.gov} \\
CIS & Controls et Benchmarks & \texttt{www.cisecurity.org} \\
ISO & Normes ISO et ISO/IEC & \texttt{www.iso.org} \\
\end{tableaucours}

\begin{boxinfo}
\textbf{La plateforme du cours.} Les fiches d'auto-évaluation et les dix labs
de \sigle{NDAP SUKULU} pour ce cours sont accessibles par le code
\texttt{SEC4062}.
\end{boxinfo}


% --------------------------------------------------------------------------
% ÉLÉMENTS FINAUX
% --------------------------------------------------------------------------
\backmatter

    \nocite{*}
    \printbibliography[heading=bibintoc, title={Références bibliographiques}]

    \printindex

% ==============================================================================
% 4ÈME DE COUVERTURE
% ==============================================================================
\clearpage
\thispagestyle{empty}
\pagecolor{CouleurPrimaire}

\begin{tikzpicture}[remember picture, overlay]
    \fill[CouleurAccent]
        (current page.south west) rectangle
        ([yshift=1.8cm] current page.south east);
    \fill[CouleurAccent]
        ([yshift=-1.4cm] current page.north west) rectangle
        ([yshift=-1.7cm] current page.north east);
\end{tikzpicture}

\vspace*{2.5cm}

\begin{center}
    {\color{CouleurAccent}\rule{0.6\linewidth}{1pt}}\\[1.5ex]
    {\color{white}\Large\bfseries À propos de ce cours}\\[1.5ex]
    {\color{CouleurAccent}\rule{0.6\linewidth}{1pt}}
\end{center}

\vspace{1cm}

\begin{center}
\begin{minipage}{0.75\textwidth}
\color{white!85!gray}\setstretch{1.4}
Tester un système comme un attaquant, avec l'autorisation de son
propriétaire, pour qu'il soit mieux défendu~: c'est le hacking éthique. Ce
cours en enseigne la démarche, le droit, les techniques et la restitution.

\medskip

Il progresse en six parties, précédées d'un prologue~: le cadre du test
(droit, règles d'engagement, méthodologies), la reconnaissance, les systèmes
et les réseaux (dont Active Directory), les applications et le nuage, les
tests au-delà des systèmes (sans fil, ingénierie sociale, équipe rouge), et
la restitution (rapport, contre-test, divulgation, mesure).

\medskip

Chaque notion est ancrée dans les standards internationaux et dans le droit
camerounais. Le cours s'adresse aux élèves ingénieurs en génie informatique
de l'ENSP Yaoundé et aux professionnels qui exercent comme pentesteurs,
auditeurs de sécurité, analystes SOC ou membres d'une équipe rouge, ou
comme RSSI. Chaque technique y est présentée avec sa détection et sa
parade.
\end{minipage}
\end{center}

\vfill

\begin{center}
    {\color{CouleurAccent}\rule{0.7\linewidth}{1.2pt}}\\[1.2ex]
    {\color{white}\large\bfseries MINKA MI NGUIDJOÏ Thierry Emmanuel}\\[0.4ex]
    {\color{white!80!gray}\small
        CISA~·~CISM~·~CGEIT~·~CRISC~·~CDPSE~·~ISO 27001~·~ISO 22301~·~ISO 9001}\\[0.3ex]
    {\color{white!70!gray}\footnotesize\itshape
        Chercheur au laboratoire d'Ingénierie Mathématiques et des Systèmes d'Information}\\[0.3ex]
    {\color{white!70!gray}\footnotesize\itshape
        Chercheur~|~Cryptologue~|~Auditeur Senior~|~Expert Judiciaire~|~GRC \& Sécurité de l'Information}\\[1.2ex]
    {\color{CouleurAccent}\rule{0.7\linewidth}{1.2pt}}
\end{center}

\vspace{1cm}

\afterpage{\nopagecolor}
\end{document}
